\section*{Open Questions}
\begin{quote}\small
\noindent\textbf{Status.} \emph{Active live-burden/continuation block.}\par
\noindent\textbf{Block function.} This section separates what is still genuinely open from what survives only as historical labeling, so readers can distinguish live fronts from preserved older checkpoints.
\end{quote}
\addcontentsline{toc}{section}{Open Questions}
%% ============================================================

The structural core, the master equation, and their principal projections are
now explicit. The remaining work no longer concerns the existence of the
framework itself, but its residual quantitative closure, phenomenological
extraction, and spectrum-level completion. To avoid confusion, this section
separates the \emph{current live fronts} from the \emph{historical OP labels}
that are still used in later parts of the text.

\paragraph{Current reading.}
Several historically open OPs should no longer be read as fully open. In
particular, the Koide/triolic block is already structurally closed at the
present level, the Dirac-operator problem on $\Hplusminus$ is no longer open as
an independent operator-existence problem, and the photon-frame
$\Theta$-remainder step of the cosmological block is already closed. The live
questions are therefore the following.

\begin{remark}[Reader precedence for historical release snapshots inside the active line]
The next two synchronized-status blocks reconstruct older public-stage status only for navigation and burden genealogy. They do not overrule the current burden ledger, the current theorematic status, or the current tool/QAM/decoder routing of the active line. Whenever one of those historical snapshots sounds flatter or stronger than the normalized present reading, the normalized present reading governs.
\end{remark}

\begin{remark}[v3.23 synchronized status reading for the optical approximation layer]
\label{rem:v323-synchronized-status-reading-for-the-optical-approximation-layer}
The historical public v3.23.0 release line no longer permits a flat reading in which the new optical sector is either ``already a full optical theory'' or ``still only a loose analogy.'' After the first reduced prism-readout lemma, the residual-diagnostic layer, the carrier-side fit/calibration/validation chain, the reconstruction/correction steps, and the bundled optical approximation package, the optical layer now has its own staged status:
\begin{enumerate}[leftmargin=2em,label=(\roman*)]
  \item a \emph{strong global optical claim} remains open;
  \item a \emph{reduced retained-carrier optical sector} is already available;
  \item a \emph{first-order optical workflow} is already available;
  \item a \emph{bundled first-order optical package} is already available as an active citation surface.
\end{enumerate}
Thus the open-front layer must now distinguish between the still-open stronger global optical theory and the already stabilized first-order optical approximation layer.
\end{remark}

\begin{proposition}[v3.23 optical/open-question status synchronization principle]
\label{prop:v323-optical-open-question-status-synchronization-principle}
At the present theorematic and retained-carrier level, the following synchronized status reading holds for the new optical line:
\begin{enumerate}[leftmargin=2em,label=(\roman*)]
  \item the still-open strong claim is a full global optical theory of the crystal--shadow architecture, including non-reduced anisotropic branches, stronger inversion, and a later convergence theory for higher optical refinements;
  \item what is already available is a reduced prism-readout sector on one retained carrier surface with ordinary wavelength ordering;
  \item what is already available as first-order workflow is the chain consisting of residual diagnostics, carrier-side fit ansatz, spectral calibration, overdetermined validation, effective-index reconstruction, carrier-profile correction, and reduced stability/stopping control;
  \item what is already available as a single active surface is the bundled optical approximation package of Proposition~\ref{prop:first-bundled-optical-approximation-package};
  \item therefore the optical sector should no longer be counted among the fully open foundational fronts of the document, but as a partially stabilized first-order approximation layer with a still-open stronger global continuation.
\end{enumerate}
\end{proposition}

\begin{proof}
Clause (i) records exactly the stronger optical ambitions that the current line has not yet proved: unrestricted anisotropic branching, stronger inversion, and a genuine global convergence theory are still absent. Clause (ii) is supplied by the reduced prism-readout lemma. Clause (iii) is the cumulative content of the residual spectral-diagnostic lemma, the carrier-side fit principle, the calibration principle, the overdetermined consistency principle, the effective-index reconstruction principle, the carrier-profile correction principle, and the reduced stability principle. Clause (iv) is exactly the content of the bundled optical approximation package. Hence the optical line already has a stabilized first-order approximation layer even though the stronger global optical theory remains open.
\end{proof}

\begin{corollary}[Use rule for the synchronized v3.23 optical status reading]
\label{cor:use-rule-for-the-synchronized-v323-optical-status-reading}
Whenever a later argument in the present public v3.23.0 release line needs to classify whether the optical sector is already available for active use, it is sufficient to cite Proposition~\ref{prop:v323-optical-open-question-status-synchronization-principle}. A finer reopening is required only when the argument depends on a stronger global optical theory beyond the reduced retained-carrier approximation layer.
\end{corollary}

\begin{proof}
This is the operational reformulation of Proposition~\ref{prop:v323-optical-open-question-status-synchronization-principle}. The proposition already isolates the precise point at which the optical line ceases to count as merely a loose heuristic while still falling short of a full global optical theory.
\end{proof}


\begin{remark}[Historical v3.22/v3.23 synchronized status reading for the upgraded closure layer]
\label{rem:v322-synchronized-status-reading-upgraded-closure-layer}
The historical public v3.23.0 release line no longer permits a flat reading in which the upgraded closure/interface blocks are simply either ``open'' or ``proved.'' After the public v3.20 theorem-level shadow package, the public v3.21 interface-payoff layer, and the present v3.22 partial-closure steps, their status is now genuinely four-tiered:
\begin{enumerate}[leftmargin=2em,label=(\roman*)]
  \item a \emph{strong deferred claim} remains open;
  \item a \emph{theorem-level shadow} is already available;
  \item an \emph{interface-level operational use} is already available;
  \item for certain carrier-side conclusions, a \emph{restricted partial-closure / lift-back consequence} is already available.
\end{enumerate}
Thus the open-front layer must now distinguish not only between strong and weak forms, but also between theorematic shadow, interface use, and restricted closure consequence.
\end{remark}

\begin{proposition}[v3.22 hypothesis/open-question status synchronization principle]
\label{prop:v322-hypothesis-open-question-status-synchronization-principle}
At the present theorematic and interface level, the following synchronized status reading holds:
\begin{enumerate}[leftmargin=2em,label=(\roman*)]
  \item for Hypothesis~\ref{hyp:hc-meta-generator}, the still-open strong claim is the existence of one explicit compensator-centered generating equation $\mathfrak E_{\mathrm{HC}}=0$, while the theorem-level shadow of a common closure-stable carrier/projection architecture is already available, already usable on the terminal interface, and now participates in the restricted lift-back layer whenever only retained-carrier conclusions are drawn;
  \item for Hypothesis~\ref{hyp:frame-dependent-defect-readout}, the still-open strong claim is the existence of a deeper explicit compensator generator $\mathbf H_C$, while the theorem-level shadow separating visible defect vanishing from carrier extinction is already available, already usable on the terminal interface, and now participates in the no-second-carrier / restricted lift-back layer;
  \item for Hypothesis~\ref{hyp:fluid-readout}, the still-open strong claim is the existence of a concrete coarse-graining map to fluid-type evolution, while the theorem-level shadow demoting fluid-like behavior to a later coarse-grained readout is already available, already usable on the terminal interface, and now participates negatively in the restricted lift-back layer by forbidding replacement of the retained carrier;
  \item for Hypothesis~\ref{hyp:bidirectional-probe-lensing}, the still-open strong claim is a full universal geometric lensing law, while the theorem-level shadow of common-carrier/two-probe coupling is already available, already usable on the terminal interface, and now participates in the no-second-carrier / partial-closure layer;
  \item for the OP~3/OP~5 line, the still-open strong claim is a full global closure theorem, while the historical public v3.23.0 release line already contained more than a weak bridge: it had a first partial closure theorem together with a restricted lift-back to retained carrier-layer structure.
\end{enumerate}
Therefore these upgraded blocks should no longer be counted among the \emph{fully open foundational fronts} of the document. Only their strong deferred forms remain fully open; their weaker theorematic, interface-level, and partial-closure forms are already part of the active line.
\end{proposition}

\begin{proof}
Each clause follows by combining the public v3.20 theorem-level shadows, the public v3.21 interface-level application principles, and where applicable the new v3.22 partial-closure consequences. For the HC meta-field generator hypothesis, this combines the common-carrier/projection shadow and its interface-level use with the present restricted lift-back layer. For frame-dependent defect readout, this combines the defect-shadow theorem, the interface-level defect-face vanishing principle, and the present no-second-carrier/lift-back layer. For fluid-type readout, this combines the fluid-shadow theorem, the interface-level fluid demotion principle, and the present lift-back restriction that forbids carrier replacement. For bidirectional probe lensing, this combines the probe-lensing shadow theorem, the interface-level probe-coupling principle, and the present no-second-carrier/partial-closure layer. Finally, the OP~3/OP~5 line is upgraded from the v3.21 bridge test to the present v3.22 partial closure theorem and restricted lift-back theorem. Hence that historical stage already supported a synchronized four-tier status reading.
\end{proof}

\begin{corollary}[Use rule for the synchronized v3.22 status reading]
\label{cor:use-rule-for-the-synchronized-v322-status-reading}
Whenever a later argument refers back to the historical public v3.23.0 release line and needs to classify whether one of the upgraded closure/interface blocks was still truly open at that stage, it is sufficient to cite Proposition~\ref{prop:v322-hypothesis-open-question-status-synchronization-principle}. A finer reopening is required only when the argument depends on the still-open strong form rather than on the already stabilized theorematic, interface-level, or restricted partial-closure form.
\end{corollary}

\begin{proof}
This is the operational reformulation of Proposition~\ref{prop:v322-hypothesis-open-question-status-synchronization-principle}. The proposition already isolates the precise point at which the upgraded blocks cease to count as fully open foundational fronts.
\end{proof}


\begin{remark}[v3.21 synchronized status reading for the upgraded hypothesis layer]
\label{rem:v321-synchronized-status-reading-upgraded-hypothesis-layer}
The present v\docversion\ release line no longer permits a flat reading in which the upgraded hypothesis blocks are simply either ``fully open'' or ``already proved.'' After the released v3.20 theorem-level shadow package and the historical v3.21 interface payoff layer, their status is now genuinely three-tiered:
\begin{enumerate}[leftmargin=2em,label=(\roman*)]
  \item a \emph{strong deferred claim} remains open;
  \item a \emph{theorem-level shadow} is already available;
  \item an \emph{interface-level operational use} is already available on the terminal surface.
\end{enumerate}
Thus the open-front layer must now distinguish between still-open strong forms and already stabilized weaker forms.
\end{remark}

\begin{proposition}[v3.21 hypothesis/open-question status synchronization principle]
\label{prop:v321-hypothesis-open-question-status-synchronization-principle}
At the present theorematic and interface level, the following synchronized status reading holds:
\begin{enumerate}[leftmargin=2em,label=(\roman*)]
  \item for Hypothesis~\ref{hyp:hc-meta-generator}, the still-open strong claim is the existence of one explicit compensator-centered generating equation $\mathfrak E_{\mathrm{HC}}=0$, while the theorem-level shadow of a common closure-stable carrier/projection architecture is already available and already usable on the terminal interface;
  \item for Hypothesis~\ref{hyp:frame-dependent-defect-readout}, the still-open strong claim is the existence of a deeper explicit compensator generator $\mathbf H_C$, while the theorem-level shadow separating visible defect vanishing from carrier extinction is already available and already usable on the terminal interface;
  \item for Hypothesis~\ref{hyp:fluid-readout}, the still-open strong claim is the existence of a concrete coarse-graining map to fluid-type evolution, while the theorem-level shadow demoting fluid-like behavior to a later coarse-grained readout is already available and already usable on the terminal interface;
  \item for Hypothesis~\ref{hyp:bidirectional-probe-lensing}, the still-open strong claim is a full universal geometric lensing law, while the theorem-level shadow of common-carrier/two-probe coupling is already available and already usable on the terminal interface.
\end{enumerate}
Therefore these four blocks should no longer be counted among the \emph{fully open foundational fronts} of the document. Only their strong deferred forms remain open; their weaker theorematic and interface-operational forms are already part of the active line.
\end{proposition}

\begin{proof}
Each clause follows by combining one theorem-level shadow result with its corresponding interface-level application principle. For Hypothesis~\ref{hyp:hc-meta-generator}, this is the theorem-level shadow of the HC meta-field generator together with the interface-level common-carrier/projection principle. For Hypothesis~\ref{hyp:frame-dependent-defect-readout}, this is the theorem-level shadow of frame-dependent defect readout together with the interface-level defect-face vanishing principle. For Hypothesis~\ref{hyp:fluid-readout}, this is the theorem-level shadow of fluid-type readout together with the interface-level fluid-readout demotion principle. For Hypothesis~\ref{hyp:bidirectional-probe-lensing}, this is the theorem-level shadow of probe-lensing together with the interface-level probe-coupling principle. Hence the four upgraded blocks now share one synchronized three-tier status reading: strong deferred form, theorem-level shadow, and interface-level operational use.
\end{proof}

\begin{corollary}[Use rule for the synchronized v3.21 status reading]
\label{cor:use-rule-for-the-synchronized-v321-status-reading}
Whenever a later argument in the present v\docversion\ release line needs to classify whether one of the four upgraded hypothesis blocks is still truly open, it is sufficient to cite Proposition~\ref{prop:v321-hypothesis-open-question-status-synchronization-principle}. A finer reopening is required only when the argument depends on the still-open strong form rather than on the already stabilized theorematic or interface-level shadow.
\end{corollary}

\begin{proof}
This is the operational reformulation of Proposition~\ref{prop:v321-hypothesis-open-question-status-synchronization-principle}. The proposition already isolates the precise point at which the upgraded blocks cease to count as fully open foundational fronts.
\end{proof}


\begin{enumerate}
  \item \textbf{Residual quantitative completion of the frame scale}
        \textbf{(historical OP~1; also absorbing OP~2(T4) and the former OP~4 coupling).}
        The continuum fixed-point structure is already identified at the
        Hopf/Tribonacci level, and the frame-offset problem is no longer a
        broad undifferentiated open front. What remains open is the explicit
        derivation of the prime-lattice / continuum-limit correction
        $\delta_{\mathrm{CL}}$ (or an equivalent arithmetic correction)
        that transports the continuum candidate to the measured matter-frame
        value.

  \item \textbf{Phenomenology: controlled emergence.}
        The probe--closure field system is now written in explicit
        variational form and closed by the residue constraint
        $\mathfrak R_C=0$.
        Structurally, the preceding sections identify Dirac-, QFT-, and
        Einstein-type regimes as distinguished projections of the same
        closure-stable equation. What remains open is the fully controlled
        derivation of these projections with explicit constants,
        approximation bounds, and quantitative phenomenology.

  \item \textbf{Spectrum: mass hierarchy.}
        The present structural reading is already strong enough to separate the
        qualitative mass-sector logic from the unresolved quantitative
        assignment. The remaining open task is the map from prime/spectral
        index to Standard-Model quantum numbers and measured mass scales.

  \item \textbf{Spectrum: three generations.}
        $S^2\supset S^1\supset S^0$ as three fermion generations is
        structurally supported, but the quantitative generation-splitting law
        is not yet derived.

  \item \textbf{Structural completion: Dixon algebra.}
        A compact reformulation in
        $\mathbb{C}\otimes\HH\otimes\OO$ (64 dimensions) may unify
        the present results into one algebraic object, but this remains a
        later consolidation step rather than a current prerequisite.

  \item \textbf{$\Lambda$-cancellation: reduced photon-frame residual problem}
        \textbf{(historical OP~5, Step~(B) only).}
        The photon-frame formulation and the structural separation
        $\Lambda_{\mathrm{fundamental}}=0$ versus quantitative balance are now
        fixed, and Step~(A) of the $\Theta$--$R$ program is already closed in
        the present text. The remaining open content is Step~(B): the exact
        $\Theta$-sum / $\Theta$--$R$ balance needed to cancel the constant term
        at the photon-frame cutoff.
\end{enumerate}

\paragraph{Historical OP synchronization.}
For the remainder of the document, the historical labels should be read with
these statuses in mind:
\begin{itemize}[leftmargin=2em]
  \item \textbf{Historical OP~3 (Koide / $\theta_K$ block):} closed at the
        present structural level; later work may still densify exposition, but
        it should no longer be listed as a primary open problem.
  \item \textbf{Historical OP~2 (Dirac operator on $\Hplusminus$):} closed as
        an independent operator problem in its existence, self-adjointness,
        and GOE-class parts; only the former frame-offset part survives, and
        it is absorbed into Item~1 above.
  \item \textbf{Historical OP~5:} reduced to Step~(B) only; Step~(A) is no
        longer open.
\end{itemize}




\subsection*{Phase IV-I: Physical Unfolding, Observer Frames, and Interface Regimes}
\addcontentsline{toc}{subsection}{Phase IV-I: Physical Unfolding, Observer Frames, and Interface Regimes}

\begin{lemma}[Observer-frame lemma]
\label{lem:observer-frame-lemma}
Within the normalized structural architecture, an observer frame is not introduced as an external background datum. Rather, it is the stabilized readout frame selected by the combined action of the cascade, the closure filter, and the admissibility-preserving transport across the grouped layers. In particular, any observer description compatible with the kernel must arise as a frame-shifted but ground-linked reading of the same underlying body.
\end{lemma}

\begin{proof}
The semantic kernel identifies all admissible descriptions as readings of a single underlying structure. Frame-shifted observation therefore does not create an ontologically new sector. By the forcing chain established in Phases~I--III, any admissible readout is first re-identified by the compensator, then filtered into a closure-stable class, then transported through the grouped architecture without ontology drift. The observer frame is therefore fixed only at the level of admissible readout, not posited independently of the structure.
\end{proof}

\begin{lemma}[Mass/spectral reading lemma]
\label{lem:mass-spectral-reading-lemma}
Any admissible mass or hierarchy reading in the Companion must be interpreted first as a spectral or sectoral readout of the normalized structural kernel, and only secondarily as a phenomenological mass assignment. In particular, the transition from spectral organization to mass hierarchy is a realization map, not a new generative axiom.
\end{lemma}

\begin{proof}
The generative layer produces admissible sectors and their closure relations, while the projection/translation layer supplies readout rules and collapse-fixed identifications. Therefore any mass reading that appears in the field/probe layer inherits its legitimacy from a prior spectral or sectoral differentiation. This excludes the reverse move of inserting an independent mass ontology into the generative layer. The mass interpretation is thus subordinate to the spectral structure and must be treated as a downstream realization.
\end{proof}

\begin{lemma}[Generation-splitting lemma]
\label{lem:generation-splitting-lemma}
If the triolic and grouped-layer architecture is preserved across admissible projections and field realizations, then any generation splitting must appear as a structured decomposition of an already admissible spectral body rather than as a free multiplication of unrelated sectors.
\end{lemma}

\begin{proof}
Triolic minimality and orthogonal anchoring prevent arbitrary proliferation of sectors at the generative level. The transport rules for $\mathcal M$ preserve admissibility and closure, and the no-ontology-drift lemma excludes the creation of new bodies during layer transfer. Hence any multiplicity such as a generation-like splitting must be read as an organized decomposition of the inherited spectral structure. It is therefore constrained by the prior closure and transport logic.
\end{proof}

\begin{theorem}[Standard-model interface theorem]
\label{thm:standard-model-interface-theorem}
The normalized Companion architecture defines a one-way interface to Standard-Model-style phenomenology: admissible structures may be projected into observer, spectral, and field realizations compatible with particle-physics language, but the Companion introduces no independent Standard-Model ontology at the generative level. Consequently, the appropriate first target of the unfolding is not a complete reconstruction of the Standard Model, but a structurally controlled interface map from the grouped kernel into physically readable sectors.
\end{theorem}

\begin{proof}
By the observer-frame, mass/spectral, and generation-splitting lemmas, all physically readable sectors arise only after the generative core has been passed through closure, confinement, cascade readout, tensor realization, and grouped transport. The field/probe layer therefore supplies an interface regime in which known physical language may be attached to admissible sectors, but the kernel itself remains prior to that language. The Standard Model can thus enter only through an interface map respecting admissibility, closure, and no-ontology drift. This proves that the correct unfolding target is an interface theorem rather than a premature full identification.
\end{proof}

\begin{corollary}[First physical unfolding chain]
\label{cor:first-physical-unfolding-chain}
The first physical unfolding chain of the Companion is
\[
\begin{aligned}
&\text{observer-frame emergence}
  \Longrightarrow \text{spectral/mass readout}\\
&\text{generation-compatible splitting}
  \Longrightarrow \text{controlled Standard-Model interface}.
\end{aligned}
\]
Thus the physical interpretation layer is admitted only after the structural, closure, tensorial, and grouped-layer chains have already been stabilized.
\end{corollary}


\subsection*{Phase IV-I$'$: Refined Physical Interface and Realization Control}
\addcontentsline{toc}{subsection}{Phase IV-I': Refined Physical Interface and Realization Control}

\begin{lemma}[Observer covariance under admissible transport]
\label{lem:observer-covariance-admissible-transport}
If two observer descriptions are related by admissibility-preserving transport through the grouped architecture, then they represent frame-shifted realizations of the same underlying body and cannot define distinct generative ontologies.
\end{lemma}

\begin{proof}
The no-ontology-drift result of Phase~II-F implies that transport across $\mathfrak G$, $\mathfrak P$, and $\mathfrak F$ preserves the underlying structural identity. By the observer-frame lemma, any observer description is already a frame-shifted readout of the same admissible kernel. Therefore two observer descriptions related by admissible transport differ only at the level of readout covariance and not at the level of generative content.
\end{proof}

\begin{proposition}[Spectral-to-mass realization map]
\label{prop:spectral-to-mass-realization-map}
There exists a structurally controlled realization map
\[
\mathcal R_{m}: \mathcal S_{\mathrm{spec}} \longrightarrow \mathcal S_{\mathrm{mass}}
\]
from admissible spectral sectors to admissible mass-readout sectors, where $\mathcal R_{m}$ is downstream of closure, cascade readout, and grouped transport.
\end{proposition}

\begin{proof}
By the mass/spectral reading lemma, any mass assignment must be subordinated to a prior spectral organization. The admissible spectral sectors are produced only after the closure-stable body has been transported through the projection/translation layer and read out through the cascade. Hence the passage to a mass language is not primitive but realized by a downstream map from already admissible spectral sectors. Denote this realization map by $\mathcal R_m$. Its domain is therefore restricted to $\mathcal S_{\mathrm{spec}}$ and its codomain inherits admissibility from the upstream structural chain.
\end{proof}

\begin{proposition}[Generation-compatible decomposition]
\label{prop:generation-compatible-decomposition}
If a generation-like multiplicity appears in the physical layer, then it must be expressible as an admissible decomposition of the image of $\mathcal R_m$ rather than as an unconstrained replication of primitive sectors.
\end{proposition}

\begin{proof}
The generation-splitting lemma already excludes free multiplication of unrelated sectors. Since the physical layer is reached only after grouped transport and spectral-to-mass realization, any multiplicity that survives into the interface regime must be inherited from the already admissible image of $\mathcal R_m$. Thus a generation-like structure can only be accepted when it is realizable as a constrained decomposition of an upstream admissible spectral body.
\end{proof}

\begin{theorem}[Controlled physical interface functor]
\label{thm:controlled-physical-interface-functor}
The first physically readable regime of the Companion is determined by a controlled interface functor
\[
\mathfrak I_{\mathrm{phys}}:
(\mathfrak G,\mathfrak P,\mathfrak F)
\longrightarrow
\mathcal O_{\mathrm{phys}},
\]
which assigns observer, spectral, mass, and generation-level readouts to admissible structural sectors while preserving closure compatibility and ontology non-drift.
\end{theorem}

\begin{proof}
The observer covariance lemma identifies physically readable frames as transport-compatible readouts of the same body. The spectral-to-mass realization proposition supplies a downstream realization map from spectral sectors to mass sectors, while the generation-compatible decomposition proposition constrains multiplicity to inherited admissible decompositions. Collecting these assignments yields a functorial interface from the grouped structural architecture to a space of physically readable observables $\mathcal O_{\mathrm{phys}}$. Since every step is subordinated to admissibility, closure, and no-ontology drift, the resulting interface is controlled rather than ontologically generative.
\end{proof}

\begin{corollary}[Refined physical interface chain]
\label{cor:refined-physical-interface-chain}
The refined physical interface chain is
\[
\begin{aligned}
&\text{observer covariance}
  \Longrightarrow \text{spectral-to-mass realization}\\
&\text{generation-compatible decomposition}
  \Longrightarrow \text{controlled interface functor}.
\end{aligned}
\]
Hence the physical language of the Companion is introduced only as a regulated readout of the normalized structural architecture.
\end{corollary}

%% ============================================================
%% APPENDIX A — Formal Proof Supplement
%% ============================================================







\subsection*{Phase IV-J: Spectral Readout, Trace Reconstruction, and Observable Selection}
\addcontentsline{toc}{subsection}{Phase IV-J: Spectral Readout, Trace Reconstruction, and Observable Selection}

\begin{lemma}[Spectral admissibility inheritance]
\label{lem:spectral-admissibility-inheritance}
If a physically readable sector is admitted by the controlled interface functor, then its spectral representative must inherit admissibility from the upstream closure-stable and transport-compatible structural body. In particular, no spectral mode may enter the readout regime unless it descends from an already admissible structural sector.
\end{lemma}

\begin{proof}
By the semantic--tensor interface and the grouped transport chain, admissibility is preserved from the generative layer through cascade readout and into the physically readable interface regime. The controlled physical interface functor therefore cannot attach an observable meaning to a mode that was not already admitted upstream. Hence spectral readability inherits admissibility rather than creating it locally.
\end{proof}

\begin{proposition}[Trace-compatible spectral decomposition]
\label{prop:trace-compatible-spectral-decomposition}
Every admissible spectral body in the readout regime admits a trace-compatible decomposition into modes whose contribution respects the closure filter, the grouped transport constraints, and the observer-covariant interface conditions.
\end{proposition}

\begin{proof}
The historical spectral and trace cores already isolate the need for admissible mode selection and trace-compatible decomposition. In the normalized architecture, these requirements are now downstream of closure, cascade, and grouped transport. Therefore any accepted spectral decomposition must preserve the admissible body identified by the upstream chain and cannot separate into trace-incompatible or closure-breaking components. This yields a trace-compatible mode decomposition of every admissible spectral body.
\end{proof}

\begin{proposition}[Observable extraction map]
\label{prop:observable-extraction-map}
There exists a structurally controlled observable extraction map
\[
\mathcal R_{\mathrm{obs}}: \mathcal S_{\mathrm{spec}}^{\mathrm{adm}} \longrightarrow \mathcal O_{\mathrm{phys}}^{\mathrm{adm}},
\]
which assigns observable readouts only to trace-compatible admissible spectral sectors.
\end{proposition}

\begin{proof}
Phase~IV-I$'$ introduced a controlled interface functor from the grouped architecture into physically readable observables. Restricting that interface to the spectrally admissible, trace-compatible body provided by Proposition~\ref{prop:trace-compatible-spectral-decomposition} yields a downstream extraction map from admissible spectral sectors to admissible observables. Since its domain is already filtered by closure, transport, and trace compatibility, the map does not generate new observables freely but only extracts structurally sanctioned ones.
\end{proof}

\begin{theorem}[Mode-selection and observable consistency theorem]
\label{thm:mode-selection-observable-consistency}
Mode selection in the Companion is admissible only when it is simultaneously spectral, trace-compatible, closure-preserving, and observer-covariant. Under these conditions, the observable extraction map produces a consistent readout regime in which selected modes can be interpreted without violating ontology non-drift.
\end{theorem}

\begin{proof}
Spectral admissibility inheritance excludes modes that do not descend from the admissible structural chain. Trace-compatible spectral decomposition excludes decompositions that would destroy the integrity of the admissible spectral body. The controlled interface functor and observer covariance from Phase~IV-I$'$ ensure that the resulting readout remains a covariance of one and the same body rather than a proliferation of ontologies. Therefore mode selection is consistent exactly when all four constraints are met, and in that regime observable extraction remains structurally controlled.
\end{proof}

\begin{corollary}[Spectral readout reconstruction chain]
\label{cor:spectral-readout-reconstruction-chain}
The next spectral unfolding chain is
\[
\begin{aligned}
&\text{spectral admissibility inheritance}
\Longrightarrow
\text{trace-compatible decomposition}\\
&\Longrightarrow
\text{observable extraction}
\Longrightarrow
\text{consistent mode-selected readout}.
\end{aligned}
\]
Thus the spectral and trace layer of the Companion is admitted only after closure, grouped transport, and controlled physical interfacing have already been stabilized.
\end{corollary}


% -- relocated late main-text phases IV-K to IV-O --


\subsection*{Phase IV-K: Resolvent Gatekeeping and Operator-Controlled Spectral Access}
\addcontentsline{toc}{subsection}{Phase IV-K: Resolvent Gatekeeping and Operator-Controlled Spectral Access}

\begin{lemma}[Resolvent admissibility lemma]\label{lem:resolvent_admissibility}
Let $\mathcal R(\lambda)=(\lambda\mathbf 1-\mathcal A)^{-1}$ denote a formal resolvent on the admissible spectral layer. If the underlying mode family is closure-stable and trace-compatible, then resolvent access is admissible only on those channels for which the induced transport preserves the grouped-layer constraints and does not reintroduce ontology drift. In particular, admissible resolvent access is a restricted rather than global spectral operation.
\end{lemma}
\begin{proof}
The semantic kernel already constrains admissible transport by closure preservation, grouped transport, and controlled readout. A resolvent therefore cannot be treated as an arbitrary inversion device: it is allowed only where the corresponding operator action remains inside the admissible closure class and respects the readout filter established in the preceding phases. Hence admissibility is inherited by resolvent access only on the structurally licensed channels.
\end{proof}

\begin{proposition}[Gatekeeping proposition for admissible modes]\label{prop:resolvent_gatekeeping}
The resolvent layer acts as a gatekeeping mechanism: admissible modes are precisely those spectral branches for which resolvent access is compatible with closure, trace reconstruction, and grouped-layer transport. Non-admissible branches are excluded before observable extraction.
\end{proposition}
\begin{proof}
By Lemma~\ref{lem:resolvent_admissibility}, resolvent access is already restricted to structure-preserving channels. Combined with the trace-compatible mode selection of Phase~IV-J, this yields a gatekeeping principle: only those modes that remain compatible with closure and transport can pass into the observable sector. The exclusion of non-admissible branches is therefore structural rather than ad hoc.
\end{proof}

\begin{theorem}[Operator-controlled trace access theorem]\label{thm:operator_controlled_trace_access}
Assume that the spectral sector is admissible, the selected modes are resolvent-gated, and the trace reconstruction remains compatible with the semantic--tensor interface. Then trace access is operator-controlled: every observable trace contribution arises through an admissible operator channel and not through uncontrolled spectral leakage.
\end{theorem}
\begin{proof}
The preceding proposition ensures that only admissible modes are available for readout. The semantic--tensor interface then identifies the operator realization of the same admissible content, while closure preservation prevents leakage into non-licensed channels. Consequently, the trace contributions that survive observable extraction are exactly those transmitted through operator-controlled admissible channels.
\end{proof}

\begin{corollary}[Stable observable channel corollary]\label{cor:stable_observable_channel}
Under admissible resolvent gatekeeping, the observable readout sector forms a stable channel: spectral access, trace reconstruction, and mode selection remain synchronized under the same structural admissibility constraints.
\end{corollary}
\begin{proof}
This follows directly from Theorem~\ref{thm:operator_controlled_trace_access}: once admissible access, admissible mode selection, and admissible trace transport coincide, the resulting observable sector is stable under the same structural constraints.
\end{proof}

\begin{remark}[Resolvent-to-readout chain]\label{rem:resolvent_to_readout_chain}
The operator-controlled spectral access chain may therefore be summarized as
\[
\begin{aligned}
\text{resolvent admissibility}
&\Longrightarrow \text{gatekept mode access}\\
&\Longrightarrow \text{operator-controlled trace access}\\
&\Longrightarrow \text{stable observable channel}.
\end{aligned}
\]
This phase extends the preceding spectral readout layer by showing that observable extraction is not only spectrally filtered but also operatorically regulated.
\end{remark}



\subsection*{Phase IV-L: Semigroup Stability and Lyapunov-Controlled Readout Persistence}
\addcontentsline{toc}{subsection}{Phase IV-L: Semigroup Stability and Lyapunov-Controlled Readout Persistence}

\begin{lemma}[Semigroup admissibility lemma]\label{lem:semigroup_admissibility}
Let $(T_t)_{t\ge 0}$ denote a formal semigroup acting on the admissible readout sector. If the initial spectral channel is closure-stable, resolvent-gated, and trace-compatible, then admissible semigroup evolution is restricted to those trajectories for which grouped transport, closure preservation, and observable compatibility remain jointly valid for all admissible times.
\end{lemma}
\begin{proof}
The preceding phases identify the observable channel only after admissibility, gatekeeping, and operator-controlled trace access have been enforced. A semigroup acting on this sector may therefore propagate states only if the induced time evolution respects the same admissibility filter. Otherwise the evolution would exit the closure-preserving and trace-compatible regime and would no longer represent a structurally licensed readout trajectory. Hence admissible semigroup evolution is constrained to the channels that preserve the previously established structural conditions for all admissible times.
\end{proof}

\begin{proposition}[Lyapunov gate proposition]\label{prop:lyapunov_gate}
Assume that an admissible Lyapunov functional $\mathcal L$ exists on the gated readout sector. If $\mathcal L$ is nonincreasing along the admissible semigroup trajectories, then unstable leakage into non-admissible channels is structurally suppressed, and the gated observable channel remains dynamically protected.
\end{proposition}
\begin{proof}
A nonincreasing Lyapunov functional excludes the spontaneous amplification of perturbations that would carry the evolution outside the admissible channel. Since admissibility already forbids ontology drift, closure violation, and non-gated spectral leakage, Lyapunov monotonicity upgrades these static constraints to a dynamical protection principle. The readout channel is therefore not merely selected once but remains protected under admissible evolution.
\end{proof}

\begin{theorem}[Persistent observable channel theorem]\label{thm:persistent_observable_channel}
Suppose that the readout sector is resolvent-gated, operator-controlled, and semigroup-admissible, and that a Lyapunov functional controls admissible perturbations. Then the observable channel is persistent: admissible mode selection, trace reconstruction, and grouped transport remain synchronized along the admissible evolution, and the readout sector stays dynamically stable against structural leakage.
\end{theorem}
\begin{proof}
By Lemma~\ref{lem:semigroup_admissibility}, admissible semigroup evolution preserves the structurally licensed readout sector. Proposition~\ref{prop:lyapunov_gate} then prevents the admissible trajectories from developing unstable excursions into non-gated channels. Combined with the operator-controlled access theorem and the stable observable channel corollary from Phase~IV-K, this yields persistence of the observable channel: the same admissibility constraints that select the readout sector also stabilize it under admissible evolution.
\end{proof}

\begin{corollary}[Long-time readout stability corollary]\label{cor:long_time_readout_stability}
Under admissible semigroup evolution and Lyapunov control, the observable readout channel remains stable at long admissible times; the selected spectral and trace-compatible sector is therefore not merely accessible but persistently readable.
\end{corollary}
\begin{proof}
This follows directly from Theorem~\ref{thm:persistent_observable_channel}: persistence of the admissible channel implies long-time stability of the same gated and closure-compatible readout sector.
\end{proof}

\begin{remark}[Semigroup-to-stability chain]\label{rem:semigroup_to_stability_chain}
The late-stage stability chain may therefore be summarized as
\[
\begin{aligned}
\text{semigroup admissibility}
&\Longrightarrow \text{Lyapunov-controlled gated evolution}\\
&\Longrightarrow \text{persistent observable channel}\\
&\Longrightarrow \text{long-time readout stability}.
\end{aligned}
\]
This phase completes the operatoric readout picture by showing that admissible channels are not only spectrally and resolvent-gated, but also dynamically stabilized over admissible evolution scales.
\end{remark}



\subsection*{Phase IV-M: Channel Fusion and Multi-Readout Consistency}
\addcontentsline{toc}{subsection}{Phase IV-M: Channel Fusion and Multi-Readout Consistency}

\begin{lemma}[Admissible channel fusion lemma]\label{lem:admissible_channel_fusion}
Suppose that two or more observable channels are individually admissible, resolvent-gated, and semigroup-stable. Then a fused readout may be admitted only if the fusion procedure preserves the common closure filter, the grouped transport discipline, and the no-ontology-drift condition. In particular, fusion is admissible only as a compatibility-preserving synthesis of already licensed channels.
\end{lemma}
\begin{proof}
The previous phases guarantee admissibility, operatoric access, and dynamical persistence for each channel separately. A fusion mechanism cannot be treated as an unrestricted superposition principle, since otherwise it could reintroduce hidden drift between channels that were only separately controlled. Therefore fusion is admissible only when the same closure, transport, and ontology constraints remain valid for the synthesized channel. In that case fusion does not create a new structural body, but re-expresses a common admissible body through multiple synchronized readouts.
\end{proof}

\begin{proposition}[Cross-channel consistency proposition]\label{prop:cross_channel_consistency}
If admissible channels are fused through a compatibility-preserving synthesis, then their spectral, trace-based, and long-time readout contents must agree on the shared admissible sector. Any discrepancy beyond the admitted transport freedom signals a non-admissible fusion rather than a genuine multi-readout effect.
\end{proposition}
\begin{proof}
Each admissible channel already carries its observable content through the same semantic--tensor interface, the same resolvent gatekeeping, and the same semigroup-stable evolution class. Hence a compatibility-preserving fusion may only combine channels whose readout contents coincide on the overlap of their admissible sectors. If an apparent fusion produces incompatible observable contents beyond the allowed transport freedom, then the synthesis has violated the common admissibility constraints and is therefore not structurally licensed.
\end{proof}

\begin{theorem}[Multi-readout consistency theorem]\label{thm:multi_readout_consistency}
Under admissible channel fusion, the Companion admits a controlled multi-readout regime: distinct readout channels may coexist, but only as synchronized presentations of one and the same admissible structural body. Consequently, multi-readout consistency is equivalent to fusion without ontology drift.
\end{theorem}
\begin{proof}
By Lemma~\ref{lem:admissible_channel_fusion}, channel fusion is allowed only as a synthesis that preserves closure, grouped transport, and non-drift. Proposition~\ref{prop:cross_channel_consistency} then forces agreement of spectral, trace, and persistent readout content on the common admissible sector. Therefore multiple channels can coexist only as mutually compatible renderings of the same admissible body. The controlled multi-readout regime is thus precisely the regime in which fusion introduces no ontology drift.
\end{proof}

\begin{corollary}[Fusion-to-consistency chain]\label{cor:fusion_to_consistency_chain}
The late-stage fusion chain may therefore be summarized as
\[
\begin{aligned}
\text{admissible channel fusion}
&\Longrightarrow \text{cross-channel consistency}\\
&\Longrightarrow \text{controlled multi-readout regime}\\
&\Longrightarrow \text{fusion without ontology drift}.
\end{aligned}
\]
This phase closes the late readout unfolding by showing that stable observable channels may be combined only through a synthesis that preserves the unique admissible structural body across all readout presentations.
\end{corollary}


\subsection*{Phase IV-N: Unified Observable Readout Synthesis}
\addcontentsline{toc}{subsection}{Phase IV-N: Unified Observable Readout Synthesis}

\begin{proposition}[Unified observable readout synthesis]
\label{prop:unified-observable-readout}
Let the late readout regime be governed by the spectral admissibility, resolvent gatekeeping, semigroup persistence, and admissible channel-fusion principles established in Phases IV-J through IV-M. Then the resulting observable sector admits a unified synthesis map in which spectral access, operator control, long-time persistence, and multi-readout consistency are read as mutually compatible aspects of one and the same admissible observable architecture.
\end{proposition}

\begin{proof}
Phase IV-J supplies the spectral and trace-compatible decomposition of admissible modes. Phase IV-K restricts this decomposition to resolvent-accessible channels only, so that observable extraction is not performed on arbitrary modes but only on structurally permitted ones. Phase IV-L then upgrades admissible access to persistence under semigroup evolution and Lyapunov control, which turns transient access into stable readout inheritance over time. Finally, Phase IV-M shows that multiple such persistent channels can be fused without inducing cross-channel inconsistency or ontological drift. Therefore the late readout regime is not a mere juxtaposition of separate spectral, operator, dynamical, and fusion statements; it forms a single admissible synthesis in which observable extraction is structurally unified.\qedhere
\end{proof}

\begin{corollary}[Late readout master chain]
The late readout architecture admits the compressed chain
\[
\begin{aligned}
&\text{spectral admissibility}
\Longrightarrow \text{resolvent-controlled access}\\
&\Longrightarrow \text{semigroup/Lyapunov persistence}
\Longrightarrow \text{multi-readout consistency}\\
&\Longrightarrow \text{unified observable synthesis.}
\end{aligned}
\]
\end{corollary}

\begin{remark}[Interpretive role of the late synthesis]
This synthesis does not introduce a new ontology behind the observable layer. Rather, it states that the different late readout mechanisms are to be interpreted as coordinated access constraints and persistence conditions on one admissible observable sector generated by the normalized structural core.
\end{remark}



\subsection*{Phase IV-O: Return Bridge to the Generative Core and the Structural Master Equation}
\addcontentsline{toc}{subsection}{Phase IV-O: Return Bridge to the Generative Core and the Structural Master Equation}

\begin{proposition}[Return bridge from unified readout to the generative core]\label{prop:return_bridge_unified_readout}
The unified observable readout synthesized in Phase IV-N introduces no new generative assumption. Rather, it is a late readout layer of the same generative structure determined by the one-body condition, triolic stabilization, orthogonal anchoring, frame-shifted observation, and the HC--closure--confinement chain.
\end{proposition}

\begin{proof}
The earlier unfolding phases establish that admissible readout channels arise only after re-identification, compensator forcing, closure stabilization, and higher confinement. Phases IV-J through IV-N then refine how admissible observable content is accessed, filtered, stabilized, fused, and synthesized. None of those later steps introduces a new generating source; they remain constrained by the same admissibility conditions inherited from the generative layer. Hence the late observable synthesis is a return readout of the original structure rather than an independent origin.
\end{proof}

\begin{theorem}[Master-equation readout closure]\label{thm:master_equation_readout_closure}
Let the structural master equation be the compressed dynamical expression of the admissible generative core. Then the late readout chain
\[
\begin{aligned}
&\text{spectral admissibility}
\Longrightarrow \text{resolvent-controlled access}
\\
&\Longrightarrow \text{semigroup/Lyapunov persistence}
\Longrightarrow \text{multi-readout consistency}
\\
&\Longrightarrow \text{unified observable synthesis}
\end{aligned}
\]
closes back into the same master-equation-governed structure, in the sense that every admissible late observable synthesis remains a governed readout of the original master-equation sector.
\end{theorem}

\begin{proof}
By the semantic/tensor interface and grouped transport phases, admissibility is preserved under the Millennium transport between the generative, projection, and realization layers. The readout phases IV-J through IV-N add only progressively refined access conditions: spectral admissibility, resolvent gatekeeping, semigroup persistence, and fusion consistency. Since each step is formulated as structure-preserving and closure-compatible, the final observable synthesis remains within the same admissible sector selected by the master equation. Therefore the readout chain closes into the same governing structure and should be interpreted as its late observable face.
\end{proof}

\begin{remark}[Interpretive closure of the unfolding program]
The point of the return bridge is to prevent the late readout machinery from being mistaken for a detached auxiliary theory or a second source. The late phases do not compete with the generative core; they complete its observable unfolding. In this sense the overall program closes as
\[
\begin{aligned}
&\text{generative core}
\Longrightarrow \text{HC/closure/confinement}
\Longrightarrow \text{master equation}
\\
&\Longrightarrow \text{spectral-resolvent-semigroup-fusion synthesis}
\\
&\Longrightarrow \text{return readout of the same structure.}
\end{aligned}
\]
\end{remark}


\subsection*{Phase IV-P: Observable Back-Projection and Kernel Reconstruction}
\addcontentsline{toc}{subsection}{Phase IV-P: Observable Back-Projection and Kernel Reconstruction}

\begin{lemma}[Admissible back-projection lemma]\label{lem:admissible_back_projection}
If a late observable synthesis is admissible in the sense of Phases IV-J through IV-O, then it admits a back-projection into the grouped architecture such that the reconstructed object remains compatible with the same closure, transport, and non-drift constraints that governed the forward readout.
\end{lemma}
\begin{proof}
By Phase IV-N, the unified observable synthesis is not an arbitrary superposition of readout data, but a compatibility-preserving synthesis of spectral, resolvent, semigroup, and fusion constraints. Phase IV-O then shows that this synthesis remains a governed readout of the same master-equation sector. Hence any admissible late synthesis can be projected back through the grouped architecture without leaving the admissible class: what is reconstructed is not a new source, but the same structural body viewed through the inverse admissibility constraints of the readout chain.
\end{proof}

\begin{proposition}[Kernel reconstruction proposition]\label{prop:kernel_reconstruction}
Admissible back-projection reconstructs not a second kernel, but a readout-identified representative of the original generative kernel. In particular, reconstruction is valid only up to the transport, projection, and realization equivalences already built into $\mathfrak G$, $\mathfrak P$, and $\mathfrak F$.
\end{proposition}
\begin{proof}
The grouped transport architecture never licenses ontology drift; it only transports one admissible body across grouped layers. Therefore any inverse passage from late observable synthesis to kernel-level description must be taken modulo the same transport equivalences that governed the forward route. The back-projected object is thus a kernel representative recovered from admissible readout data, not an independent second kernel standing alongside the original generative source.
\end{proof}

\begin{theorem}[Observable-to-kernel closure theorem]\label{thm:observable_to_kernel_closure}
The late unfolding program closes in both directions: every admissible generative structure yields a governed late observable synthesis, and every admissible late observable synthesis yields a back-projected kernel representative of that same structure. Consequently, the Companion admits a controlled observable-to-kernel closure rather than a one-way readout only.
\end{theorem}
\begin{proof}
The forward direction is the content of the unfolding chain from the generative core through the master equation to the unified observable synthesis. By Lemma~\ref{lem:admissible_back_projection}, an admissible late synthesis can be projected back into the grouped architecture while preserving closure, transport, and non-drift. Proposition~\ref{prop:kernel_reconstruction} then shows that the reconstructed object is to be read as a representative of the original generative kernel, not as a second ontological origin. Thus the unfolding closes in both directions: generation governs readout, and admissible readout reconstructs the governed kernel representative.
\end{proof}

\begin{corollary}[Bidirectional closure chain]\label{cor:bidirectional_closure_chain}
The late-stage closure may therefore be compressed as
\[
\begin{aligned}
&\text{generative core}
\Longrightarrow \text{master-equation-governed readout}
\Longrightarrow \text{unified observable synthesis}\\
&\Longrightarrow \text{admissible back-projection}
\Longrightarrow \text{kernel representative of the same structure.}
\end{aligned}
\]
\end{corollary}


\subsection*{Phase IV-Q: Global Unfolding Synthesis and Architectural Compression}
\addcontentsline{toc}{subsection}{Phase IV-Q: Global Unfolding Synthesis and Architectural Compression}

The preceding phases may now be compressed into a single Companion-level statement: the full unfolding chain from the generative kernel to observable back-projection is not a juxtaposition of independent modules, but a single admissibility-controlled architecture with forward, lateral, and return transport. The role of the late synthesis is therefore not to append a second theory or a new axiom block to the generative core, but to expose the full reversible structural path from formation to readout and back.

\begin{proposition}[Global unfolding synthesis]\label{prop:global_unfolding_synthesis}
The Companion unfolding program may be read as one admissibility-preserving architecture
\[
\begin{aligned}
\text{generative formation}
&\Longrightarrow \text{HC/closure/confinement forcing}\\
&\Longrightarrow \text{cascade and readout admissibility}\\
&\Longrightarrow \text{semantic/tensor transport}\\
&\Longrightarrow \text{grouped-layer realization}\\
&\Longrightarrow \text{probe and field realization}\\
&\Longrightarrow \text{physical interface control}\\
&\Longrightarrow \text{spectral/resolvent/semigroup/fusion synthesis}\\
&\Longrightarrow \text{observable back-projection}\\
&\Longrightarrow \text{kernel reconstruction of the same structure}.
\end{aligned}
\]
In particular, the later readout layers do not introduce a second ontological source, but provide a controlled return map into the same structural kernel whose forcing chain generated the admissible dynamics in the first place.
\end{proposition}

\begin{proof}
Phase~I establishes the generative forcing chain from one-body discipline to HC, closure, and higher confinement. Phase~II transports this material through cascade-admissible readout, semantic/tensor realization, and grouped-layer architecture without ontological drift. Phase~III realizes the same admissible structure in probe and field form, culminating in the structural master equation and its exact-closure regime. Phase~IV then unfolds the physical interface, spectral access, resolvent control, semigroup stability, fusion consistency, and unified readout synthesis, before Phase~IV-O and Phase~IV-P return the resulting observables to the master-equation-governed kernel and reconstruct a representative of the same admissible structure. The complete chain is therefore structurally circular only in the admissible sense of closure and readout return, not in the sense of a second independent origin.\qedhere
\end{proof}

\begin{corollary}[Companion architectural compression rule]
Any later compression of the Companion is admissible only if it preserves explicit access to each of the following four structural roles:
\[
\text{formation},\qquad
\text{transport},\qquad
\text{realization/readout},\qquad
\text{return reconstruction}.
\]
A compressed presentation that suppresses one of these roles may remain heuristically useful, but it no longer represents the full unfolding architecture proved in the Companion.
\end{corollary}

\paragraph{Interpretive role of the global synthesis.}
This proposition gives the document a single citeable architectural apex: the Companion does not merely collect historical cores, transport layers, field realizations, and readout mechanisms, but exhibits them as one normalized unfolding from generative structure to controlled observability and admissible back-return.



\subsection*{Phase IV-R: Proof Compression and Structural Admissibility Testing}
\addcontentsline{toc}{subsection}{Phase IV-R: Proof Compression and Structural Admissibility Testing}

The global synthesis may now be turned into a practical proof-saving rule. Once the full unfolding architecture has been exhibited, later arguments need not re-derive every intermediate transport, readout, and return step from scratch. Instead, it becomes sufficient to verify that a proposed construction remains inside the admissible chain already established by Phases~I through~IV-Q.

\begin{lemma}[Structural admissibility test]\label{lem:structural_admissibility_test}
Let $X$ be a proposed extension, readout, or reformulation of the Companion architecture. If $X$ preserves the four structural roles of formation, transport, realization/readout, and return reconstruction, and if $X$ introduces no ontology drift relative to the grouped architecture, then $X$ may be tested against the unfolding program by admissibility checking rather than by rebuilding the entire derivation chain from the beginning.
\end{lemma}
\begin{proof}
By Corollary~[Companion architectural compression rule], any admissible compression or reformulation must preserve access to the four defining structural roles. Since Phase~II-F already excludes ontology drift across grouped layers, and Phase~IV-P returns admissible readout data to a kernel representative of the same structure, a later proposal that respects these constraints remains inside the already proved architecture. The burden of proof is therefore reduced from full regeneration to admissibility testing: one checks whether the proposal preserves the established chain instead of reconstructing the whole chain anew.
\end{proof}

\begin{proposition}[Proof-compression proposition]\label{prop:proof_compression_proposition}
The Companion permits controlled proof compression: whenever a later statement is structurally subordinate to an already proved admissible chain, it may be justified by citing preservation of that chain together with the specific additional admissibility condition that differentiates the later statement.
\end{proposition}
\begin{proof}
Phases~I through~IV-Q provide the explicit full chain from formation to observable back-return. If a later statement does not alter this chain, but only specializes one of its stages---for example by restricting admissible spectral modes, selecting a probe sector, or refining an operator gate---then the full derivation need not be repeated. It is enough to show that the specialization preserves the inherited admissibility constraints and adds only the stated extra condition. The proof is therefore compressed, but not weakened: the omitted steps are inherited from the already established chain.
\end{proof}

\begin{theorem}[Meta-theorem of admissible structural compression]\label{thm:meta_admissible_structural_compression}
Any later expansion of the Companion is mathematically legitimate in compressed form if and only if it can be located inside the established unfolding architecture without breaking formation, transport, realization/readout, or return reconstruction. In particular, a compressed argument is valid exactly when it is a structurally admissible sub-chain of the global unfolding synthesis.
\end{theorem}
\begin{proof}
Necessity follows from the architectural compression rule: if one of the four structural roles is broken, the proposal no longer belongs to the full unfolding architecture and cannot inherit its proof force. Sufficiency follows from Lemma~\ref{lem:structural_admissibility_test} and Proposition~\ref{prop:proof_compression_proposition}: once a proposal is shown to preserve the admissible chain and differs only by an explicit additional admissibility condition, the rest of the proof is inherited from the already established unfolding program. Thus admissible compression is neither heuristic abbreviation nor informal shortcut, but a precise meta-rule derived from the structure of the Companion itself.
\end{proof}

\begin{corollary}[Late meta-compression chain]\label{cor:late_meta_compression_chain}
The full late-stage architecture may therefore be compressed as
\[
\begin{aligned}
&\text{global unfolding synthesis}
\Longrightarrow \text{admissibility testing of later proposals}\\
&\Longrightarrow \text{controlled proof compression}
\Longrightarrow \text{valid structural sub-chain expansions.}
\end{aligned}
\]
\end{corollary}

\paragraph{Interpretive role of the meta-compression theorem.}
This final theorem turns the Companion from a one-off expansion into a reusable proof engine. Later developments no longer need to restart from the beginning whenever they remain inside the admissible architecture; they need only show where they sit in the chain and which extra admissibility condition they impose.



% --- inserted late main-text phases IV-S and IV-T ---

%% ============================================================
%%  s_op1_op4_hopf_ckm.tex
%%  STATUS: PROOF MODE — Erster strukturierter Beweisentwurf
%%
%%  OP 1: β_sym = 1/√3 aus Hopf-Fixpunkt-Geometrie
%%  OP 4: CKM/PMNS aus G₂ Euler-Winkeln (Faktor ≈ 2 schließen)
%%
%%  Marc Brendecke, Bochum, März 2026
%%  Drop-In für quantum_sphaera v2.x (nach Section 11)
%%
%%  SCHLÜSSELEINSICHT: OP 1 und OP 4 sind dasselbe Problem.
%%  Der Faktor ≈ 2 in den Mischungswinkeln (OP 4) ist dieselbe
%%  5.1%-Lücke wie in β_sym (OP 1): beides ist die
%%  Kontinuumslimit-Korrektur der diskreten Primgitter-Geometrie.
%% ============================================================

\subsection*{Phase IV-S: Hopf Fixed-Point Geometry and CKM/PMNS Residuals}
\addcontentsline{toc}{subsection}{Phase IV-S: Hopf Fixed-Point Geometry and CKM/PMNS Residuals}
\label{sec:hopf_ckm}

\begin{remark}[Transition from Phase IV-R to the Hopf/CKM block]\label{rem:transition_r_to_s}
Phase~IV-R turned the late synthesis into a reusable admissibility and proof-compression rule. The present phase uses that rule in a more specialized way: instead of introducing a new generative source, it tests whether the residual Hopf/CKM discrepancies can be read as admissible geometric corrections inside the already normalized operatoric readout architecture.
\end{remark}

\subsection{The single underlying gap}

\begin{remark}[OP 1 and OP 4 are coupled]
The two open problems
\begin{itemize}
\item OP 1: why does $\bsym \approx 0.5493$ rather than $1/\sqrt{3}
  \approx 0.5774$? (gap: $5.1\%$)
\item OP 4: why does the $G_2$-Euler-angle derivation of the Cabibbo
  angle yield $\theta_{12} \approx \arcsin(\gamma_{\mathrm{string}})
  \approx 7.3^\circ$ instead of the observed $\theta_C \approx 13.0^\circ$?
  (gap: factor $\approx 1.78 \approx 2$)
\end{itemize}
share a single origin: the \emph{continuum-limit correction} from the
discrete prime lattice $\Gamma_P$ to the smooth manifold geometry.
\begin{equation}
\boxed{
  \delta_{\mathrm{CL}}
  := \frac{\beta_{\mathrm{sym}}^{\mathrm{measured}}}{\beta_{\mathrm{sym}}^{\mathrm{Hopf}}}
  = \frac{0.5493}{1/\sqrt{3}}
  = 0.5493 \cdot \sqrt{3}
  \approx 0.9515,
  \quad
  \text{so } \delta_{\mathrm{CL}} \approx 1 - 4.85\%.
}
\end{equation}
The ART--QFT computation (Section 10) identifies the same gap as a
$1.4\%$ effect in $\gamma_L$: the ratio
$\gamma_L^{\mathrm{measured}}/\gamma_L^{\mathrm{Hopf}}
= 1.864 / (2/\sqrt{3}) \approx 0.9853$,
i.e.\ a $1.47\%$ shift in $\gamma_L$ propagating to a $4.85\%$ shift
in $\bsym$ via the nonlinear formula
$\bsym = \sqrt{(\gamma_L-1)/(\gamma_L+1)}$.
\end{remark}

%% ============================================================
\subsection{OP 1: \texorpdfstring{$\bsym = 1/\sqrt{3}$}{β\_sym = 1/√3}
  from Hopf fixed-point geometry}
\label{sec:op1}
%% ============================================================

\subsubsection{The Hopf setup (already established)}

From Theorems~11.1--11.2 (Companion v2.19):
\begin{enumerate}
\item The cascade $\KC: S^{15} \to S^8 \to S^4 \to S^2$ identifies
  the photon ground state with the north pole of $S^3 \cong SU(2)$
  ($\theta_H = 0$, $z^- = 0$).
\item The matter observer sits at Hopf angle $\theta_H \approx 66.6^\circ$,
  giving $\bsym = \sin(\theta_H/2) \approx 0.5493$.
\item The triolic angle $2\pi/3$ is a useful geometric comparison, but it
  is not the correct Hopf candidate for $\bsym$; the sign-corrected value
  is discussed in Remark~\ref{rem:hopf-candidate-fix}.
\end{enumerate}

\begin{remark}[Sign-corrected Hopf candidate]
\label{rem:hopf-candidate-fix}
Careful: the Hopf candidate in the open problem list is
$\bsym = 1/\sqrt{3} \approx 0.5774$, which corresponds to
$\sin(\theta_H/2) = 1/\sqrt{3}$, i.e. $\theta_H/2 = \arcsin(1/\sqrt{3}) \approx 35.26^\circ$,
so $\theta_H \approx 70.53^\circ$.
This is \emph{not} the triolic angle $2\pi/3$.

It is instead the \textbf{tetrahedral angle}: on $S^3 \cong SU(2)$,
the regular tetrahedron inscribed in $S^3$ has its vertices at angles
separated by $\arccos(-1/3)$ along great circles. The tetrahedral
vertex at unit distance from the north pole has half-angle
$\arctan(\sqrt{2}) \approx 54.74^\circ$ (the Euler characteristic
angle), and the supplementary Hopf projection gives $\sin(\theta/2)
= 1/\sqrt{3}$.

So the \textbf{Hopf candidate is the tetrahedral fixed point of
$\mathbb{Z}/3\mathbb{Z}$ acting on $S^3$}, not the equator or the
triolic angle.
\end{remark}

\subsubsection[The Z/3Z fixed-point structure]{The \texorpdfstring{$\mathbb{Z}/3\mathbb{Z}$}{Z/3Z} fixed-point structure}

\begin{proposition}[$\mathbb{Z}/3\mathbb{Z}$ orbits on $S^3$]
\label{prop:z3z-orbits}
The triolic $\mathbb{Z}/3\mathbb{Z}$ action on $S^3$ (cyclic
generation permutation $\mathrm{Gen}_1 \to \mathrm{Gen}_2 \to
\mathrm{Gen}_3 \to \mathrm{Gen}_1$ via $G_2$-rotation
$e_4 \to e_5 \to e_6 \to e_4$) has fixed points at:
\begin{itemize}
\item North pole: $\theta_H = 0$ (photon, $\bsym = 0$).
\item South pole: $\theta_H = \pi$ (tachyon, $\bsym = 1$).
\item Tetrahedral orbit: $\cos(\theta_H) = -1/3$, i.e. $\theta_H = \arccos(-1/3) \approx 109.47^\circ$, giving $\sin(\theta_H/2) = \sqrt{(1 - (-1/3))/2} = \sqrt{2/3} \approx 0.8165$. \emph{This is not the candidate.}
\item \textbf{The tetrahedral half-angle orbit:} $\theta_H/2 =
  \arcsin(1/\sqrt{3})$, i.e.\ the unique $\mathbb{Z}/3\mathbb{Z}$
  orbit on $S^3$ whose Hopf projection to $S^2$ lands at the
  tetrahedral vertex $(1/3, 1/3, 1/3)/\|(1,1,1)\|$.
\end{itemize}
\end{proposition}

\begin{proof}
Under the $\mathbb{Z}/3\mathbb{Z}$ rotation $R_{2\pi/3}$ on the
triplet $(e_4, e_5, e_6)$, the Hopf-invariant point of $S^3$
(invariant under both the $U(1)$ fibre and $R_{2\pi/3}$) satisfies
$q = R_{2\pi/3} q$ in quaternion coordinates. Writing
$q = a + b\mathbf{i} + c\mathbf{j} + d\mathbf{k}$ with $a^2 + b^2
+ c^2 + d^2 = 1$, the constraint $R_{2\pi/3} q = q$ forces
$b = c = d$ (equal imaginary parts) and $|b| = 1/\sqrt{3 + a^2/(1-a^2)}$.
At the unit sphere, $a^2 + 3b^2 = 1$ with $b = c = d$ gives
$b^2 = (1 - a^2)/3$. The Hopf projection $\pi: S^3 \to S^2$ maps
$q \mapsto \sin(\theta_H/2)$ where $\cos(\theta_H) = 2a^2 - 1$.
The unique non-polar fixed point satisfies $\sin(\theta_H/2) =
|b|\sqrt{3} = \sqrt{1 - a^2}$; combined with the normalization,
$\sin^2(\theta_H/2) = 1 - a^2 = 3b^2$. The Cauchy--Schwarz minimum
of the action on the orbit gives $a = \sqrt{2/3}$, so
$\sin^2(\theta_H/2) = 1/3$, i.e.\ $\sin(\theta_H/2) = 1/\sqrt{3}$.
\end{proof}

\subsubsection{The continuum-limit correction $\delta_{\mathrm{CL}}$}

The Hopf candidate $\bsym^{\mathrm{Hopf}} = 1/\sqrt{3}$ is derived
from the \emph{continuous} $S^3$ geometry. The actual framework sits
on the discrete prime lattice $\Gamma_P = \{\log p : p \text{ prime}\}$.

\begin{proposition}[Effective closure-side continuum-limit shift]
\label{prop:cl-shift}
The observable continuum-limit correction on the closure side is still
read through a multiplicative factor
\[
  \bsym = \bsym^{\mathrm{Hopf}} \cdot \delta_{\mathrm{CL}},
  \quad
  \delta_{\mathrm{CL}} = \frac{0.5493}{1/\sqrt{3}} \approx 0.9515.
\]
However, in the post-v3.31.0 reading this scalar factor is not treated as
the primitive prime-lattice correction object itself. Instead, it is read
as the \emph{effective closure-side projection} of a deeper coupled
prime-lattice correction datum
\[
  \mathfrak{C}_{\Gamma_P}^{\mathrm{cpl}}
  =
  \bigl(
  \Delta_{\mathrm{CL}},
  \Delta_{\Theta R}(T),
  \rho_A
  \bigr),
\]
whose closure-side readout yields the observed scalar factor
$\delta_{\mathrm{CL}}$.
\end{proposition}

\begin{proof}[Proof sketch --- corrected structural reading]
The continuum-limit shift $\delta_{\mathrm{CL}} \approx 0.9515$ is
numerically established from the observed ratio
$\bsym/(1/\sqrt{3}) \approx 0.9515$.
A natural candidate expression is the ratio of Weyl spectral densities:
\[
  \delta_{\mathrm{CL}}
  \stackrel{?}{=}
  \frac{\rho_{\mathrm{photon}}(\gamma_1)}{\rho_{\mathrm{continuum}}(\gamma_1)}
  = \frac{\log(\gamma_1/(2\pi\gL))/(2\pi\gL)}{\log(\gamma_1/(2\pi))/(2\pi)}.
\]
\emph{Correction (v2.29.1):} Numerical evaluation gives
$\rho_{\mathrm{photon}}(\gamma_1)/\rho_{\mathrm{continuum}}(\gamma_1) \approx 0.124$,
which does \emph{not} reproduce the observed $\delta_{\mathrm{CL}} \approx 0.951$.
The corrected effective fit for the closure-side factor is given in
Remark~\ref{conj:cl-correction}:
$\delta_{\mathrm{CL}} = \exp(-\gamma_1/(4\pi\gamma_L^5)) \approx 0.9512$,
matching to $0.018\%$.

The structural reading, however, is no longer purely scalar. On the
current closure-phase transfer line, the first-principles prime-lattice
correction is read through the coupled datum
\[
  \mathfrak{C}_{\Gamma_P}^{\mathrm{cpl}}
  =
  \bigl(
  \Delta_{\mathrm{CL}},
  \Delta_{\Theta R}(T),
  \rho_A
  \bigr),
\]
with $\delta_{\mathrm{CL}}$ interpreted as its effective closure-side
projection. Thus the exponential formula remains a valid observed fit,
while its deeper origin is sought in the coupled closure--Theta correction
problem under admissible packet-side control.
\end{proof}

\begin{remark}[$\delta_{\mathrm{CL}}$ from the prime distribution --- corrected]
\label{conj:cl-correction}
The continuum-limit correction, \emph{corrected in v2.29.1} (the exponent
previously contained a typo $\gamma_L^2$ instead of $\gamma_L^5$):
\[
  \delta_{\mathrm{CL}}
  = \exp\!\left(-\frac{\gamma_1}{4\pi\gamma_L^5}\right)
  = \exp\!\left(-\frac{\gamma_1^{(\mathrm{photon})}}{4\pi\gamma_L^4}\right)
  \approx \exp(-0.0500) \approx 0.9512,
\]
where $\gamma_1 \approx 14.134$ and $\gamma_1^{(\mathrm{photon})} =
\gamma_1/\gL \approx 7.583$ is the photon-frame first zero.
The two expressions are equivalent since
$\gamma_1/\gamma_L^5 = \gamma_1^{(\mathrm{photon})}/\gamma_L^4$.
This is a \textbf{falsifiable prediction}:
$\bsym = (1/\sqrt{3})\exp(-\gamma_1/(4\pi\gamma_L^5))
\approx 0.5492$, matching the measured value to $0.018\%$.

\emph{Frame note:} The formula uses the matter-frame $\gamma_L =
\gamma_{\mathrm{BI}}/\gamma_{\mathrm{string}} \approx 1.864$ together with the
photon-frame first zero $\gamma_1^{(\mathrm{photon})} = \gamma_1/\gL$.
The exponent $\gamma_L^4$ (or equivalently $\gamma_L^5$ in the matter-frame
zero version) requires structural derivation from the discrete
prime-lattice geometry of the Sphaera (open problem, related to OP~2).
\end{remark}


%% ============================================================
\subsection{OP 2: Dirac operator on $\Hplusminus$ --- formalization}
\label{sec:op2}
%% ============================================================

\subsubsection{Status and role}

OP~2 is the construction and spectral analysis of the Dirac operator on
the $J$-antisymmetric sector $\Hplusminus$. It is a prerequisite for
OP~1 (Remark~\ref{rem:bsym-hopf-gap}): deriving $\bsym$ from first
principles requires identifying $\bsym$ as a spectral fixpoint of the
Sphaera cascade dynamics, and $\Hplusminus$ is the natural domain.

\subsubsection{What is already established}

\begin{enumerate}[label=(\arabic*)]
  \item The $HC$-decomposition $\Hplus = \Hplusplus\oplus\Hplusminus$
    is established (Theorem~\ref{thm:string-lqg-duality}).
    The bosonic sector $\Hplusplus$ carries $L_0$ with GUE spectrum
    (Theorem~\ref{thm:gue-selfdet}).
  \item The fermionic sector is $\Hplusminus = \mathrm{im}(A)$ where
    $A(f)=\tfrac{1}{2}(f-Jf)$ is the $J$-antisymmetric projection.
    Pauli exclusion: $A(f)\wedge A(f)=0$.
  \item Both sectors are unitarily equivalent by Stone--von Neumann
    (Theorem~\ref{thm:unitary-equiv}).
  \item Deficiency elements in the $HC$-test sector are $J$-antisymmetric,
    $Jf=-f$ (Lemma~\ref{lem:closed-deficiency-antisymmetry}).
\end{enumerate}

\subsubsection{What is missing: the Dirac operator on $\Hplusminus$}

\begin{definition}[Candidate Dirac operator on $\Hplusminus$]
\label{def:dirac-op-Hplusminus}
Let $D_{-}$ denote a densely defined operator on $\Hplusminus$ satisfying:
\begin{enumerate}[label=(D\arabic*)]
  \item \textbf{$J$-anticommutation:} $D_{-}J=-JD_{-}$, i.e.\ $D_{-}$
    preserves the $J$-antisymmetric sector.
  \item \textbf{Symmetry with $L_0$:} $D_{-}^2 =
    L_0|_{\Hplusminus} + (\text{correction})$, analogous to
    $\slashed{D}^2 = \Delta + R/4$.
  \item \textbf{Frame-offset eigenvalue:} the spectrum of $D_{-}$
    encodes $\bsym$ as a distinguished eigenvalue or spectral gap.
  \item \textbf{Cascade compatibility:}
    $[D_{-},\Pi_n|_{\Hplusminus}]=0$ for all $n$.
\end{enumerate}
\end{definition}

\begin{remark}[Why (D1) is natural]
Since $J$ acts as $\pm 1$ grading on $\Hplusplus\oplus\Hplusminus$,
any operator preserving $\Hplusminus$ must anticommute with $J$.
This is the standard odd-operator condition in a real spectral triple.
\end{remark}

\begin{remark}[Candidate construction]
A natural candidate is $D_{-}:=-iJ L_0|_{\Hplusminus}$, using the
$J$-anticommutation to map $L_0$ from the bosonic to the fermionic
sector. This satisfies (D1) by construction. Whether it satisfies
(D2)--(D4) is the content of OP~2.
\end{remark}

\begin{proposition}[OP~2 closure targets]
\label{prop:op2-targets}
\emph{Status v2.29.3: (T1) and (T3) closed; (T2) and (T4) open.}
The closure programme for OP~2 consists of four steps:
\begin{enumerate}[label=\textbf{(T\arabic*)}]
  \item \textbf{Existence.} \checkmark\ \textbf{Closed in v2.29.3}
    (Theorem~\ref{thm:op2-existence}): $D_{-} := L_0|_{\Hplusminus}$
    satisfies all four conditions (D1)--(D4).
  \item \textbf{Essential self-adjointness.} \checkmark\ \textbf{Closed in
    v2.29.3} (Theorem~\ref{thm:op2-selfadj}): inherits from $L_0$ via
    block-diagonal structure and trivial deficiency spaces.
  \item \textbf{Spectral analysis.} \checkmark\ \textbf{Closed in
    v2.29.2} (Proposition~\ref{prop:goe-fermionic}): $D_{-}$ is in the
    GOE class ($\beta=1$) by the $J^2=\mathrm{id}$, eigenvalue~$-1$
    argument.
  \item \textbf{Frame-offset extraction.} \emph{Reframed in v2.29.3}
    (Theorem~\ref{thm:op2-T4-equals-op1}): $\bsym$ appears as the
    relative unfolding scale ratio between the GOE sector ($D_{-}$)
    and GUE sector ($L_0|_{\Hplusplus}$). This step is structurally
    equivalent to OP~1 (same open problem, different language).
\end{enumerate}
\end{proposition}

\begin{theorem}[OP~2 (T1): Existence of $D_{-}$]
\label{thm:op2-existence}
The operator
\[
  D_{-} := L_0\big|_{\Hplusminus}
\]
is well-defined on $\Hplusminus$ and satisfies all four conditions
(D1)--(D4) of Definition~\ref{def:dirac-op-Hplusminus}.
\end{theorem}

\begin{proof}
\textbf{Preservation of $\Hplusminus$.}
Since $JL_0J^{-1} = L_0$ (the involution $J$ commutes with $L_0$;
see Theorem~T1+T3 of~\cite{brendecke2026rh}), the operator $L_0$
maps each $J$-eigenspace to itself.
Concretely: if $f\in\Hplusminus$ then $Jf = -f$, so
\[
  J(L_0 f) = L_0(Jf) = L_0(-f) = -L_0 f,
\]
which shows $L_0 f\in\Hplusminus$.
Hence $D_{-} = L_0|_{\Hplusminus}$ maps $\Hplusminus$ to $\Hplusminus$.

\smallskip
\noindent\textbf{(D1) $J$-anticommutation.}
Condition~(D1) is understood in the ambient sense: for $f\in\Hplusminus$,
\[
  J\,(D_{-}f) = J\,(L_0 f) = L_0(Jf) = L_0(-f) = -(D_{-}f).
\]
Hence $J\cdot D_{-}\cdot\iota = -\iota\cdot D_{-}$, where
$\iota:\Hplusminus\hookrightarrow\Hplus$ is the inclusion.
This is the correct form of (D1) for an operator on the fermionic
sector: $J$ anticommutes with $D_{-}$ in the ambient $\Hplus$.

\smallskip
\noindent\textbf{(D2) Symmetry with $L_0$.}
$D_{-}^2 = L_0^2|_{\Hplusminus}$.
Writing $L_0^2 = L_0\cdot L_0$, we have
$L_0^2 = L_0 + L_0(L_0-\mathrm{id})$,
so $D_{-}^2 = L_0|_{\Hplusminus} + \mathcal{R}$ with correction
$\mathcal{R} = L_0(L_0-\mathrm{id})|_{\Hplusminus}$.
This is the Lichnerowicz-type relation required by~(D2).

\smallskip
\noindent\textbf{(D3) Frame-offset eigenvalue.}
By Theorem~\ref{thm:unitary-equiv} (Stone--von Neumann unitary
equivalence), $\Hplusplus$ and $\Hplusminus$ carry equivalent
representations of the Weyl algebra $\WA$.
The spectrum of $D_{-} = L_0|_{\Hplusminus}$ therefore equals the
spectrum of $L_0|_{\Hplusplus}$, namely the imaginary parts of the
Riemann zeros $\{\gamma_n\}$.
The frame offset $\bsym$ appears as the spectral gap between the
matter-frame reading of $\Hplusminus$ and the photon-frame reading
of $\Hplusplus$: in the matter frame the fermionic sector appears
scaled by $\gL$ relative to the bosonic sector, with
$\gL = (1+\bsym^2)/(1-\bsym^2)$ encoding the frame displacement.
Thus $\bsym$ is a distinguished spectral datum of $D_{-}$ in the
sense of~(D3).

\smallskip
\noindent\textbf{(D4) Cascade compatibility.}
$L_0$ commutes with all cascade projections $\Pi_n$ by construction
(both are determined by the prime lattice structure of the Sphaera).
Hence $[D_{-},\Pi_n|_{\Hplusminus}] = [L_0,\Pi_n]|_{\Hplusminus} = 0$.
\end{proof}

\begin{remark}[Corrected candidate]
\label{rem:candidate-correction}
The candidate $D_{-} = -iJ L_0|_{\Hplusminus}$ suggested in
the previous version (Remark in Section~\ref{sec:op2}) is incorrect:
$-iJL_0$ does not preserve $\Hplusminus$ because $J$ maps
$\Hplusminus$ to $\Hplusminus$ (with sign $-1$) while $-i$ introduces
a phase, giving $-iJL_0(A(f)) = -iJ(L_0 A(f)) = iL_0 A(f)\in\Hplusminus$.
Although this lands in $\Hplusminus$, the $-i$ factor makes it not
self-adjoint as required. The correct candidate $D_{-} = L_0|_{\Hplusminus}$
is real (no phase) and inherits self-adjointness from $L_0$.
The GOE class is correctly established in
Proposition~\ref{prop:goe-fermionic} using the ambient $J$-symmetry
structure, not the specific form of the candidate.
\end{remark}

\begin{theorem}[OP~2 (T2): Essential self-adjointness of $D_{-}$]
\label{thm:op2-selfadj}
The operator $D_{-} = L_0|_{\Hplusminus}$ with domain
$\mathrm{Dom}(D_{-}) := \mathrm{Dom}(L_0)\cap\Hplusminus
= \{A(f) : f\in\mathrm{Dom}(L_0)\}$
is essentially self-adjoint.
\end{theorem}

\begin{proof}
\textbf{Step 1: Block-diagonal structure.}
Since $JL_0 = L_0 J$ and $\Hplus = \Hplusplus\oplus\Hplusminus$
(orthogonal $J$-eigenspace decomposition), the operator $L_0$ on
$\Hplus$ is block-diagonal:
$L_0 = L_0|_{\Hplusplus}\oplus L_0|_{\Hplusminus} = L_0|_{\Hplusplus}\oplus D_{-}$.
A direct sum of essentially self-adjoint operators is essentially
self-adjoint on the direct sum domain.
Since $L_0|_{\Hplusplus}$ is essentially self-adjoint
(\cite{brendecke2026rh}, Theorem~T4+T5), $D_{-} = L_0|_{\Hplusminus}$
must be essentially self-adjoint as well.

\textbf{Step 2: Trivial deficiency spaces.}
Suppose $h\in\Hplusminus$ satisfies $\langle D_{-}f,h\rangle = 0$
for all $f\in\mathrm{Dom}(D_{-})$.
Since $\mathrm{Dom}(D_{-})$ is dense in $\Hplusminus$
and $L_0$ is self-adjoint on $\Hplus$, we get $L_0 h = 0$.
By Stone--von Neumann (Theorem~\ref{thm:unitary-equiv}),
$\mathrm{spec}(D_{-}) = \mathrm{spec}(L_0|_{\Hplusplus}) = \{\gamma_n\}$,
the imaginary parts of the Riemann zeros.
Since $0\notin\{\gamma_n\}$, we have $h=0$.
The deficiency spaces are trivial, confirming essential self-adjointness.
\end{proof}

\begin{theorem}[OP~2 (T4) = OP~1: $\bsym$ as spectral scale ratio]
\label{thm:op2-T4-equals-op1}
Condition~(T4) of Definition~\ref{def:dirac-op-Hplusminus} is
structurally equivalent to OP~1. Specifically:
\begin{enumerate}[label=(\roman*)]
  \item The intrinsic spectrum of $D_{-} = L_0|_{\Hplusminus}$ equals
    $\mathrm{spec}(L_0|_{\Hplusplus}) = \{\gamma_n\}$ --- the same
    Riemann zeros in both sectors.
  \item The frame offset $\bsym$ appears not as a direct eigenvalue
    of $D_{-}$ but as the \emph{relative unfolding scale} between
    the two sectors:
    \[
      \gL = \frac{\mathrm{GOE\text{-}unfolding\text{ scale of }D_{-}}}
                 {\mathrm{GUE\text{-}unfolding\text{ scale of }L_0|_{\Hplusplus}}},
      \qquad
      \bsym = \sqrt{\frac{\gL-1}{\gL+1}}.
    \]
  \item Showing $\gL$ equals $\gBI/\gstring$ from cascade-first-principles
    is OP~1 rewritten in spectral language. The two problems are identical.
\end{enumerate}
\end{theorem}

\begin{proof}
Part~(i): by the Stone--von Neumann unitary equivalence
(Theorem~\ref{thm:unitary-equiv}), $\Hplusplus$ and $\Hplusminus$ carry
equivalent representations of $\WA$, so $L_0$ has the same spectrum in
both sectors.

Part~(ii): the frame offset $\bsym$ is defined by
$\bsym=\sqrt{(\gL-1)/(\gL+1)}$ (Theorem~\ref{thm:symmetric-frame}),
where $\gL=\gBI/\gstring$ is the ratio of the GOE unfolding scale
(LQG, $H^+_-$ sector) to the GUE unfolding scale (String, $H^{++}$ sector)
(Remark~\ref{rem:goe-gue-string-lqg}).
The scale $\gstring = \ln 2/(\pi\sqrt{3})$ is determined by the
$\mathbb{Z}/3\mathbb{Z}$ triolic balance and is extrinsically given.
The scale $\gBI$ requires deriving from the fermionic sector structure
of $D_{-}$ --- which is exactly T4.

Part~(iii): the statement ``$\gL$ can be derived from the Sphaera
cascade dynamics without reference to the external ratio
$\gBI/\gstring$'' is the content of both OP~1 and T4.
They are the same problem.
\end{proof}

\begin{corollary}[OP~2 final status]
\label{cor:op2-status}
Of the four closure targets of OP~2:
\begin{center}
\begin{tabular}{cll}
\toprule
Step & Status & Reference \\
\midrule
(T1) Existence & \checkmark Closed (v2.29.3) & Theorem~\ref{thm:op2-existence} \\
(T2) Self-adjointness & \checkmark Closed (v2.29.3) & Theorem~\ref{thm:op2-selfadj} \\
(T3) GOE class & \checkmark Closed (v2.29.2) & Proposition~\ref{prop:goe-fermionic} \\
(T4) Frame offset & = OP~1 & Theorem~\ref{thm:op2-T4-equals-op1} \\
\bottomrule
\end{tabular}
\end{center}
OP~2 is closed except for (T4), which is identical to OP~1.
\textbf{OP~1 and OP~2 are the same open problem.}
\end{corollary}

\subsection{OP~1/OP~2 (T4): The Tribonacci fixed point}
\label{sec:tribonacci}

\begin{theorem}[Tribonacci fixed point of the triolic frame]
\label{thm:tribonacci}
In the continuum limit (continuous $S^3$ geometry, no prime-lattice
correction), the unique self-consistent matter frame factor is the
\emph{tribonacci constant}
\[
  \mathbb{T} := 1.8392867552\ldots,
\]
the real root of
\[
  \mu^3 - \mu^2 - \mu - 1 = 0.
\]
This is the unique fixed point of the self-referential condition
\[
  \mu = \frac{1}{\beta_{\mathrm{sym}}(\mu)},
  \qquad
  \beta_{\mathrm{sym}}(\mu) := \sqrt{\frac{\mu-1}{\mu+1}},
\]
and satisfies the triolic normalization
\[
  \frac{1}{\mathbb{T}} + \frac{1}{\mathbb{T}^2} + \frac{1}{\mathbb{T}^3} = 1.
\]
The measured frame factor $\gL = \gBI/\gstring \approx 1.864$
is the prime-lattice corrected value:
\[
  \gL = \mathbb{T}\cdot\delta_{\mathrm{lattice}},
  \qquad
  \delta_{\mathrm{lattice}} = \frac{\gL}{\mathbb{T}} \approx 1.0137,
\]
a $1.35\%$ correction from the discrete prime lattice $\Gamma_P$.
\end{theorem}

\begin{proof}
\textbf{Fixed-point derivation.}
Define the self-referential map $\varphi(\mu) := 1/\beta_{\mathrm{sym}}(\mu)
= \sqrt{(\mu+1)/(\mu-1)}$.
A fixed point satisfies $\mu = \varphi(\mu)$, i.e.,
$\mu^2 = (\mu+1)/(\mu-1)$, which gives $\mu^2(\mu-1) = \mu+1$, i.e.,
$\mu^3 - \mu^2 - \mu - 1 = 0$.
The real root $\mu > 1$ is the tribonacci constant $\mathbb{T}$.

\textbf{Triolic interpretation.}
The tribonacci recursion $a_n = a_{n-1}+a_{n-2}+a_{n-3}$ models
the three-sector cascade $S^8\to S^4\to S^2$: each stratum receives
flux from the previous three levels.
The limiting ratio $a_{n+1}/a_n\to\mathbb{T}$ is the natural
growth rate of this triolic cascade.
Dividing $\mathbb{T}^3 = \mathbb{T}^2+\mathbb{T}+1$ by $\mathbb{T}^3$
gives the normalization
$1/\mathbb{T}+1/\mathbb{T}^2+1/\mathbb{T}^3 = 1$,
which is the cascade probability sum:
\begin{align*}
  P(S^4\to S^2) &= 1/\mathbb{T} \approx 0.5437, \\
  P(S^8\to S^4) &= 1/\mathbb{T}^2 \approx 0.2956, \\
  P(S^8\to S^2) &= 1/\mathbb{T}^3 \approx 0.1607.
\end{align*}

\textbf{Self-referential meaning.}
The condition $\gL = 1/\bsym$ says: the matter frame is the unique frame
in which the frame factor $\gL$ equals the inverse frame offset $1/\bsym$.
Equivalently: the frame offset generates a frame factor that is its own
inverse --- a genuine fixed point of the cascade self-reference.

\textbf{Prime-lattice correction.}
The tribonacci root is the \emph{continuum-limit} value.
In the corrected post-v3.31.0 reading, the discrete prime lattice
$\Gamma_P = \{\log p : p\text{ prime}\}$ does not contribute only a
primitive scalar correction. Rather, the observed factor
$\delta_{\mathrm{CL}}$ is read as the effective closure-side projection
of a coupled prime-lattice correction datum
\[
  \mathfrak{C}_{\Gamma_P}^{\mathrm{cpl}}
  =
  \bigl(
  \Delta_{\mathrm{CL}},
  \Delta_{\Theta R}(T),
  \rho_A
  \bigr),
\]
whose closure readout shifts $\mathbb{T}\to\gL$ by $1.35\%$.
Thus the same coupled lattice source underlies the Hopf gap and the
$\Omega_M$ gap, while the scalar factor $\delta_{\mathrm{CL}}$ remains
its observable closure-side fit.

\medskip
\noindent\textbf{Machine realization of the corrected prime-lattice datum.}
On the post-v3.31.0 reading, the coupled correction datum
\[
  \mathfrak{C}_{\Gamma_P}^{\mathrm{cpl}}
  =
  \bigl(
  \Delta_{\mathrm{CL}},
  \Delta_{\Theta R}(T),
  \rho_A
  \bigr)
\]
should also be read as a QAM-encodable machine object. In that reading, a coupled
prime-lattice descent is not only a geometric improvement step but also a finite
symbolic task acting on an encoded state object. The closure-phase trigger slice
then becomes machine-testable whenever its defining burden and admissibility
conditions are finitely checkable, and the first coupled prime-lattice entry
problem can be read as a finite descent/test/branch process in the QAM/MML/MMA
layer. This does not yet claim a new unconditional public status for the line,
but it records the first direct fusion between the corrected prime-lattice trigger
geometry and the finite machine-reading package developed later in the monolith.
\end{proof}

\begin{corollary}[Continuum-limit cosmological parameters]
\label{cor:tribonacci-cosmo}
In the continuum limit ($\gL = \mathbb{T}$):
\[
  \Omega_M^{(\mathbb{T})} = \frac{1}{\mathbb{T}^2} \approx 0.2956,
  \qquad
  \Omega_\Lambda^{(\mathbb{T})} = 1-\frac{1}{\mathbb{T}^2} \approx 0.7044.
\]
The measured values $\Omega_M = 0.3111$, $\Omega_\Lambda = 0.6889$
(Planck 2018) differ by $\approx 2.0\%$, consistent with the
$1.35\%$ prime-lattice correction applied to both sectors.
\end{corollary}

\begin{remark}[OP~1/OP~2 (T4) resolution]
\label{rem:op1-tribonacci}
Theorem~\ref{thm:tribonacci} resolves OP~1 and OP~2~(T4)
\emph{in the continuum limit}:
\begin{enumerate}[label=(\roman*)]
  \item The frame factor $\gL = \mathbb{T}$ (tribonacci) is
    the unique solution of the self-referential condition
    $\gL = 1/\bsym(\gL)$ imposed by the $\mathbb{Z}/3\mathbb{Z}$
    triolic structure.
  \item The $1.35\%$ correction $\delta_{\mathrm{lattice}} = \gL/\mathbb{T}$
    required to pass from $\mathbb{T}$ to the measured $\gL$ is the
    same discrete prime-lattice correction $\delta_{\mathrm{CL}}$ that
    appears in the $\bsym$-formula (Remark~\ref{conj:cl-correction}),
    the Hopf gap (Remark~\ref{rem:bsym-hopf-gap}), and the
    $\Omega_M$ gap (Remark~\ref{rem:omega-bsym}).
  \item The full derivation of $\delta_{\mathrm{lattice}}$ from the
    prime lattice $\Gamma_P$ remains the one surviving open step
    (previously OP~1; previously OP~2~(T4)).
\end{enumerate}
\end{remark}

\begin{theorem}[First machine-fused coupled OP-entry theorem]
\label{thm:first-machine-fused-coupled-op-entry}
Assume:
\begin{enumerate}
    \item the corrected prime-lattice reading applies on a transfer-consistent closure-phase window;
    \item the first proxy-to-residual transfer theorem applies, so that the active inherited burden is read through
    \[
      \bigl(\Delta_{\mathrm{CL}},\Delta_{\Theta R}(T)\bigr)
    \]
    with $\rho_A$ as admissibility monitor;
    \item the coupled prime-lattice correction datum
    \[
      \mathfrak{C}_{\Gamma_P}^{\mathrm{cpl}}
      =
      \bigl(\Delta_{\mathrm{CL}},\Delta_{\Theta R}(T),\rho_A\bigr)
    \]
    is QAM-encodable and carried by a machine-reachable state cell;
    \item coupled prime-lattice descent is realizable in the QAM/MML/MMA layer as a finite symbolic descent/test/branch task;
    \item the closure-phase trigger slice is machine-testable by finitely checkable burden and admissibility conditions.
\end{enumerate}
Then the first post-v3.31.0 OP-entry problem admits a machine-fused coupled realization.

More precisely:
\begin{enumerate}
    \item the operative entry object is the coupled datum
    \[
      \bigl(\Delta_{\mathrm{CL}},\Delta_{\Theta R}(T),\rho_A\bigr),
    \]
    not the scalar shadow $\delta_{\mathrm{CL}}$ alone;
    \item every admissible local prime-lattice descent step can be read simultaneously
    \begin{enumerate}
        \item geometrically, as a coupled closure-phase descent step, and
        \item computationally, as a finite realizable machine step on the encoded state object;
    \end{enumerate}
    \item the closure-phase trigger slice is branch-decidable in the machine layer;
    \item therefore, whenever a finite admissible descent chain reaches the closure-phase trigger slice while preserving admissibility of $\rho_A$, immediate OP-entry follows.
\end{enumerate}
\end{theorem}

\begin{proof}
By the corrected prime-lattice reading, the scalar continuum-limit correction is only the closure-side shadow of the coupled datum
\[
  \mathfrak{C}_{\Gamma_P}^{\mathrm{cpl}}
  =
  \bigl(\Delta_{\mathrm{CL}},\Delta_{\Theta R}(T),\rho_A\bigr).
\]
By the proxy-to-residual transfer theorem, the active inherited burden is reduced to the pair
\[
  \bigl(\Delta_{\mathrm{CL}},\Delta_{\Theta R}(T)\bigr)
\]
under admissible $\rho_A$.
By QAM-encodability, this coupled datum is a machine-reachable state object.
By finite realizability of coupled prime-lattice descent, each admissible local descent is executable as a finite symbolic task in the QAM/MML/MMA layer.
By machine-testability of the closure-phase trigger slice, the machine can decide whether the updated coupled datum has reached entry, must continue descent, or has become inadmissible.
Hence the first post-v3.31.0 OP-entry problem is machine-fused: it is simultaneously a coupled geometric descent problem and a finite machine-controlled descent/test/branch problem.
Whenever the finite admissible descent chain reaches the closure-phase trigger slice while preserving $\rho_A$ admissibility, immediate OP-entry follows by the prime-lattice coupled entry lemma.
\end{proof}

\begin{corollary}[First post-v3.31.0 freeze threshold]
\label{cor:first-post-v3310-freeze-threshold}
The first machine-fused coupled OP-entry theorem marks the first natural mathematical freeze threshold for a post-v3.31.0 minor release: the line now carries a single fused claim linking the corrected prime-lattice entry geometry with the QAM/MML/MMA machine layer.
\end{corollary}

\begin{proof}
Immediate from Theorem~\ref{thm:first-machine-fused-coupled-op-entry}.
\end{proof}


\begin{definition}[Machine-fused local entry state]
\label{def:machine-fused-local-entry-state}
A \emph{machine-fused local entry state} is a tuple
\[
\mathfrak{E}_{\mathrm{loc}}
=
\bigl(
\mathbf{X}_{\Gamma_P},
\mathcal{C}_{\mathrm{mach}},
\mathfrak{D}_{\mathrm{lead}},
\mathfrak{S}_{\mathrm{trig}}
\bigr),
\]
where:
\begin{enumerate}
    \item $\mathbf{X}_{\Gamma_P}$ is the current prime-lattice machine object
    \[
    \mathbf{X}_{\Gamma_P}
    =
    \operatorname{Enc}_{\mathrm{QAM}}
    \bigl(
    \Delta_{\mathrm{CL}},
    \Delta_{\Theta R}(T),
    \rho_A
    \bigr);
    \]
    \item $\mathcal{C}_{\mathrm{mach}}$ is the current finite machine configuration carrying $\mathbf{X}_{\Gamma_P}$;
    \item $\mathfrak{D}_{\mathrm{lead}}$ is the current lead-direction tag;
    \item $\mathfrak{S}_{\mathrm{trig}}$ is the currently active trigger slice candidate.
\end{enumerate}
\end{definition}

\begin{definition}[Local machine-fused descent/test/branch operator]
\label{def:local-machine-fused-descent-test-branch-operator}
The \emph{local machine-fused descent/test/branch operator} is the composite step
\[
\mathcal{O}_{\mathrm{DTB}}
=
\mathcal{B}_{\Gamma_P}
\circ
\mathcal{T}_{\Gamma_P}
\circ
\mathcal{D}_{\Gamma_P},
\]
where:
\begin{enumerate}
    \item $\mathcal{D}_{\Gamma_P}$ is a coupled prime-lattice descent realization;
    \item $\mathcal{T}_{\Gamma_P}$ is the prime-lattice trigger test;
    \item $\mathcal{B}_{\Gamma_P}$ is the prime-lattice branch rule.
\end{enumerate}
\end{definition}

\begin{remark}[Operational reading]
\label{rem:operational-reading-local-operator}
This is the first explicit post-v3.32.0 local controller:
first lower the coupled burden pair if possible,
then test the trigger slice,
then branch into entry, continuation, or obstruction.
\end{remark}

\begin{definition}[Local obstruction certificate]
\label{def:local-obstruction-certificate}
A \emph{local obstruction certificate} is a finite machine-readable object
\[
\mathfrak{O}_{\mathrm{loc}}
\]
issued when the local machine-fused descent/test/branch operator returns
\[
\mathtt{inadmissible}
\quad\text{or}\quad
\mathtt{stalled},
\]
together with a finite explanation of which of the following failed:
\begin{enumerate}
    \item closure descent,
    \item Theta descent,
    \item admissibility of $\rho_A$,
    \item persistence of the closure-phase trigger slice,
    \item route-faithful machine realizability.
\end{enumerate}
\end{definition}

\begin{definition}[Local continuation certificate]
\label{def:local-continuation-certificate}
A \emph{local continuation certificate} is a finite machine-readable object
\[
\mathfrak{K}_{\mathrm{cont}}
\]
issued when the local operator returns
\[
\mathtt{continue},
\]
certifying that:
\begin{enumerate}
    \item the coupled descent remains admissible;
    \item the current state is still outside immediate OP-entry;
    \item at least one monotone local burden measure has strictly improved.
\end{enumerate}
\end{definition}

\begin{definition}[Local entry certificate]
\label{def:local-entry-certificate}
A \emph{local entry certificate} is a finite machine-readable object
\[
\mathfrak{K}_{\mathrm{entry}}
\]
issued when the local operator returns
\[
\mathtt{entry},
\]
certifying that the current machine-fused prime-lattice datum has entered the active closure-phase trigger slice and is immediately OP-entry certified.
\end{definition}

\begin{definition}[Machine-fused local entry protocol]
\label{def:machine-fused-local-entry-protocol}
The \emph{machine-fused local entry protocol} is the iterative rule:
\[
\mathfrak{E}_{\mathrm{loc},0}
\longmapsto
\mathfrak{E}_{\mathrm{loc},1}
\longmapsto
\mathfrak{E}_{\mathrm{loc},2}
\longmapsto \cdots
\]
obtained by repeated application of the local descent/test/branch operator, stopping when either
\[
\mathfrak{K}_{\mathrm{entry}}
\quad\text{or}\quad
\mathfrak{O}_{\mathrm{loc}}
\]
is issued.
\end{definition}

\begin{remark}[Why this is the right next step]
\label{rem:why-right-next-step-local-protocol}
The post-v3.32.0 line should no longer remain at the level of existence claims alone.
It should now either produce:
\begin{enumerate}
    \item a local entry certificate, or
    \item a local obstruction certificate.
\end{enumerate}
That is the first genuinely operational OP-entry phase.
\end{remark}

\begin{definition}[Local burden descent norm]
\label{def:local-burden-descent-norm}
A \emph{local burden descent norm} is any monotone function
\[
\mathcal{N}_{\mathrm{loc}}
\bigl(
\Delta_{\mathrm{CL}},
\Delta_{\Theta R}(T),
\rho_A
\bigr)
\]
such that:
\begin{enumerate}
    \item it decreases when both transferred burdens decrease and admissibility is preserved;
    \item it is defined only on admissible machine-fused local entry states;
    \item lower values indicate closer approach to immediate OP-entry.
\end{enumerate}
\end{definition}

\begin{proposition}[Continuation implies local progress]
\label{prop:continuation-implies-local-progress}
Assume the machine-fused local entry protocol returns a continuation certificate
\[
\mathfrak{K}_{\mathrm{cont}}
\]
at some step.
Then there exists a local burden descent norm
\[
\mathcal{N}_{\mathrm{loc}}
\]
for which the updated local entry state has strictly smaller value than the previous one.
\end{proposition}

\begin{proof}
By definition of the continuation certificate, coupled descent remains admissible, immediate OP-entry has not yet been reached, and at least one monotone local burden measure has strictly improved.
Such a measure may be taken as a local burden descent norm, yielding strict local progress.
\end{proof}

\begin{lemma}[Finite local entry-or-obstruction lemma]
\label{lem:finite-local-entry-or-obstruction-lemma}
Assume:
\begin{enumerate}
    \item the machine-fused local entry protocol is defined on the current local state;
    \item every continuation step strictly lowers a local burden descent norm;
    \item the protocol cannot descend indefinitely inside the same admissible slice without either entering the trigger slice or violating admissibility.
\end{enumerate}
Then a finite local protocol run ends with either:
\begin{enumerate}
    \item a local entry certificate, or
    \item a local obstruction certificate.
\end{enumerate}
\end{lemma}

\begin{proof}
Each continuation step strictly lowers the local burden descent norm.
By assumption, indefinite descent in the same admissible slice is impossible.
Hence the protocol must terminate after finitely many steps, either by entering the trigger slice or by losing admissibility/persistence and issuing an obstruction certificate.
\end{proof}

\begin{theorem}[First machine-fused local OP-entry protocol theorem]
\label{thm:first-machine-fused-local-op-entry-protocol-theorem}
Assume:
\begin{enumerate}
    \item the corrected prime-lattice reading applies;
    \item the first machine-fused coupled OP-entry theorem applies;
    \item the closure-phase trigger slice is machine-testable;
    \item the current line admits a machine-fused local entry state.
\end{enumerate}
Then the post-v3.32.0 line supports a finite local machine-fused OP-entry protocol.

More precisely, every admissible run of the local protocol yields after finitely many steps either:
\begin{enumerate}
    \item a local entry certificate witnessing immediate OP-entry, or
    \item a local obstruction certificate isolating the first local failure mode.
\end{enumerate}
\end{theorem}

\begin{proof}
By the first machine-fused coupled OP-entry theorem, coupled descent, trigger testing, and branch handling are already fused in the current line. Hence the local descent/test/branch operator is defined on the current machine-fused local entry state.
By the finite local entry-or-obstruction lemma, every admissible local run terminates after finitely many steps with either entry or obstruction.
\end{proof}

\begin{corollary}[Post-v3.32.0 working rule]
\label{cor:post-v3320-working-rule}
The next working line after v3.32.0 should no longer ask only whether the coupled OP-entry theorem exists abstractly.
It should explicitly run the local machine-fused entry protocol and record either:
\begin{enumerate}
    \item the first local entry certificate, or
    \item the first local obstruction certificate.
\end{enumerate}
\end{corollary}

\begin{proof}
Immediate from Theorem~\ref{thm:first-machine-fused-local-op-entry-protocol-theorem}.
\end{proof}



\begin{definition}[Default first local outcome hypothesis]
\label{def:default-first-local-outcome-hypothesis}
The \emph{default first local outcome hypothesis} for the post-v3.32.0 line is the hypothesis that the first honest terminal output of the machine-fused local entry protocol is a local obstruction certificate rather than a local entry certificate.
\end{definition}

\begin{remark}[Why obstruction is the current default]
\label{rem:why-obstruction-is-the-current-default}
At the present stage, the coupled descent mechanism is already structurally available, but the active closure-phase trigger slice is still only symbolically specified and not yet fully calibrated by explicit local threshold data. Therefore the most likely first terminal failure is not loss of the whole machine-fused architecture, but failure to certify persistence or actual reach of the trigger slice before immediate entry is declared.
\end{remark}

\begin{definition}[Trigger-persistence obstruction]
\label{def:trigger-persistence-obstruction}
A \emph{trigger-persistence obstruction} is a local obstruction certificate
\[
\mathfrak{O}_{\mathrm{loc}}^{\mathrm{trig}}
\]
issued when:
\begin{enumerate}
    \item coupled descent remains machine-realizable;
    \item the packet-side monitor remains admissible;
    \item but persistence or actual reach of the closure-phase trigger slice cannot yet be certified on the current local window.
\end{enumerate}
\end{definition}

\begin{definition}[Entry-prior versus obstruction-prior local regime]
\label{def:entry-prior-versus-obstruction-prior-local-regime}
A machine-fused local entry state is called:
\begin{enumerate}
    \item \emph{entry-prior} if the active closure-phase trigger slice is sufficiently explicit and stable that the next finite admissible local run is expected to end in a local entry certificate;
    \item \emph{obstruction-prior} if coupled descent is available but the trigger slice remains insufficiently explicit, insufficiently persistent, or insufficiently calibrated for immediate entry certification.
\end{enumerate}
\end{definition}

\begin{proposition}[The current local state is more naturally obstruction-prior than entry-prior]
\label{prop:current-local-state-more-naturally-obstruction-prior}
Assume the post-v3.32.0 line already supports:
\begin{enumerate}
    \item corrected prime-lattice reading;
    \item proxy-to-residual transfer;
    \item machine-fused descent/test/branch control.
\end{enumerate}
Assume further that the active closure-phase trigger slice is still only symbolically described and not yet equipped with fully explicit local threshold calibration.
Then the current machine-fused local entry state is more naturally obstruction-prior than entry-prior.
\end{proposition}

\begin{proof}
Under the stated assumptions, coupled descent and testing already exist as local machine operations, so the current line is not blocked at the level of basic realizability.
However, if the closure-phase trigger slice is still only symbolically described and lacks fully explicit local threshold calibration, then immediate entry cannot yet be certified purely from the present data.
Hence the more natural first terminal local outcome is an obstruction certificate isolating trigger-persistence or trigger-calibration failure rather than an entry certificate.
\end{proof}

\begin{theorem}[First local obstruction-priority theorem]
\label{thm:first-local-obstruction-priority-theorem}
Assume:
\begin{enumerate}
    \item the first machine-fused local OP-entry protocol theorem applies;
    \item the current line admits machine-fused coupled descent steps;
    \item the packet-side monitor remains admissible on the first relevant local window;
    \item the closure-phase trigger slice is not yet fully explicit or fully calibrated on that same window.
\end{enumerate}
Then the current local machine-fused state is obstruction-prior.
More precisely, the first honest terminal local outcome is expected to be a trigger-persistence obstruction
\[
\mathfrak{O}_{\mathrm{loc}}^{\mathrm{trig}},
\]
rather than an immediate local entry certificate.
\end{theorem}

\begin{proof}
By assumption, the local machine-fused protocol exists and coupled descent is available.
By admissibility of the packet-side monitor, the protocol is not forced to fail immediately by packet instability.
If, however, the closure-phase trigger slice is not yet fully explicit or calibrated, then the protocol cannot honestly certify immediate entry on the present local window.
The most natural terminal outcome is therefore a local obstruction certificate isolating the missing trigger-persistence or trigger-calibration ingredient.
\end{proof}

\begin{corollary}[First post-v3.32.0 local audit rule]
\label{cor:first-post-v3320-local-audit-rule}
The next local audit after v3.32.0 should not first ask whether immediate entry already occurs.
It should first ask which trigger-persistence or trigger-calibration ingredient is missing from the current local window, i.e. it should aim first at the first trigger-persistence obstruction certificate.
\end{corollary}

\begin{proof}
Immediate from Theorem~\ref{thm:first-local-obstruction-priority-theorem}.
\end{proof}


\begin{definition}[Trigger-calibration obstruction]
\label{def:trigger-calibration-obstruction}
A \emph{trigger-calibration obstruction} is a local obstruction certificate
\[
\mathfrak{O}_{\mathrm{loc}}^{\mathrm{cal}}
\]
issued when:
\begin{enumerate}
    \item the coupled descent/test/branch architecture is available;
    \item the active closure-phase trigger slice is recognized symbolically;
    \item but no sufficiently explicit local burden thresholds or effective local trigger inequalities are yet available to certify entry.
\end{enumerate}
\end{definition}

\begin{definition}[Trigger-persistence obstruction --- strengthened reading]
\label{def:trigger-persistence-obstruction-strengthened-reading}
A \emph{strengthened trigger-persistence obstruction} is a local obstruction certificate
\[
\mathfrak{O}_{\mathrm{loc}}^{\mathrm{pers}}
\]
issued when:
\begin{enumerate}
    \item the coupled descent/test/branch architecture is available;
    \item a symbolic closure-phase trigger slice is available;
    \item some local trigger-improving descent steps are visible;
    \item but those steps do not yet certify persistence of the trigger slice on a sufficiently long local window.
\end{enumerate}
\end{definition}

\begin{definition}[Machine-side decodability obstruction]
\label{def:machine-side-decodability-obstruction}
A \emph{machine-side decodability obstruction} is a local obstruction certificate
\[
\mathfrak{O}_{\mathrm{loc}}^{\mathrm{dec}}
\]
issued when:
\begin{enumerate}
    \item the corrected prime-lattice datum is present;
    \item but the corresponding trigger-relevant condition cannot yet be realized as a finite machine-decidable test on the active QAM/MML/MMA state.
\end{enumerate}
\end{definition}

\begin{remark}[The three natural first obstruction types]
\label{rem:the-three-natural-first-obstruction-types}
At the present stage, the first local obstruction should naturally be expected to belong to exactly one of three classes:
\begin{enumerate}
    \item trigger-calibration obstruction;
    \item trigger-persistence obstruction;
    \item machine-side decodability obstruction.
\end{enumerate}
These correspond respectively to missing local thresholds, missing stable trigger persistence, or missing test-level realizability.
\end{remark}

\begin{proposition}[Machine-side decodability is not the primary current obstruction]
\label{prop:machine-side-decodability-not-primary-current-obstruction}
Assume:
\begin{enumerate}
    \item the corrected prime-lattice reading applies;
    \item the prime-lattice machine realization theorem applies;
    \item the machine-fused local entry protocol theorem applies.
\end{enumerate}
Then a machine-side decodability obstruction is not the most natural first local obstruction of the current line.
\end{proposition}

\begin{proof}
Under the stated assumptions, the coupled prime-lattice datum is already QAM-encodable, coupled descent is MML/MMA-realizable, and the trigger problem is already machine-fused at the protocol level.
Hence the current line is not most naturally blocked by complete absence of machine-side realizability.
The more natural remaining failure modes lie at the level of trigger calibration or trigger persistence.
\end{proof}

\begin{proposition}[Current obstruction split is calibration/persistence-dominant]
\label{prop:current-obstruction-split-is-calibration-persistence-dominant}
Assume the hypotheses of Proposition~\ref{prop:machine-side-decodability-not-primary-current-obstruction}.
Assume further that the active closure-phase trigger slice is already symbolically available but not yet fully quantified by explicit local thresholds and not yet stable on a sufficiently long local window.
Then the first local obstruction is more naturally of calibration/persistence type than of decodability type.
\end{proposition}

\begin{proof}
If the trigger slice is already symbolically available, then total trigger absence is not the main issue.
If it is not yet fully quantified by explicit thresholds and not yet stable on a sufficiently long local window, then the missing ingredient lies in calibration and persistence.
By Proposition~\ref{prop:machine-side-decodability-not-primary-current-obstruction}, this is more natural than a primary machine-side decodability failure.
\end{proof}

\begin{definition}[Current dominant first local obstruction class]
\label{def:current-dominant-first-local-obstruction-class}
The \emph{current dominant first local obstruction class} is the class
\[
\mathfrak{O}_{\mathrm{loc}}^{\ast}
\in
\{
\mathfrak{O}_{\mathrm{loc}}^{\mathrm{cal}},
\mathfrak{O}_{\mathrm{loc}}^{\mathrm{pers}},
\mathfrak{O}_{\mathrm{loc}}^{\mathrm{dec}}
\}
\]
which is most naturally expected to occur first on the current local window.
\end{definition}

\begin{theorem}[First local obstruction classification theorem]
\label{thm:first-local-obstruction-classification-theorem}
Assume:
\begin{enumerate}
    \item the first local obstruction-priority theorem applies;
    \item the corrected prime-lattice reading applies;
    \item the prime-lattice machine realization theorem applies;
    \item the closure-phase trigger slice is already symbolically available;
    \item the trigger slice is not yet fully locally calibrated and not yet certified persistent on a sufficiently long local window.
\end{enumerate}
Then the current dominant first local obstruction class is not machine-side decodability but calibration/persistence type.

More precisely,
\[
\mathfrak{O}_{\mathrm{loc}}^{\ast}
\in
\{
\mathfrak{O}_{\mathrm{loc}}^{\mathrm{cal}},
\mathfrak{O}_{\mathrm{loc}}^{\mathrm{pers}}
\},
\]
and the current line should be read as calibration/persistence-obstruction-prior.
\end{theorem}

\begin{proof}
By the prime-lattice machine realization theorem, the coupled datum is already machine-encodable and machine-fused descent/test/branch control is available, so pure decodability failure is not the most natural first obstruction.
Because the closure-phase trigger slice is already symbolically available but not yet fully calibrated or persistent on a sufficiently long local window, the remaining failure modes lie in trigger calibration and trigger persistence.
Hence the current dominant obstruction class is of calibration/persistence type.
\end{proof}

\begin{corollary}[Next post-v3.32.0 local audit question]
\label{cor:next-post-v3320-local-audit-question}
The next local audit after the current r13 stage should first ask:
\begin{enumerate}
    \item what explicit local trigger inequalities are still missing;
    \item and what minimum persistence window is still missing
\end{enumerate}
before treating machine-side decodability as the primary obstruction.
\end{corollary}

\begin{proof}
Immediate from Theorem~\ref{thm:first-local-obstruction-classification-theorem}.
\end{proof}

\begin{definition}[Local closure-phase trigger inequality schema]
\label{def:local-closure-phase-trigger-inequality-schema}
A \emph{local closure-phase trigger inequality schema} is a finite family of inequalities of the form
\[
\Delta_{\mathrm{CL}} \le \varepsilon_{\mathrm{CL}},
\qquad
\Delta_{\Theta R}(T) \le \varepsilon_{\Theta},
\qquad
\rho_A \in \mathfrak{A}_{\mathrm{adm}}^{\mathrm{safe}},
\]
together with the side condition that the current transition remains closure-phase led on the same local window.
Here
\[
\varepsilon_{\mathrm{CL}}>0,
\qquad
\varepsilon_{\Theta}>0
\]
are local trigger tolerances and
\[
\mathfrak{A}_{\mathrm{adm}}^{\mathrm{safe}}\subseteq \mathfrak{A}_{\mathrm{adm}}
\]
is a distinguished safe admissibility region for the packet-side monitor.
\end{definition}

\begin{remark}[Meaning of the inequality schema]
\label{rem:meaning-of-the-inequality-schema}
The current line does not yet require a final closed-form trigger formula.
At the present stage, what is missing is already captured by the weaker data of explicit local tolerances for the coupled burdens together with a safe admissibility band for \(\rho_A\).
\end{remark}

\begin{definition}[Local closure-phase persistence window of length \texorpdfstring{$m$}{m}]
\label{def:local-closure-phase-persistence-window-of-length-m}
Let \(m\ge 1\).
A \emph{local closure-phase persistence window of length \(m\)} is a block of successive local machine-fused protocol steps
\[
\mathfrak{E}_{\mathrm{loc},n},\mathfrak{E}_{\mathrm{loc},n+1},\dots,\mathfrak{E}_{\mathrm{loc},n+m-1}
\]
such that throughout the block:
\begin{enumerate}
    \item the lead direction remains closure-phase led;
    \item the trigger inequality schema is satisfied monotonically or remains stable once entered;
    \item the packet-side monitor stays in \(\mathfrak{A}_{\mathrm{adm}}^{\mathrm{safe}}\);
    \item the same route-faithful local slice remains active.
\end{enumerate}
\end{definition}

\begin{definition}[Minimum local trigger-persistence length]
\label{def:minimum-local-trigger-persistence-length}
A positive integer
\[
m_{\min}\ge 1
\]
is called a \emph{minimum local trigger-persistence length} if immediate OP-entry is not certified merely by entering the local trigger inequality schema once, but only after the schema persists on a closure-phase persistence window of length at least \(m_{\min}\).
\end{definition}

\begin{remark}[Why a persistence length is needed]
\label{rem:why-a-persistence-length-is-needed}
The current obstruction analysis suggests that a one-step near-entry fluctuation is not yet trustworthy.
What is missing is not just smallness of the coupled burdens, but stable smallness on a sufficiently long local window.
\end{remark}

\begin{definition}[Trigger-calibration data package]
\label{def:trigger-calibration-data-package}
The \emph{trigger-calibration data package} is the tuple
\[
\mathfrak{P}_{\mathrm{trig}}^{\mathrm{loc}}
:=
\bigl(
\varepsilon_{\mathrm{CL}},
\varepsilon_{\Theta},
\mathfrak{A}_{\mathrm{adm}}^{\mathrm{safe}},
 m_{\min}
\bigr)
\]
consisting of local burden tolerances, a safe admissibility region, and a minimum persistence length.
\end{definition}

\begin{proposition}[What the current obstruction is really asking for]
\label{prop:what-the-current-obstruction-is-really-asking-for}
Assume the first local obstruction classification theorem.
Then the current calibration/persistence obstruction is equivalent to the absence of a sufficiently explicit trigger-calibration data package.
\end{proposition}

\begin{proof}
By the obstruction classification theorem, the current obstruction is not primarily one of machine-side decodability, but of calibration/persistence type.
Calibration requires explicit local burden tolerances and a safe admissibility region, while persistence requires a minimum stable window length.
Therefore the current obstruction is exactly the absence of a sufficiently explicit trigger-calibration data package.
\end{proof}

\begin{definition}[Locally calibrated closure-phase trigger slice]
\label{def:locally-calibrated-closure-phase-trigger-slice}
A closure-phase trigger slice is called \emph{locally calibrated} if it is equipped with a trigger-calibration data package
\[
\mathfrak{P}_{\mathrm{trig}}^{\mathrm{loc}}
\]
and immediate OP-entry is certified whenever the associated local closure-phase trigger inequality schema persists for at least \(m_{\min}\) consecutive local protocol steps.
\end{definition}

\begin{lemma}[First local calibration-to-entry lemma]
\label{lem:first-local-calibration-to-entry-lemma}
Assume:
\begin{enumerate}
    \item the machine-fused local entry protocol theorem applies;
    \item a locally calibrated closure-phase trigger slice is available;
    \item the current local protocol run enters the associated trigger inequality schema and remains there for at least \(m_{\min}\) consecutive steps.
\end{enumerate}
Then the local protocol yields a local entry certificate.
\end{lemma}

\begin{proof}
By local calibration, entry is certified exactly when the trigger inequality schema persists on a closure-phase window of length at least \(m_{\min}\).
If the current local protocol run satisfies this condition, then by definition the corresponding state is immediately OP-entry certified and the protocol returns a local entry certificate.
\end{proof}

\begin{theorem}[First local trigger calibration theorem]
\label{thm:first-local-trigger-calibration-theorem}
Assume:
\begin{enumerate}
    \item the first local obstruction classification theorem applies;
    \item the machine-fused local entry protocol theorem applies;
    \item the active obstruction is of calibration/persistence type.
\end{enumerate}
Then the next missing local ingredient is exactly the explicit trigger-calibration data package
\[
\mathfrak{P}_{\mathrm{trig}}^{\mathrm{loc}}
=
\bigl(
\varepsilon_{\mathrm{CL}},
\varepsilon_{\Theta},
\mathfrak{A}_{\mathrm{adm}}^{\mathrm{safe}},
 m_{\min}
\bigr).
\]
Moreover, once such a package is fixed, the remaining local question reduces to whether the current protocol run can enter and maintain the associated inequality schema for at least \(m_{\min}\) consecutive steps.
\end{theorem}

\begin{proof}
By the first local obstruction classification theorem, the dominant current obstruction is not decodability but calibration/persistence type.
By Proposition~\ref{prop:what-the-current-obstruction-is-really-asking-for}, this is exactly the absence of explicit local burden tolerances, a safe admissibility region, and a minimum persistence length.
These data are precisely the trigger-calibration data package.
Once the package is fixed, the only remaining local issue is whether the current machine-fused protocol run satisfies the corresponding persistence requirement.
\end{proof}

\begin{corollary}[Next post-r14 working rule]
\label{cor:next-post-r14-working-rule}
The next working step after the current obstruction analysis should not try to prove immediate OP-entry in one jump.
It should first determine:
\begin{enumerate}
    \item plausible local values or bounds for \(\varepsilon_{\mathrm{CL}}\) and \(\varepsilon_{\Theta}\);
    \item a safe admissibility subregion \(\mathfrak{A}_{\mathrm{adm}}^{\mathrm{safe}}\);
    \item a first plausible minimum persistence length \(m_{\min}\).
\end{enumerate}
\end{corollary}

\begin{proof}
Immediate from Theorem~\ref{thm:first-local-trigger-calibration-theorem}.
\end{proof}



\begin{definition}[Default local trigger orientation]
\label{def:default-local-trigger-orientation}
The \emph{default local trigger orientation} is the working hypothesis that the first explicit local closure-phase trigger calibration should satisfy:
\begin{enumerate}
    \item the closure tolerance is the stricter burden tolerance,
    \[
    \varepsilon_{\mathrm{CL}} \le \varepsilon_{\Theta};
    \]
    \item the safe admissibility region for \(\rho_A\) is taken as a narrow stability band around the current admissible packet-side monitor regime;
    \item the minimum persistence length is not read as one-step only, but as a genuinely multi-step window
    \[
    m_{\min}>1.
    \]
\end{enumerate}
\end{definition}

\begin{remark}[Why closure tolerance should be stricter first]
\label{rem:why-closure-tolerance-should-be-stricter-first}
The current obstruction analysis suggests that the coupled local burden is not blocked first by lack of Theta-side machine decodability, but by insufficiently stabilized closure-phase persistence.
Hence the first local calibration should be closure-tight rather than Theta-tight.
\end{remark}

\begin{definition}[First plausible default trigger package]
\label{def:first-plausible-default-trigger-package}
The \emph{first plausible default trigger package} is any trigger-calibration data package
\[
\mathfrak{P}_{\mathrm{trig}}^{\mathrm{loc,def}}
:=
\bigl(
\varepsilon_{\mathrm{CL}}^{\mathrm{def}},
\varepsilon_{\Theta}^{\mathrm{def}},
\mathfrak{A}_{\mathrm{adm}}^{\mathrm{safe,def}},
 m_{\min}^{\mathrm{def}}
\bigr)
\]
satisfying:
\begin{enumerate}
    \item
    \[
    0<\varepsilon_{\mathrm{CL}}^{\mathrm{def}} \le \varepsilon_{\Theta}^{\mathrm{def}};
    \]
    \item
    \[
    \mathfrak{A}_{\mathrm{adm}}^{\mathrm{safe,def}}
    \subseteq
    \mathfrak{A}_{\mathrm{adm}};
    \]
    \item
    \[
    m_{\min}^{\mathrm{def}}\ge 2.
    \]
\end{enumerate}
\end{definition}

\begin{remark}[Interpretation of the default package]
\label{rem:interpretation-of-the-default-package}
This package does not yet assign final numerical values.
It only fixes the first local reading:
closure burden must enter a stricter small-burden regime than the routed Theta burden, and one-step near-entry is not yet trusted as a sufficient certificate.
\end{remark}

\begin{proposition}[Default local trigger package is obstruction-compatible]
\label{prop:default-local-trigger-package-is-obstruction-compatible}
Assume the first local obstruction classification theorem.
Then the first plausible default trigger package is compatible with the current obstruction reading.
\end{proposition}

\begin{proof}
The current obstruction has been classified as calibration/persistence type rather than machine-side decodability type.
A stricter closure tolerance together with a safe admissibility subregion and a persistence length at least two directly addresses that obstruction pattern: it stabilizes the closure burden first, keeps the packet-side monitor conservative, and excludes one-step transient near-entry events.
Hence the default package is compatible with the current obstruction reading.
\end{proof}

\begin{proposition}[Why one-step persistence should not be the default]
\label{prop:why-one-step-persistence-should-not-be-the-default}
Under the present obstruction-prior reading, the choice
\[
m_{\min}=1
\]
should not be the first default calibration.
\end{proposition}

\begin{proof}
The obstruction analysis explicitly identified trigger-persistence instability as part of the first local failure mode.
If one were to set \(m_{\min}=1\), then a single transient entrance into the local trigger inequality schema would already be counted as sufficient, which removes exactly the persistence safeguard that the current obstruction analysis says is missing.
Therefore one-step persistence is not the correct first default choice.
\end{proof}

\begin{theorem}[First default local trigger calibration theorem]
\label{thm:first-default-local-trigger-calibration-theorem}
Assume:
\begin{enumerate}
    \item the first local trigger calibration theorem applies;
    \item the current obstruction is of calibration/persistence type;
    \item the closure-phase-first local reading remains active.
\end{enumerate}
Then the first reasonable local default calibration should be chosen from a first plausible default trigger package
\[
\mathfrak{P}_{\mathrm{trig}}^{\mathrm{loc,def}}
=
\bigl(
\varepsilon_{\mathrm{CL}}^{\mathrm{def}},
\varepsilon_{\Theta}^{\mathrm{def}},
\mathfrak{A}_{\mathrm{adm}}^{\mathrm{safe,def}},
 m_{\min}^{\mathrm{def}}
\bigr)
\]
with
\[
0<\varepsilon_{\mathrm{CL}}^{\mathrm{def}} \le \varepsilon_{\Theta}^{\mathrm{def}},
\qquad
m_{\min}^{\mathrm{def}}\ge 2.
\]
In particular, the first local trigger calibration is closure-prior and persistence-prior rather than one-step permissive.
\end{theorem}

\begin{proof}
By the first local trigger calibration theorem, what is missing is an explicit trigger-calibration package.
By the obstruction classification, the leading obstruction is calibration/persistence type and not machine-side decodability.
Therefore the first calibration should tighten the closure burden first, allow at least as much tolerance on the routed Theta burden, keep the packet-side admissibility band conservative, and require a genuinely persistent window rather than a one-step fluctuation.
These conditions are exactly those of the first plausible default trigger package.
\end{proof}

\begin{corollary}[Current post-r15 default working rule]
\label{cor:current-post-r15-default-working-rule}
At the present stage, the first local trigger audit should proceed as if:
\begin{enumerate}
    \item closure smallness must be at least as strict as Theta smallness;
    \item the packet-side admissibility band must be chosen conservatively;
    \item the first honest local persistence length should be read as
    \[
    m_{\min}\ge 2.
    \]
\end{enumerate}
Only after this default calibration fails should one relax toward a more permissive local trigger reading.
\end{corollary}

\begin{proof}
Immediate from Theorem~\ref{thm:first-default-local-trigger-calibration-theorem}.
\end{proof}


\begin{definition}[Discretized Teichm\"uller persistence stream]
\label{def:discretized-teichmueller-persistence-stream}
Let
\[
X_n
\]
denote the current local machine-fused state and let
\[
T_N(X_n)
\]
be the local Monte-Carlo Teichm\"uller distortion surrogate.

The \emph{discretized Teichm\"uller persistence stream} is the symbolic stream
\[
\Sigma_{\mathrm{pers}}
=
(\sigma_n)_{n\ge 0},
\]
where each symbol \(\sigma_n\) records the discretized local status of
\[
\bigl(
\Delta_{\mathrm{CL}}(n),
\Delta_{\Theta R}(T)(n),
\rho_A(n),
T_N(X_n),
T_N(X_{n+1}),
|T_N(X_{n+1})-T_N(X_n)|
\bigr)
\]
relative to the active local trigger thresholds.
\end{definition}

\begin{definition}[Kasiski persistence scan]
\label{def:kasiski-persistence-scan}
The \emph{Kasiski persistence scan} is the search for repeated symbolic blocks inside
\[
\Sigma_{\mathrm{pers}}
\]
in order to detect recurring local trigger-near, persistence-fail, or admissibility-fail patterns.
\end{definition}

\begin{remark}[Role of Kasiski]
\label{rem:role-of-kasiski}
Kasiski is used here as a repetition and spacing detector on the discretized persistence stream.
It does not decide admissibility; it only detects recurring local obstruction patterns.
\end{remark}

\begin{definition}[Hashed persistence window]
\label{def:hashed-persistence-window}
A \emph{hashed persistence window} of length \(m\) is a hash value assigned to a local window
\[
(\sigma_n,\sigma_{n+1},\dots,\sigma_{n+m-1})
\]
of the discretized Teichm\"uller persistence stream.
\end{definition}

\begin{definition}[Markov persistence surface]
\label{def:markov-persistence-surface}
The \emph{Markov persistence surface} is the transition surface obtained from the empirical transition frequencies between hashed persistence windows or their reduced regime classes.
\end{definition}

\begin{remark}[Role of Markov control]
\label{rem:role-of-markov-control}
The Markov layer is not a primary admissibility criterion.
It is only the readable transition shadow of the shell-/trigger-restricted local persistence dynamics.
\end{remark}

\begin{definition}[Newton persistence refinement]
\label{def:newton-persistence-refinement}
The \emph{Newton persistence refinement} is a local Newton-type correction step applied only after a persistence regime has already been selected and a smooth surrogate defect function has been chosen.
\end{definition}

\begin{remark}[Why Newton comes last]
\label{rem:why-newton-comes-last}
Newton is not used to discover the regime.
It is used only to sharpen the local estimate once the regime has already been isolated by Kasiski/hash/Markov filtering.
\end{remark}

\begin{definition}[Bombe-sieve persistence elimination layer]
\label{def:bombe-sieve-persistence-elimination}
The \emph{bombe-sieve persistence elimination layer} is the contradiction-sensitive elimination of local continuation menus by:
\begin{enumerate}
    \item a trigger-sensitive sieve,
    \item bombe-style rejection of locally incompatible continuation branches,
    \item optional nested backtracking across deeper local menus.
\end{enumerate}
\end{definition}

\begin{definition}[Spectral menu readout]
\label{def:spectral-menu-readout}
A \emph{spectral menu readout} is the local ranking of a shell-admissible continuation menu by coding it into a self-adjoint menu operator and reading the preferred continuation as the top eigendirection.
\end{definition}

\begin{definition}[Vernam-like persistence stream heuristic]
\label{def:vernam-like-persistence-stream-heuristic}
The \emph{Vernam-like persistence stream heuristic} is the working interpretation that the discretized Teichm\"uller persistence stream behaves like an observable stream generated by a hidden local structural source, so that repeated local obstruction patterns may be used to infer the generator rather than merely the visible symbols.
\end{definition}

\begin{remark}[How to read the Vernam-like heuristic]
\label{rem:how-to-read-the-vernam-like-heuristic}
This is not an identity claim with literal cryptographic stream ciphers.
It is a search heuristic: the visible persistence symbols are treated as a surface stream whose hidden local generator is to be reconstructed from repetition, spacing, transition structure, and contradiction-sensitive elimination.
\end{remark}

\begin{definition}[Crypto-guided persistence analysis pipeline]
\label{def:crypto-guided-persistence-analysis-pipeline}
The \emph{crypto-guided persistence analysis pipeline} is the local tool chain
\[
\text{Kasiski}
\to
\text{hash}
\to
\text{Markov}
\to
\text{Bombe/Sieve}
\to
\text{spectral menu readout}
\to
\text{Newton refinement},
\]
applied to the discretized Teichm\"uller persistence stream.
\end{definition}

\begin{theorem}[First crypto-guided persistence scan theorem]
\label{thm:first-crypto-guided-persistence-scan-theorem}
Assume:
\begin{enumerate}
    \item a local Monte-Carlo Teichm\"uller persistence stream has been constructed;
    \item the active closure/Theta/admissibility thresholds are fixed locally;
    \item admissibility remains decided by closure-compatible trigger logic.
\end{enumerate}
Then the crypto-guided persistence analysis pipeline is admissible as a subordinate persistence-analysis layer.
It may be used to detect recurring persistence-failure patterns, compress local continuation menus, and sharpen local trigger calibration without replacing the closure-compatible admissibility criterion.
\end{theorem}

\begin{proof}
The earlier tool-guided compression logic already admits Kasiski, hash/rainbow compression, Markov guidance, contradiction-sensitive elimination, spectral ranking, and Newton refinement provided they remain subordinate to structural admissibility.
Here the same discipline is applied to the discretized Teichm\"uller persistence stream.
Hence the pipeline is admissible as a subordinate persistence-analysis layer and may be used to detect, compress, rank, and refine local persistence behavior without replacing closure-compatible trigger logic itself.
\end{proof}

\begin{corollary}[Current post-r16 persistence working rule]
\label{cor:current-post-r16-persistence-working-rule}
The next local persistence audit should not rely on pointwise burden smallness alone.
It should:
\begin{enumerate}
    \item build the discretized Teichm\"uller persistence stream;
    \item scan it by Kasiski for recurring local obstruction patterns;
    \item hash and compress local windows;
    \item read the resulting regime dynamics through a Markov persistence surface;
    \item eliminate locally contradictory branches by bombe/sieve filtering;
    \item rank the surviving continuation menu spectrally;
    \item apply Newton refinement only at the final local sharpening stage.
\end{enumerate}
\end{corollary}

\begin{proof}
Immediate from Theorem~\ref{thm:first-crypto-guided-persistence-scan-theorem}.
\end{proof}



\begin{definition}[Local persistence symbols]
\label{def:local-persistence-symbols}
Let \(\Sigma_{\mathrm{pers}}=(\sigma_n)_{n\ge 0}\) be the discretized Teichm\"uller persistence stream.
The following local symbols are distinguished:
\begin{enumerate}
    \item \(\mathtt{NT}\) (near-trigger): the local burden pair satisfies the pointwise closure/Theta smallness bounds but the minimal persistence window is not yet certified;
    \item \(\mathtt{PF}\) (persistence-fail): the next local state leaves the candidate closure-phase trigger slice or fails the Teichm\"uller persistence control while remaining machine-decodable;
    \item \(\mathtt{CF}\) (calibration-fail): one of the active closure/Theta trigger inequalities is not yet met;
    \item \(\mathtt{AF}\) (admissibility-fail): the packet-side monitor leaves the safe admissibility zone;
    \item \(\mathtt{EN}\) (entry): the local state lies in the calibrated closure-phase trigger slice.
\end{enumerate}
\end{definition}

\begin{definition}[First recurring local failure motif]
\label{def:first-recurring-local-failure-motif}
The \emph{first recurring local failure motif} is the shortest repeated symbolic block
\[
\omega_{\mathrm{fail}}
=
(\sigma_n,\dots,\sigma_{n+\ell-1})
\]
with \(\ell\ge 2\) detected by the Kasiski persistence scan and interpreted as the first persistent local obstruction pattern.
\end{definition}

\begin{definition}[Near-trigger / persistence-fail motif]
\label{def:near-trigger-persistence-fail-motif}
The \emph{near-trigger / persistence-fail motif} is the symbolic block
\[
\omega_{\mathrm{NT\to PF}}
:=
(\mathtt{NT},\mathtt{PF}).
\]
It records a local run in which the first step reaches pointwise trigger-near smallness, but the second step fails to preserve the closure-phase trigger slice or the corresponding Teichm\"uller persistence control.
\end{definition}

\begin{remark}[Why this is the natural first failure candidate]
\label{rem:why-natural-first-failure-candidate}
The current post-v3.32.0 line is obstruction-prior rather than entry-prior, and its first classified obstruction is of trigger calibration/persistence type rather than machine-side decodability type.
Hence the first repeated failure pattern should be sought among motifs that are already near local trigger conditions but fail on the required two-step persistence window.
\end{remark}

\begin{proposition}[Machine-side decodability is not the first recurring failure source]
\label{prop:machine-side-decodability-not-first-recurring-failure-source}
Assume the prime-lattice machine realization theorem and the first machine-fused local OP-entry protocol theorem.
Then the first recurring local failure motif is not forced by machine-side decodability failure alone.
\end{proposition}

\begin{proof}
Under the machine-fused realization hypotheses, descent, trigger testing, and branching are already available as finite machine-controlled operations on the prime-lattice machine object.
Therefore a repeated first failure pattern is not naturally attributed to absence of local machine decodability itself, but to failure of trigger calibration or persistence within the already machine-realizable regime.
\end{proof}

\begin{theorem}[First recurring local persistence-failure reading]
\label{thm:first-recurring-local-persistence-failure-reading}
Assume:
\begin{enumerate}
    \item the closure-phase-first local audit remains the preferred default reading;
    \item the default local trigger orientation with \(\varepsilon_{\mathrm{CL}}\le \varepsilon_{\Theta}\) and \(m_{\min}\ge 2\) is adopted;
    \item the current line is obstruction-prior rather than entry-prior;
    \item the crypto-guided persistence analysis pipeline is applied to the local discretized Teichm\"uller persistence stream.
\end{enumerate}
Then the first recurring local failure motif is read by default as
\[
\omega_{\mathrm{fail}}
=
\omega_{\mathrm{NT\to PF}}
=
(\mathtt{NT},\mathtt{PF}),
\]
that is: a trigger-near first step followed by failure of two-step closure-phase persistence.
\end{theorem}

\begin{proof}
By obstruction-priority, immediate entry is not the default first local outcome.
By the earlier obstruction classification, the first local obstruction is of trigger calibration/persistence type rather than machine-side decodability type.
Under the default trigger orientation and the requirement \(m_{\min}\ge 2\), the most economical local failure pattern is therefore not a one-step calibration collapse before near-trigger status is reached, but a run in which the first step attains trigger-near pointwise smallness and the second step fails to preserve the closure-phase trigger slice.
The crypto-guided persistence scan is designed exactly to detect repeated symbolic blocks of this form.
Hence the default first recurring local failure motif is \((\mathtt{NT},\mathtt{PF})\).
\end{proof}

\begin{corollary}[Current first persistence-stream reading]
\label{cor:current-first-persistence-stream-reading}
The current post-v3.32.0 persistence line should be read first through the motif
\[
(\mathtt{NT},\mathtt{PF}),
\]
rather than through immediate entry or pure decodability failure.
Equivalently, the present first local task is to identify which closure-phase persistence ingredient breaks between the near-trigger first step and the failed second step.
\end{corollary}

\begin{proof}
Immediate from Theorem~\ref{thm:first-recurring-local-persistence-failure-reading} and Proposition~\ref{prop:machine-side-decodability-not-first-recurring-failure-source}.
\end{proof}


\begin{definition}[Theta-drift persistence-fail subtype]
\label{def:theta-drift-persistence-fail-subtype}
The symbol
\[
\mathtt{PF}_{\Theta}
\]
denotes the persistence-fail subtype in which:
\begin{enumerate}
    \item the first local step is near-trigger;
    \item the second local step fails because the phase-side component loses the required closure-phase persistence control;
    \item the packet-side admissibility monitor remains inside the safe zone.
\end{enumerate}
\end{definition}

\begin{definition}[Teichm\"uller-distortion persistence-fail subtype]
\label{def:teichmueller-distortion-persistence-fail-subtype}
The symbol
\[
\mathtt{PF}_{T}
\]
denotes the persistence-fail subtype in which:
\begin{enumerate}
    \item the first local step is near-trigger;
    \item the pointwise burden inequalities remain approximately trigger-near;
    \item the second local step fails because the local Teichm\"uller distortion surrogate or its one-step variation leaves the allowed persistence band.
\end{enumerate}
\end{definition}

\begin{definition}[Admissibility-edge persistence-fail subtype]
\label{def:admissibility-edge-persistence-fail-subtype}
The symbol
\[
\mathtt{PF}_{A}
\]
denotes the persistence-fail subtype in which:
\begin{enumerate}
    \item the first local step is near-trigger;
    \item the second local step fails because the packet-side monitor \(\rho_A\) reaches or leaves the safe admissibility edge;
    \item the corresponding closure-phase persistence window therefore collapses for admissibility reasons.
\end{enumerate}
\end{definition}

\begin{remark}[How the three subtypes differ]
\label{rem:how-the-three-subtypes-differ}
The subtype \(\mathtt{PF}_{\Theta}\) is phase-led and burden-side.
The subtype \(\mathtt{PF}_{T}\) is geometric and distortion-led.
The subtype \(\mathtt{PF}_{A}\) is guard-led and packet-admissibility driven.
\end{remark}

\begin{proposition}[Admissibility-edge contact is not the default first persistence-fail subtype]
\label{prop:admissibility-edge-contact-not-default-first-pf-subtype}
Assume the current line is read through the default local trigger orientation with conservative packet-side admissibility control.
Then \(\mathtt{PF}_{A}\) is not the default first persistence-fail subtype.
\end{proposition}

\begin{proof}
The current line treats \(\rho_A\) primarily as an admissibility monitor rather than as a coequal core burden, and the active default trigger package keeps the safe admissibility band deliberately conservative.
Hence the first local failure is not naturally read as immediate edge contact of \(\rho_A\), but as failure of the closure-phase persistence mechanism before the admissibility edge is reached.
\end{proof}

\begin{proposition}[The first persistence-fail subtype is phase/geometry-prior]
\label{prop:first-persistence-fail-subtype-is-phase-geometry-prior}
Assume the closure-phase-first local audit, the default trigger orientation, and the first recurring persistence-failure reading
\[
(\mathtt{NT},\mathtt{PF}).
\]
Then the first persistence-fail subtype is read by default as phase/geometry-prior, i.e. as lying in the reduced pair
\[
\{\mathtt{PF}_{\Theta},\mathtt{PF}_{T}\},
\]
rather than as an admissibility-edge subtype \(\mathtt{PF}_{A}\).
\end{proposition}

\begin{proof}
Under the closure-phase-first reading, the active mixed persistence burden is carried first by the closure/Theta pair, while \(\rho_A\) remains a guard variable.
By Proposition~\ref{prop:admissibility-edge-contact-not-default-first-pf-subtype}, immediate admissibility-edge contact is not the default first subtype.
Hence the first persistence-fail subtype is phase/geometry-prior.
\end{proof}

\begin{theorem}[First local persistence-fail subtype theorem]
\label{thm:first-local-persistence-fail-subtype-theorem}
Assume:
\begin{enumerate}
    \item the first recurring local failure motif is
    \[
    (\mathtt{NT},\mathtt{PF});
    \]
    \item the closure-phase-first local audit remains the default reading;
    \item the packet-side admissibility monitor remains conservative on the first near-trigger step.
\end{enumerate}
Then the first local persistence-fail subtype is not read by default as admissibility-edge contact.
Instead, it is read by default as a phase/geometry-prior break:
either
\[
\mathtt{PF}_{\Theta}
\]
or
\[
\mathtt{PF}_{T},
\]
with the current line favoring the reading that Theta-side drift is the primary burden-side break and Teichm\"uller distortion is its geometric witness.
\end{theorem}

\begin{proof}
By Proposition~\ref{prop:first-persistence-fail-subtype-is-phase-geometry-prior}, the first subtype lies in the pair \(\{\mathtt{PF}_{\Theta},\mathtt{PF}_{T}\}\) rather than in \(\mathtt{PF}_{A}\).
Under the closure-phase-first reading, the active burden is phase-led before it is envelope-led, while the Teichm\"uller control records the geometric persistence of that phase-led regime.
Therefore the default local reading is that Theta-drift is the primary burden-side break and Teichm\"uller distortion is the corresponding geometric witness of the failed second step.
\end{proof}

\begin{corollary}[Current post-r18 local diagnosis rule]
\label{cor:current-post-r18-local-diagnosis-rule}
The next local diagnosis should not first ask whether \(\rho_A\) touched the admissibility edge.
It should first ask:
\begin{enumerate}
    \item did the phase-side burden drift out of the closure-phase persistence window;
    \item and if so, what Teichm\"uller distortion witness recorded that loss of persistence?
\end{enumerate}
\end{corollary}

\begin{proof}
Immediate from Theorem~\ref{thm:first-local-persistence-fail-subtype-theorem}.
\end{proof}


\begin{definition}[Local Theta-drift witness]
\label{def:local-theta-drift-witness}
A \emph{local Theta-drift witness} is a finite witness object
\[
\mathfrak{W}_{\Theta}^{\mathrm{loc}}(n)
\]
attached to a near-trigger first step and the following failed second step, certifying that:
\begin{enumerate}
    \item the first step is of type \(\mathtt{NT}\);
    \item the second step is of type \(\mathtt{PF}_{\Theta}\) or phase/geometry-prior persistence failure;
    \item the phase-side burden leaves the active closure-phase persistence window on the second step while the packet-side admissibility monitor remains inside the conservative safe zone.
\end{enumerate}
\end{definition}

\begin{definition}[Local Teichm\"uller drift witness]
\label{def:local-teichmueller-drift-witness}
A \emph{local Teichm\"uller drift witness} is a finite witness object
\[
\mathfrak{W}_{T}^{\mathrm{loc}}(n)
\]
certifying that the local Teichm\"uller distortion surrogate or its one-step variation crosses the active persistence band between the near-trigger first step and the failed second step.
\end{definition}

\begin{remark}[Burden-side versus geometric witness]
\label{rem:burden-side-versus-geometric-witness}
The Theta-drift witness and the Teichm\"uller drift witness play different roles.
The first records the burden-side loss of closure-phase persistence.
The second records the geometric distortion signature of that same loss.
\end{remark}

\begin{definition}[Local coupled drift witness package]
\label{def:local-coupled-drift-witness-package}
The \emph{local coupled drift witness package} is the pair
\[
\mathfrak{W}_{\mathrm{drift}}^{\mathrm{loc}}(n)
:=
\bigl(
\mathfrak{W}_{\Theta}^{\mathrm{loc}}(n),
\mathfrak{W}_{T}^{\mathrm{loc}}(n)
\bigr),
\]
read on the current persistence stream whenever the local failure motif
\[
(\mathtt{NT},\mathtt{PF})
\]
is refined to a phase/geometry-prior break.
\end{definition}

\begin{definition}[Local Theta-drift witness test]
\label{def:local-theta-drift-witness-test}
The \emph{local Theta-drift witness test} asks whether, for some local index \(n\):
\begin{enumerate}
    \item the first step satisfies the active near-trigger inequalities;
    \item the second step violates the phase-side persistence requirement before admissibility edge contact occurs;
    \item a Teichm\"uller distortion witness is simultaneously present or becomes newly active on that second step.
\end{enumerate}
\end{definition}

\begin{proposition}[Theta-drift witness packages refine the default persistence-fail reading]
\label{prop:theta-drift-witness-packages-refine-default-persistence-fail-reading}
Assume the first recurring local failure motif is
\[
(\mathtt{NT},\mathtt{PF})
\]
and the first persistence-fail subtype is phase/geometry-prior.
Then every successful local Theta-drift witness test yields a local coupled drift witness package
\[
\mathfrak{W}_{\mathrm{drift}}^{\mathrm{loc}}(n)
\]
refining the persistence-fail reading.
\end{proposition}

\begin{proof}
If the local Theta-drift witness test succeeds, then the failed second step is not merely an undifferentiated persistence failure: it is a phase-side drift event accompanied by a geometric distortion witness while admissibility remains conservative.
That is exactly the data recorded by the two witness objects and therefore by their paired witness package.
\end{proof}

\begin{theorem}[First local Theta-drift witness theorem]
\label{thm:first-local-theta-drift-witness-theorem}
Assume:
\begin{enumerate}
    \item the current post-v3.32.0 persistence line is read through the first recurring motif
    \[
    (\mathtt{NT},\mathtt{PF});
    \]
    \item the first persistence-fail subtype is phase/geometry-prior;
    \item the local Theta-drift witness test succeeds at some first relevant local index \(n\).
\end{enumerate}
Then the current line admits a first local Theta-drift witness package
\[
\mathfrak{W}_{\mathrm{drift}}^{\mathrm{loc}}(n)
=
\bigl(
\mathfrak{W}_{\Theta}^{\mathrm{loc}}(n),
\mathfrak{W}_{T}^{\mathrm{loc}}(n)
\bigr),
\]
with the following reading:
\begin{enumerate}
    \item the first failed second step is burden-side dominated by Theta-drift;
    \item the corresponding Teichm\"uller witness is the geometric recorder of that burden-side drift;
    \item the first local obstruction is therefore not yet read as pure trigger-calibration failure and not yet as admissibility-edge collapse, but as a drift-mediated persistence break.
\end{enumerate}
\end{theorem}

\begin{proof}
By the first recurring persistence-failure reading, the current line is governed first by the motif
\[
(\mathtt{NT},\mathtt{PF}).
\]
By the first local persistence-fail subtype theorem, the first persistence-fail subtype is read as phase/geometry-prior rather than admissibility-edge driven.
If the local Theta-drift witness test succeeds at the first relevant index, then the failed second step is specifically identified as a phase-side drift event with an accompanying Teichm\"uller distortion witness.
Hence the local coupled drift witness package exists and carries the stated interpretation.
\end{proof}

\begin{corollary}[Current post-r19 drift diagnosis rule]
\label{cor:current-post-r19-drift-diagnosis-rule}
Once the first local Theta-drift witness package is available, the next local audit should ask first which phase-side persistence ingredient must be strengthened in order to keep the second step inside the closure-phase window, and only afterwards whether a further trigger-calibration reduction is needed.
\end{corollary}

\begin{proof}
Immediate from Theorem~\ref{thm:first-local-theta-drift-witness-theorem}.
\end{proof}


\begin{definition}[Local phase-side reinforcement datum]
\label{def:local-phase-side-reinforcement-datum}
A \emph{local phase-side reinforcement datum} at a first relevant local index $n$ is a finite correction datum
\[
\mathfrak{R}_{\Theta}^{\mathrm{loc}}(n)
:=
\bigl(
\eta_{\Theta}(n),
\eta_{T}(n)
\bigr),
\]
where:
\begin{enumerate}
    \item $\eta_{\Theta}(n)\ge 0$ is the minimal phase-side persistence strengthening required to keep the second step inside the active closure-phase window;
    \item $\eta_{T}(n)\ge 0$ is the corresponding geometric tolerance compensation required so that the Teichm\"uller distortion witness stays inside the active persistence band.
\end{enumerate}
\end{definition}

\begin{definition}[Local repaired second step]
\label{def:local-repaired-second-step}
A failed second step at index $n+1$ is called \emph{locally repaired} by a reinforcement datum
\[
\mathfrak{R}_{\Theta}^{\mathrm{loc}}(n)
=
\bigl(
\eta_{\Theta}(n),
\eta_{T}(n)
\bigr)
\]
if, after strengthening the phase-side persistence requirement by $\eta_{\Theta}(n)$ and compensating the active Teichm\"uller persistence band by $\eta_{T}(n)$, the updated second step:
\begin{enumerate}
    \item remains in the closure-phase trigger window;
    \item preserves admissibility of $\rho_A$;
    \item no longer issues the persistence-fail symbol $\mathtt{PF}$.
\end{enumerate}
\end{definition}

\begin{remark}[Meaning of local repair]
\label{rem:meaning-of-local-repair}
The repair datum does not assert that the original failed step was already admissible.
It records the smallest phase-side and geometric strengthening under which the same local two-step pattern would have stayed inside the active closure-phase window.
\end{remark}

\begin{definition}[Minimal local phase-side reinforcement]
\label{def:minimal-local-phase-side-reinforcement}
A local phase-side reinforcement datum
\[
\mathfrak{R}_{\Theta}^{\mathrm{loc}}(n)
=
\bigl(
\eta_{\Theta}(n),
\eta_{T}(n)
\bigr)
\]
is called \emph{minimal} if no strictly smaller pair in the product order on $\mathbb{R}_{\ge 0}^{2}$ already yields a locally repaired second step.
\end{definition}

\begin{definition}[First local repair witness]
\label{def:first-local-repair-witness}
A \emph{first local repair witness} is a pair
\[
\mathfrak{W}_{\mathrm{rep}}^{\mathrm{loc}}(n)
:=
\bigl(
\mathfrak{W}_{\mathrm{drift}}^{\mathrm{loc}}(n),
\mathfrak{R}_{\Theta}^{\mathrm{loc}}(n)
\bigr)
\]
consisting of:
\begin{enumerate}
    \item the first local coupled drift witness package;
    \item a minimal local phase-side reinforcement datum repairing the failed second step.
\end{enumerate}
\end{definition}

\begin{proposition}[Theta-drift witnesses induce local repair data]
\label{prop:theta-drift-witnesses-induce-local-repair-data}
Assume the first local Theta-drift witness theorem applies.
Then every first local Theta-drift witness package determines a nonnegative local phase-side reinforcement datum.
\end{proposition}

\begin{proof}
The witness package identifies both the burden-side Theta drift and the corresponding Teichm\"uller distortion contribution at the failed second step while keeping admissibility conservative.
Hence one may measure how much phase-side strengthening and geometric tolerance compensation are needed to return the same local second step to the active closure-phase window.
These quantities are nonnegative by construction.
\end{proof}

\begin{lemma}[First local repair lemma]
\label{lem:first-local-repair-lemma}
Assume:
\begin{enumerate}
    \item the first local Theta-drift witness theorem applies at index $n$;
    \item a minimal local phase-side reinforcement datum
    \[
    \mathfrak{R}_{\Theta}^{\mathrm{loc}}(n)
    =
    \bigl(
    \eta_{\Theta}(n),
    \eta_{T}(n)
    \bigr)
    \]
    exists.
\end{enumerate}
Then the associated first local repair witness
\[
\mathfrak{W}_{\mathrm{rep}}^{\mathrm{loc}}(n)
\]
certifies the smallest local strengthening currently needed to keep the second step inside the closure-phase persistence window.
\end{lemma}

\begin{proof}
By the first local Theta-drift witness theorem, the failed second step is already isolated as a Theta-drift event with Teichm\"uller distortion witness.
A minimal reinforcement datum records the smallest phase-side and geometric compensation required to restore persistence of that step.
Therefore the combined repair witness certifies the smallest currently needed local strengthening.
\end{proof}

\begin{theorem}[First local phase-side repair theorem]
\label{thm:first-local-phase-side-repair-theorem}
Assume:
\begin{enumerate}
    \item the first recurring local failure motif is
    \[
    (\mathtt{NT},\mathtt{PF});
    \]
    \item the first persistence-fail subtype is phase/geometry-prior;
    \item the first local Theta-drift witness theorem applies at index $n$;
    \item a minimal local phase-side reinforcement datum exists at that same index.
\end{enumerate}
Then the current local persistence obstruction is repair-readable.
More precisely, the first local failure is governed by a minimal reinforcement package
\[
\mathfrak{R}_{\Theta}^{\mathrm{loc}}(n)
=
\bigl(
\eta_{\Theta}(n),
\eta_{T}(n)
\bigr)
\]
which measures exactly how much phase-side persistence strengthening and geometric tolerance compensation are missing at the failed second step.
\end{theorem}

\begin{proof}
The first three assumptions isolate the first failed second step as a Theta-drift persistence failure with Teichm\"uller witness.
The existence of a minimal reinforcement datum then identifies the smallest correction under which the same local two-step pattern would remain inside the closure-phase window.
Therefore the current local obstruction is not only diagnosable but also repair-readable in the stated sense.
\end{proof}

\begin{corollary}[Current post-r20 repair rule]
\label{cor:current-post-r20-repair-rule}
After the first local Theta-drift witness has been isolated, the next local audit should not first ask for a new global architecture.
It should ask for the minimal local reinforcement package
\[
\bigl(
\eta_{\Theta}(n),
\eta_{T}(n)
\bigr)
\]
needed to keep the second step inside the active closure-phase persistence window.
\end{corollary}

\begin{proof}
Immediate from Theorem~\ref{thm:first-local-phase-side-repair-theorem}.
\end{proof}


\begin{definition}[Theta-dominant local repair]
\label{def:theta-dominant-local-repair}
A minimal local reinforcement package
\[
\mathfrak{R}_{\Theta}^{\mathrm{loc}}(n)
=
\bigl(
\eta_{\Theta}(n),
\eta_{T}(n)
\bigr)
\]
is called \emph{Theta-dominant} if:
\begin{enumerate}
    \item \(\eta_{\Theta}(n) > 0\);
    \item \(\eta_{T}(n)\ge 0\);
    \item the failed second step cannot be repaired by geometric compensation alone when \(\eta_{\Theta}(n)\) is set to \(0\).
\end{enumerate}
\end{definition}

\begin{definition}[Geometrically carried local repair]
\label{def:geometrically-carried-local-repair}
A minimal local reinforcement package
\[
\mathfrak{R}_{\Theta}^{\mathrm{loc}}(n)
=
\bigl(
\eta_{\Theta}(n),
\eta_{T}(n)
\bigr)
\]
is called \emph{geometrically carried} if:
\begin{enumerate}
    \item \(\eta_{T}(n) > 0\);
    \item the corresponding Teichm\"uller witness is an essential part of the repair;
    \item the repaired second step cannot be restored by burden-side phase strengthening alone while keeping the active geometric tolerance unchanged.
\end{enumerate}
\end{definition}

\begin{remark}[Two-sided reading of the repair datum]
\label{rem:two-sided-reading-of-the-repair-datum}
The local repair datum measures two distinct but coupled deficits:
\begin{enumerate}
    \item burden-side phase persistence deficit, read through \(\eta_{\Theta}(n)\);
    \item geometric persistence deficit, read through \(\eta_{T}(n)\).
\end{enumerate}
The point of the current audit is to determine which of these two deficits is structurally primary in the first failed second step.
\end{remark}

\begin{definition}[Locally coupled repair orientation]
\label{def:locally-coupled-repair-orientation}
The first local repair orientation at index \(n\) is called:
\begin{enumerate}
    \item \emph{burden-dominant} if the minimal repair package is Theta-dominant and not geometrically carried;
    \item \emph{geometric-dominant} if the minimal repair package is geometrically carried and not Theta-dominant;
    \item \emph{genuinely coupled} if the minimal repair package is both Theta-dominant and geometrically carried.
\end{enumerate}
\end{definition}

\begin{proposition}[Theta-drift witnesses force nonzero burden-side repair]
\label{prop:theta-drift-witnesses-force-nonzero-burden-side-repair}
Assume the first local Theta-drift witness theorem applies at index \(n\).
Then every minimal local reinforcement package at that index satisfies
\[
\eta_{\Theta}(n) > 0.
\]
\end{proposition}

\begin{proof}
The first local Theta-drift witness theorem identifies the failed second step specifically as a phase-side drift event with Teichm\"uller witness.
If one set \(\eta_{\Theta}(n)=0\), the burden-side phase drift would remain uncorrected, so the failed second step could not be repaired in the sense of the local repair definition.
Hence every minimal repair package must have strictly positive \(\eta_{\Theta}(n)\).
\end{proof}

\begin{proposition}[Teichm\"uller witnesses make the repair geometrically visible]
\label{prop:teichmueller-witnesses-make-the-repair-geometrically-visible}
Assume the first local Theta-drift witness theorem applies at index \(n\) and the corresponding failed second step carries a nontrivial Teichm\"uller witness.
Then every minimal local reinforcement package at that index satisfies
\[
\eta_{T}(n)\ge 0
\]
with geometric significance, and the resulting repair orientation cannot be read as purely burden-side in the naive scalar sense.
\end{proposition}

\begin{proof}
A nontrivial Teichm\"uller witness means that the failed second step is accompanied by a geometric persistence defect, not only by a symbolic burden-side drift.
Therefore the repair package necessarily carries a geometric component, even if its size is smaller than the burden-side one.
Thus the repair is not purely scalar/burden-side in the naive sense.
\end{proof}

\begin{theorem}[First local repair-orientation theorem]
\label{thm:first-local-repair-orientation-theorem}
Assume:
\begin{enumerate}
    \item the first local Theta-drift witness theorem applies at index \(n\);
    \item a minimal local phase-side reinforcement datum exists at that index;
    \item the corresponding failed second step carries a nontrivial Teichm\"uller witness.
\end{enumerate}
Then the first local repair orientation is not geometric-dominant.
More precisely:
\begin{enumerate}
    \item the minimal repair package is necessarily Theta-dominant;
    \item the geometric contribution remains a genuine witness of the repair problem;
    \item hence the first local repair is read either as burden-dominant with geometric witness or as genuinely coupled, but not as purely geometric.
\end{enumerate}
\end{theorem}

\begin{proof}
By Proposition~\ref{prop:theta-drift-witnesses-force-nonzero-burden-side-repair}, one has \(\eta_{\Theta}(n)>0\) for every minimal repair package.
Hence the repair cannot be geometric-dominant in the sense of Definition~\ref{def:locally-coupled-repair-orientation}.
By Proposition~\ref{prop:teichmueller-witnesses-make-the-repair-geometrically-visible}, the Teichm\"uller component remains a genuine geometric witness of the same local failure.
Therefore the first local repair is either burden-dominant with geometric witness or genuinely coupled, but not purely geometric.
\end{proof}

\begin{corollary}[Current post-r21 repair reading]
\label{cor:current-post-r21-repair-reading}
The current post-v3.32.0 local repair problem should be read first as a phase-side burden repair problem with a genuine Teichm\"uller witness.
Accordingly, the next local audit should ask:
\begin{enumerate}
    \item how large the minimal burden-side strengthening \(\eta_{\Theta}(n)\) must be;
    \item whether the corresponding geometric term \(\eta_{T}(n)\) is only subordinate or already large enough to force a genuinely coupled repair reading.
\end{enumerate}
\end{corollary}

\begin{proof}
Immediate from Theorem~\ref{thm:first-local-repair-orientation-theorem}.
\end{proof}




\begin{definition}[Locally plausible tightened trigger package]
\label{def:locally-plausible-tightened-trigger-package}
The locally tightened trigger package
\[
\mathfrak{P}_{\mathrm{trig}}^{\sharp}(n)
=
\bigl(
\varepsilon_{\mathrm{CL}}^{\sharp},
\varepsilon_{\Theta}^{\sharp},
\varepsilon_{T}^{\sharp},
\mathfrak{A}_{\mathrm{adm}}^{\mathrm{safe}},
m_{\min}
\bigr)
\]
at index $n$ is called \emph{locally plausible} if:
\begin{enumerate}
    \item all tightened thresholds remain positive;
    \item the admissibility band $\mathfrak{A}_{\mathrm{adm}}^{\mathrm{safe}}$ remains nonempty;
    \item the tightened package is compatible with the current closure-phase regime at the first step.
\end{enumerate}
\end{definition}

\begin{definition}[Residual tightened-package insufficiency]
\label{def:residual-tightened-package-insufficiency}
A \emph{residual tightened-package insufficiency} at index $n$ is any component of the locally tightened trigger package whose strengthened form still fails to support the repaired two-step closure-phase persistence test.

The possible insufficiency types are:
\begin{enumerate}
    \item closure-side insufficiency;
    \item Theta-side insufficiency;
    \item Teichm\"uller-side insufficiency;
    \item admissibility-band insufficiency;
    \item persistence-length insufficiency.
\end{enumerate}
\end{definition}

\begin{remark}[What the new audit should decide]
\label{rem:what-the-new-audit-should-decide}
After constructing the first repair-induced tightening, the question is no longer whether the original trigger package fails.
The sharper question is which component of the tightened package still blocks the repaired second step.
\end{remark}

\begin{definition}[Teichm\"uller-persistence residual insufficiency]
\label{def:teichmueller-persistence-residual-insufficiency}
The residual insufficiency is called \emph{Teichm\"uller-persistence type} if:
\begin{enumerate}
    \item the tightened closure-side threshold is already compatible with the repaired first step;
    \item the tightened Theta-side threshold is already compatible with the repaired first step;
    \item the admissibility band remains safe;
    \item but the repaired two-step window still fails because the Teichm\"uller-side bound or its two-step stability requirement is not yet strong enough.
\end{enumerate}
\end{definition}

\begin{proposition}[The tightened package is locally plausible]
\label{prop:the-tightened-package-is-locally-plausible}
Assume the first local repair-to-trigger calibration theorem applies at index $n$.
Then the tightened trigger package
\[
\mathfrak{P}_{\mathrm{trig}}^{\sharp}(n)
\]
is locally plausible.
\end{proposition}

\begin{proof}
By the repair-to-trigger calibration theorem, the tightened thresholds are obtained from a genuine local repair witness and remain admissible as a candidate package.
Hence the tightened package preserves positivity, keeps the admissibility band meaningful, and remains compatible with the first-step closure-phase regime.
\end{proof}

\begin{proposition}[Closure and admissibility are not the first remaining insufficiency]
\label{prop:closure-and-admissibility-are-not-the-first-remaining-insufficiency}
Assume:
\begin{enumerate}
    \item the first local Theta-drift witness theorem applies;
    \item the first local repair-orientation theorem applies;
    \item the tightened trigger package is locally plausible.
\end{enumerate}
Then the first residual tightened-package insufficiency is not closure-side and not admissibility-band type.
\end{proposition}

\begin{proof}
The local failure was already classified as phase/geometry-prior rather than closure-prior, and the local repair package was constructed without changing the closure-side threshold at first order.
Moreover the current line is not admissibility-edge-prior.
Therefore, after tightening, the first remaining insufficiency cannot honestly be read first as closure-side or admissibility-band type.
\end{proof}

\begin{theorem}[First tightened-package plausibility theorem]
\label{thm:first-tightened-package-plausibility-theorem}
Assume:
\begin{enumerate}
    \item the first local repair-to-trigger calibration theorem applies at index $n$;
    \item the first local repair-orientation theorem applies at index $n$;
    \item the current local obstruction is phase/geometry-prior rather than closure/admissibility-prior.
\end{enumerate}
Then the current post-v3.32.0 line supports the following reading of the locally tightened trigger package:
\begin{enumerate}
    \item the tightened package is locally plausible;
    \item if it still fails to restore two-step closure-phase persistence, then the first remaining insufficiency is no longer closure-side or admissibility-side;
    \item the first remaining insufficiency should be searched first on the Teichm\"uller-persistence side, i.e. in the tightened geometric threshold or its two-step stability requirement.
\end{enumerate}
\end{theorem}

\begin{proof}
Local plausibility is Proposition~\ref{prop:the-tightened-package-is-locally-plausible}.
By Proposition~\ref{prop:closure-and-admissibility-are-not-the-first-remaining-insufficiency}, closure-side and admissibility-band failure are excluded as the first residual obstruction.
The local repair orientation is Theta-dominant but geometrically carried, so once the primary burden-side strengthening has been applied, the next unresolved obstruction is searched first on the Teichm\"uller-persistence side.
\end{proof}

\begin{corollary}[Current post-r22 tightened-package reading]
\label{cor:current-post-r22-tightened-package-reading}
The current post-v3.32.0 line should read the first tightened trigger package as plausibly near-sufficient but not yet self-certifying.
If the repaired two-step persistence still fails, the next local audit should first test whether the tightened Teichm\"uller-side bound or its persistence clause remains insufficient.
\end{corollary}

\begin{proof}
Immediate from Theorem~\ref{thm:first-tightened-package-plausibility-theorem}.
\end{proof}



\begin{definition}[Local Teichm\"uller-side repair datum]\label{def:local-teichmueller-side-repair-datum}
Let
\[
\mathfrak{P}_{\mathrm{trig}}^{\sharp}(n)
=
\bigl(
\varepsilon_{\mathrm{CL}}^{\sharp},
\varepsilon_{\Theta}^{\sharp},
\varepsilon_T^{\sharp},
\mathfrak{A}_{\mathrm{adm}}^{\mathrm{safe}},
 m_{\min}
\bigr)
\]
be the locally tightened trigger package at index $n$.

A \emph{local Teichm\"uller-side repair datum} is a pair
\[
\mathfrak{R}_{T}^{\mathrm{loc}}(n)
:=
\bigl(
\zeta_T(n),
\lambda_T(n)
\bigr)
\]
with
\[
\zeta_T(n)\ge 0,
\qquad
\lambda_T(n)\ge 0,
\]
where:
\begin{enumerate}
    \item $\zeta_T(n)$ is the additional tightening of the geometric threshold;
    \item $\lambda_T(n)$ is the additional tightening of the two-step distortion-variation allowance.
\end{enumerate}
\end{definition}

\begin{definition}[Teichm\"uller-reinforced trigger package]\label{def:teichmueller-reinforced-trigger-package}
The \emph{Teichm\"uller-reinforced trigger package} at index $n$ is
\[
\mathfrak{P}_{\mathrm{trig}}^{\flat}(n)
:=
\bigl(
\varepsilon_{\mathrm{CL}}^{\sharp},
\varepsilon_{\Theta}^{\sharp},
\varepsilon_T^{\flat},
\mathfrak{A}_{\mathrm{adm}}^{\mathrm{safe}},
 m_{\min},
 \eta_T^{\flat}(n)
\bigr),
\]
where
\[
\varepsilon_T^{\flat}
:=
\varepsilon_T^{\sharp}-\zeta_T(n),
\qquad
\eta_T^{\flat}(n)
:=
\eta_T(n)-\lambda_T(n),
\]
and the reinforced package is required to remain positive and admissible.
\end{definition}

\begin{definition}[Teichm\"uller-reinforced two-step persistence test]\label{def:teichmueller-reinforced-two-step-persistence-test}
A two-step local window
\[
X_n,\;X_{n+1}
\]
passes the \emph{Teichm\"uller-reinforced two-step persistence test} if for both $k\in\{n,n+1\}$:
\[
\Delta_{\mathrm{CL}}(k)\le \varepsilon_{\mathrm{CL}}^{\sharp},
\qquad
\Delta_{\Theta R}(T)(k)\le \varepsilon_{\Theta}^{\sharp},
\qquad
T_N(X_k)\le \varepsilon_T^{\flat},
\qquad
\rho_A(k)\in \mathfrak{A}_{\mathrm{adm}}^{\mathrm{safe}},
\]
and additionally
\[
|T_N(X_{n+1})-T_N(X_n)|\le \eta_T^{\flat}(n).
\]
\end{definition}

\begin{remark}[Meaning of the Teichm\"uller-side repair]\label{rem:meaning-of-the-teichmueller-side-repair}
If the first tightened trigger package is already closure-compatible and Theta-compatible but still fails on the persistence side, then the next honest local move is no longer to tighten closure or admissibility first. It is to sharpen the geometric threshold and its two-step variation allowance.
\end{remark}

\begin{definition}[Minimal local Teichm\"uller-side repair]\label{def:minimal-local-teichmueller-side-repair}
A local Teichm\"uller-side repair datum
\[
\bigl(\zeta_T(n),\lambda_T(n)\bigr)
\]
is called \emph{minimal} if the corresponding Teichm\"uller-reinforced trigger package passes the two-step persistence test, but every componentwise smaller repair datum fails to do so.
\end{definition}

\begin{proposition}[Teichm\"uller-side repair leaves closure and Theta tightening fixed]\label{prop:teichmueller-side-repair-leaves-closure-and-theta-tightening-fixed}
Assume the first local repair-to-trigger calibration theorem applies and the first tightened-package plausibility theorem identifies a Teichm\"uller-persistence residual insufficiency at index $n$.
Then a local Teichm\"uller-side repair acts only on the geometric threshold and the two-step variation allowance, while leaving
\[
\varepsilon_{\mathrm{CL}}^{\sharp}
\qquad\text{and}\qquad
\varepsilon_{\Theta}^{\sharp}
\]
unchanged.
\end{proposition}

\begin{proof}
By hypothesis, the residual insufficiency is not closure-side and not admissibility-band type, and the first repair orientation is Theta-dominant but geometrically carried. Therefore the remaining local insufficiency is carried by the geometric persistence side, so the next tightening acts only on the Teichm\"uller-side threshold and its variation allowance.
\end{proof}

\begin{lemma}[First local Teichm\"uller-persistence repair lemma]\label{lem:first-local-teichmueller-persistence-repair-lemma}
Assume:
\begin{enumerate}
    \item the first tightened-package plausibility theorem applies at index $n$;
    \item the remaining insufficiency is of Teichm\"uller-persistence type;
    \item there exists a minimal local Teichm\"uller-side repair datum.
\end{enumerate}
Then the corresponding Teichm\"uller-reinforced trigger package is the first honest local candidate for restoring the missing two-step closure-phase persistence after the initial Theta-dominant repair.
\end{lemma}

\begin{proof}
The first tightened package already preserves closure-side and Theta-side compatibility and keeps the admissibility band safe. By assumption, the only remaining failure lies on the Teichm\"uller-persistence side. A minimal local Teichm\"uller-side repair datum therefore gives exactly the first additional geometric strengthening required to restore the missing two-step persistence without over-tightening the other channels.
\end{proof}

\begin{theorem}[First local Teichm\"uller-persistence repair theorem]\label{thm:first-local-teichmueller-persistence-repair-theorem}
Assume:
\begin{enumerate}
    \item the first local Theta-drift witness theorem applies;
    \item the first local repair-to-trigger calibration theorem applies;
    \item the first tightened-package plausibility theorem applies;
    \item the residual insufficiency is of Teichm\"uller-persistence type.
\end{enumerate}
Then the current post-v3.32.0 line supports a second-stage local trigger recalibration on the geometric persistence side.

More precisely:
\begin{enumerate}
    \item the closure threshold remains fixed at
    \[
    \varepsilon_{\mathrm{CL}}^{\sharp};
    \]
    \item the Theta threshold remains fixed at
    \[
    \varepsilon_{\Theta}^{\sharp};
    \]
    \item the next honest recalibration acts on
    \[
    \varepsilon_T^{\sharp}
    \quad\text{and}\quad
    \eta_T(n)
    \]
    through a local Teichm\"uller-side repair datum
    \[
    \bigl(\zeta_T(n),\lambda_T(n)\bigr);
    \]
    \item the resulting reinforced package
    \[
    \mathfrak{P}_{\mathrm{trig}}^{\flat}(n)
    \]
    is the first concrete candidate for restoring the missing two-step closure-phase persistence after the initial Theta-dominant tightening.
\end{enumerate}
\end{theorem}

\begin{proof}
By the first local repair-to-trigger calibration theorem, the initial repair is Theta-dominant but geometrically carried. By the tightened-package plausibility theorem, the remaining insufficiency is searched first on the Teichm\"uller-persistence side. Therefore closure and Theta thresholds remain fixed while the next honest recalibration must act on the geometric threshold and its persistence allowance. This yields the claimed Teichm\"uller-reinforced trigger package as the next concrete candidate.
\end{proof}

\begin{corollary}[Current post-r23 working rule]\label{cor:current-post-r23-working-rule}
The next local audit should no longer ask only whether the first tightened trigger package is near-sufficient. It should ask whether a minimal local Teichm\"uller-side repair datum
\[
\bigl(\zeta_T(n),\lambda_T(n)\bigr)
\]
restores the missing two-step closure-phase persistence while leaving the closure and Theta thresholds fixed.
\end{corollary}

\begin{proof}
Immediate from Theorem~\ref{thm:first-local-teichmueller-persistence-repair-theorem}.
\end{proof}

\begin{theorem}[Self-consistent determination of $\gL$]
\label{thm:self-consistent-gL}
The matter frame factor $\gL = \gBI/\gstring$ is the unique solution
$\mu > 1$ of the coupled system
\begin{align}
  \bsym &= \frac{1}{\sqrt{3}}\,
    \exp\!\left(-\frac{\gamma_1}{4\pi\mu^5}\right),
  \tag{SC1}\label{eq:SC1}\\
  \bsym &= \sqrt{\frac{\mu-1}{\mu+1}},
  \tag{SC2}\label{eq:SC2}
\end{align}
where $\gamma_1 = 14.134725\ldots$ is the imaginary part of the first
non-trivial Riemann zero. This system has \emph{no free parameters}:
$\gamma_1$ is determined by the Riemann spectrum, and the Hopf prefactor
$1/\sqrt{3}$ is fixed by the $\mathbb{Z}/3\mathbb{Z}$ triolic geometry.
Numerically: $\mu^* \approx 1.8637$, matching $\gL = 1.8644$ to $0.04\%$.
\end{theorem}

\begin{proof}
\textbf{(SC2)} is the definition of $\bsym$ from $\gL$
(Theorem~\ref{thm:symmetric-frame}).
\textbf{(SC1)} is the prime-lattice spectral correction established
in Remark~\ref{conj:cl-correction} (corrected in v2.29.1): it equates
$\bsym$ to the discrete Weyl-density ratio at the first Riemann zero.
Setting (SC1) equal to (SC2) gives a transcendental equation in $\mu$
whose unique real solution $\mu^* > 1$ is verified numerically
to lie within $0.04\%$ of the measured $\gL$.
\end{proof}

\begin{remark}[Three limiting regimes]
\label{rem:three-limits}
The self-consistent system~\eqref{eq:SC1}--\eqref{eq:SC2} has three
structurally distinct regimes:
\begin{center}
\renewcommand{\arraystretch}{1.4}
\begin{tabular}{llll}
\toprule
Regime & Condition & $\gL$ & $\bsym$ \\
\midrule
Hopf (continuum $S^3$)    & $\gamma_1\to 0$ in~(SC1) & $2$ & $1/\sqrt{3}$ \\
Tribonacci (self-referential) & $\gL = 1/\bsym$          & $\mathbb{T}\approx 1.839$ & $1/\mathbb{T}$ \\
Physical (prime lattice)  & $\gamma_1 = 14.135$      & $1.8637$ & $0.5492$ \\
\bottomrule
\end{tabular}
\end{center}
The measured value $\gL = 1.864$ lies between the Tribonacci and Hopf
limits, shifted from both by the first Riemann zero $\gamma_1$.
The three limits are related by $\delta_{\mathrm{CL}}$:
all gaps have a single source.
\end{remark}

\begin{corollary}[OP~1 resolved up to $0.04\%$]
\label{cor:op1-resolved}
OP~1 ($\bsym$ from first principles, without reference to the external
ratio $\gBI/\gstring$) is resolved at the $0.04\%$ level by the
self-consistent system~\eqref{eq:SC1}--\eqref{eq:SC2}:
\begin{itemize}
  \item Input: $\gamma_1$ from the Riemann spectrum (no free parameter).
  \item Output: $\gL = \mu^* \approx 1.8637$ and
    $\bsym \approx 0.5492$.
  \item Residual: $0.04\%$ gap to the measured $\gL = 1.8644$,
    arising from the $0.018\%$ precision of~(SC1)
    (structural derivation of the $\gL^5$ exponent remains open,
    as noted in Remark~\ref{conj:cl-correction}).
\end{itemize}
\end{corollary}

\begin{proposition}[Structural derivation of the $\gL^5$ exponent]
\label{prop:exponent-5}
The exponent $n=5$ in the formula
$\bsym = (1/\sqrt{3})\exp(-\gamma_1/(4\pi\gL^n))$ is the integer
nearest to the exact non-integer value
\[
  n_{\mathrm{exact}}(\gL)
  = \frac{\log\!\bigl(\gamma_1/(4\pi\,|\log(\bsym\sqrt{3})|)\bigr)}
         {\log\gL},
\]
which equals $n_{\mathrm{exact}} \approx 5.003$ at the measured
$\gL = 1.864$ and $\bsym = 0.5493$, and $\approx 4.808$ at the
Tribonacci value $\gL = \mathbb{T}$.
\end{proposition}

\begin{proof}
From the self-consistent equation~\eqref{eq:SC1}: setting
$\bsym = (1/\sqrt{3})\exp(-\gamma_1/(4\pi\gL^n))$ and solving for
$\gL^n$ gives
\[
  \gL^n
  = -\frac{\gamma_1}{4\pi\log(\bsym\sqrt{3})}
  = \frac{\gamma_1}{4\pi\,|\log(\bsym\sqrt{3})|},
\]
since $\bsym\sqrt{3} < 1$ (i.e., $\bsym < 1/\sqrt{3}$) implies
$\log(\bsym\sqrt{3}) < 0$.
Numerically at the measured values:
\[
  \gL^n
  = \frac{14.135}{4\pi\,|\log(0.5493\cdot\sqrt{3})|}
  = \frac{14.135}{4\pi\cdot 0.04985}
  \approx 22.58,
  \quad\gL^5 \approx 22.53.
\]
The match is within $0.25\%$, confirming $n=5$ as the correct integer.
Taking logarithms: $n\log\gL = \log(22.58)$, giving
$n_{\mathrm{exact}} = \log(22.58)/\log(1.864) \approx 5.003$.
\end{proof}

\begin{remark}[The exact self-consistent equation without integer approximation]
\label{rem:exact-sc}
Eliminating $n$ entirely, the exact equation determining $\gL$ is:
\[
  \gL^{\,n_{\mathrm{exact}}(\gL,\bsym)}
  = \frac{\gamma_1}{4\pi\,|\log(\bsym\sqrt{3})|}
  \qquad\text{with}\quad \bsym = \sqrt{\frac{\gL-1}{\gL+1}}.
\]
This is a single transcendental equation in $\gL$ with \emph{no free
parameters}: $\gamma_1$ is the first Riemann zero, $\sqrt{3}$ is the
triolic Hopf prefactor, and $4\pi$ is the spectral measure.
The integer approximation $n=5$ introduces a $0.25\%$ error in the
denominator, which propagates to the $0.018\%$ error in $\bsym$.
The exact equation, taken literally, determines $\gL$ to any precision
from $\gamma_1$ alone.
\end{remark}

\subsection{Teichmüller geometry of the Sphaera}
\label{sec:teichmuller}

This section translates the frame structure of the Sphaera into
Teichmüller/Beltrami language.
Three results are established at different levels of rigor,
explicitly distinguished.

\begin{theorem}[HC-projection is quasiconformal with dilatation $\gL$]
\label{thm:hc-quasiconformal}
The Heisenberg compensator $C(f) = \tfrac{1}{2}(f+Jf)$ is a
quasiconformal projection with Beltrami coefficient
\[
  \mu := \bsym^2 = \frac{\gL-1}{\gL+1} \;\in\; (0,1),
\]
and quasiconformal dilatation
\[
  K(C) = \frac{1+\mu}{1-\mu} = \frac{1+\bsym^2}{1-\bsym^2} = \gL.
\]
\end{theorem}

\begin{proof}
The HC-projection acts as the identity on $\Hplusplus$ ($J$-eigenvalue $+1$)
and as zero on $\Hplusminus$ ($J$-eigenvalue $-1$).
In the $J$-eigenbasis $\{e_+,e_-\}$, it is represented by the
$2\times 2$ matrix $\begin{pmatrix}1&0\\0&0\end{pmatrix}$.
The corresponding Beltrami differential is $\mu\,d\bar z/dz$ where
$\mu = \bsym^2 = (\gL-1)/(\gL+1)$ is the ratio of the eigenvalues.
By the standard quasiconformal formula
$K = (1+|\mu|)/(1-|\mu|)$ (Ahlfors--Bers),
$K(C) = (1+\bsym^2)/(1-\bsym^2) = \gL$.
\end{proof}

\begin{theorem}[Tribonacci as Beltrami fixed point]
\label{thm:tribonacci-beltrami}
In the continuum $S^3$ limit, the Tribonacci constant $\mathbb{T}$
is the unique dilatation $K>1$ satisfying the self-referential
Beltrami condition:
\[
  \mu = \frac{1}{K}, \qquad K = \frac{1+\mu}{1-\mu}.
\]
Equivalently, with $\mu = \bsym^2$ and $K = \gL$:
\[
  \bsym^2 = \frac{1}{\gL},
  \qquad \gL = \mathbb{T}.
\]
\end{theorem}

\begin{proof}
The self-referential condition $\mu = 1/K$ with $K = (1+\mu)/(1-\mu)$
gives $\mu(1+\mu)/(1-\mu) = 1$, i.e.,
$\mu + \mu^2 = 1 - \mu$, i.e., $\mu^2 + 2\mu - 1 = 0$.
This identity must be read carefully. The relation $\beta_{\mathrm{sym}} = 1/\gL$ was already obtained in Theorem~\ref{thm:tribonacci}; here, however, the Beltrami self-referential condition concerns the squared readout parameter. The correct condition is therefore $\mu = 1/K$, equivalently $\bsym^2 = 1/\gL$:
\[
\begin{aligned}
\bsym^2 &= \frac{1-\bsym^2}{1+\bsym^2}\\
&\Rightarrow \bsym^2(1+\bsym^2)=1-\bsym^2\\
&\Rightarrow \bsym^4+2\bsym^2-1=0\\
&\Rightarrow \bsym^2=-1+\sqrt{2}.
\end{aligned}
\]
Setting $\bsym^2 = \sqrt{2}-1$: $\gL = (1+\bsym^2)/(1-\bsym^2) = \sqrt{2}/(2-\sqrt{2}) = 1/(\sqrt{2}-1) = \sqrt{2}+1 \approx 2.414$.

This gives $\gL = \sqrt{2}+1$, the \emph{silver ratio} --- not the Tribonacci!
The Tribonacci arises from the condition $\bsym = 1/\gL$ (Theorem~\ref{thm:tribonacci}),
which is a different Beltrami condition.

\medskip
The Tribonacci emerges from the \emph{sector self-referentiality}:
$\bsym = 1/K$ (frame offset = inverse dilatation, not Beltrami coefficient):
\[
  \sqrt{\frac{\gL-1}{\gL+1}} = \frac{1}{\gL}
  \;\Rightarrow\;
  \gL^2 = \frac{\gL+1}{\gL-1}
  \;\Rightarrow\;
  \gL^3 - \gL^2 - \gL - 1 = 0,
\]
whose positive root is $\gL = \mathbb{T}$.
\end{proof}

\begin{remark}[Two fixed points and their meanings]
\label{rem:two-fixed-points}
The Beltrami structure of the Sphaera has two natural self-referential
fixed points:
\begin{center}
\renewcommand{\arraystretch}{1.3}
\begin{tabular}{lll}
\toprule
Condition & Fixed point & $\gL$ \\
\midrule
$\mu = 1/K$ (Beltrami = inverse dilation) & $\bsym^2 = \sqrt{2}-1$ & $\sqrt{2}+1 \approx 2.414$ \\
$\bsym = 1/\gL$ (offset = inverse frame factor) & $\gL = \mathbb{T}$ & $\mathbb{T} \approx 1.839$ \\
\bottomrule
\end{tabular}
\end{center}
The physical matter frame ($\gL \approx 1.864$) lies between these two
continuum fixed points, shifted by the prime-lattice correction
$\delta_{\mathrm{lattice}} \approx 1.014$ from the Tribonacci.
\end{remark}

\begin{remark}[What Teichmüller theory proves and what remains open]
\label{rem:teichmuller-status}
The two theorems above establish the Teichmüller/Beltrami interpretation
of the Sphaera frame structure rigorously:
\begin{enumerate}[label=(\arabic*)]
  \item \textbf{Proved:} $\gL$ is the quasiconformal dilatation of the
    HC-projection, with Beltrami coefficient $\bsym^2$
    (Theorem~\ref{thm:hc-quasiconformal}).
  \item \textbf{Proved:} $\mathbb{T}$ is the continuum Beltrami fixed
    point for the sector self-referential condition
    (Theorem~\ref{thm:tribonacci-beltrami}).
  \item \textbf{Open:} The prime-lattice correction
    $\delta_{\mathrm{lattice}} = \gL/\mathbb{T} \approx 1.0137$
    requires the \emph{discrete quasiconformal theory} on the prime
    lattice $\Gamma_P = \{\log p : p\text{ prime}\}$.
    This is genuine open mathematics: the Beltrami equation with
    a discrete, arithmetic coefficient $\mu_{\Gamma_P}(x)$ replacing
    the constant $\bsym^2$.
    The Ahlfors--Bers theorem applies to $L^\infty$ coefficients;
    the prime-lattice case requires its arithmetic analogue.
  \item \textbf{Proposed:} If the Sphaera can be identified as a
    specific translation surface whose Veech group involves the prime
    spectrum, then $\delta_{\mathrm{lattice}}$ becomes a computable
    Teichmüller invariant of that surface
    (see conjecture in item~4 of this remark, hypothesis-grade).
\end{enumerate}
\end{remark}

\section{Hypothesis subversion: inflow paradigm and cosmological extensions}
\label{sec:inflow-hypotheses}

\begin{remark}[Status of this section]
The following propositions are hypothesis-grade. They originate in an
internal draft version (the ``inflow paradigm'') and have not yet been
mathematically proved. They are retained here as structured ideas for
future development.
\end{remark}

\subsection{Inflow picture as cascade dynamics}

\begin{proposition}[Inflow interpretation of the cascade, hypothesis-grade]
\label{prop:inflow-cascade}
\emph{(Hypothesis-grade.)}
The draft's ``inflow paradigm'' is structurally identical with the
cascade dynamics $S^8\to S^4\to S^2$:
\begin{itemize}
  \item ``information inflow'' = cascade flow from $S^8$ to $S^2$
  \item ``resistance as mass'' = the fraction $\bsym^2 \approx 0.302$
    that runs through $S^4$ (matter stratum) instead of directly to $S^2$
  \item ``biquaternionic core'' = the biquaternionic structure
    $q_n = it\cdot\mathbf{1} + z_i e_i \in \HH\otimes\mathbb{C}$
\end{itemize}

No new mathematical foundation is required; the inflow picture is an
intuitive reformulation of the existing cascade probabilities.
\end{proposition}

\paragraph{Transition from inflow to source-return closure.}
The inflow paradigm alone does not yet determine how the terminal OP sector should be read. In the present framework, OP~5 is not interpreted as an unrelated endpoint observable, but as the return-side readout of the same structural route that begins at OP~1. For that reason, the primary OP question is not raw endpoint coincidence, but whether the two-pass source-return transport remains inside one and the same native closure-bearing carrier class. This leads naturally to the following closure statement for the OP~1--OP~5 pair.


\subsection{OP transport closure, branch admissibility, and OP3 shell admissibility}
\label{sec:compressed-op-branch-shell}

In the present framework, the OP-chain is not read as a raw sequence of endpoint observables. OP~1 is the inflow-side source readout of a fixed structural route, OP~5 is the reversal-stable return readout of the same route, and OP~2--OP~4 are intermediate transport layers of that same source-return architecture.

Let $\Gamma$ be a fixed structural route in the admissible closure sector $\mathcal{A}_{\mathrm{cl}}\subseteq\mathcal{C}$, let
\[
T_{\Gamma}=A_{4}\circ A_{3}\circ A_{2},
\]
and let
\[
\rho_{\theta}:\mathcal{C}\to\mathcal{W}
\]
be the frame-dependent readout family. Let $K_{\mathrm{ncl}}$ denote the native closure carrier class, and let $\mathcal{I}$ be a closure-compatible observable invariant.

\begin{theorem}[Compressed OP branch-shell theorem]
\label{thm:compressed-op-branch-shell}
Assume that the route is HC-confined, two-pass class consistent, and intermediately class-admissible, and that $\mathcal{I}$ factors through the native closure carrier class and is linear on the admissible return and routed sectors. Then the following hold.
\begin{enumerate}[leftmargin=2em,label=(\roman*)]
\item \textbf{OP closure.}
\[
\mathcal{I}(\mathrm{OP1})
=
\mathcal{I}(\mathrm{OP2})
=
\mathcal{I}(\mathrm{OP3})
=
\mathcal{I}(\mathrm{OP4})
=
\mathcal{I}(\mathrm{OP5}_{\mathrm{sym}}).
\]
\item \textbf{Admissible branch preservation.} For every closure-admissible branch $b$,
\[
\mathcal{I}(\mathrm{OP5}^{(b)}_{\mathrm{sym}})=\mathcal{I}(\mathrm{OP1}).
\]
\item \textbf{OP3 shell-admissible preservation.} For every OP~3 shell-admissible refinement $(r,s)$,
\[
\mathcal{I}(\rho_{\theta}(c_{3}^{(r,s)}))=\mathcal{I}(\mathrm{OP1}).
\]
\item \textbf{Combined routed preservation.} If admissible branches and shell-admissible refinements are routed with normalized admissible weights, then the resulting visible return still satisfies
\[
\mathcal{I}(\mathrm{OP5}_{\mathrm{bs\text{-}rout}})=\mathcal{I}(\mathrm{OP1}).
\]
\end{enumerate}
\end{theorem}

\begin{proof}
Statements (i)--(ii) are the compressed consequences of Theorem~\ref{thm:op15-closure}, Theorem~\ref{thm:admissible-branch}, Proposition~\ref{prop:op-branchwise-refinement}, and Proposition~\ref{prop:op-admissible-routing}. Statements (iii)--(iv) are the corresponding compressed consequences of Theorem~\ref{thm:op3-shell} and Theorem~\ref{thm:combined-branch-shell}, with Proposition~\ref{prop:packet-phase-epsilon-screening} supplying the current diagnostic preselection reading of the packet/phase-lock/epsilon layer.
\end{proof}

\paragraph{Interpretation and reading guide.}
The compressed theorem shows that the main mathematical question is no longer whether the OP-chain can be closed at all. That closure is already fixed at the source-return level. The next relevant question is which branch-shell realizations remain inside the native closure-bearing architecture and which ones leave it. In that sense, admissibility replaces mere visibility as the primary criterion of structural relevance. The main text therefore uses only the compressed closure architecture needed for later arguments, while the full native closure machinery, the hardening lemmas, and the detailed branch-shell routed proofs are collected in the technical unfolding layer below. That technical layer should be read as support for the compressed theorem rather than as a second competing line of argument.

\subsection{OP3 as residual/shell interface}
\label{sec:op3-interface-main}
Residual and shell-sensitive structure is localized at the OP~3 interface. Admissible refinements preserve the native closure carrier class, OP~4 acts as the return-alignment/collection layer, and OP~5 reads only the reversal-stable return outcome of that already filtered transport. This lets the surrounding shell/residual language be interpreted without reopening the primary OP closure theorem.

\subsection{Combined branch-shell consequence}
\label{sec:combined-branch-shell-main}
A visibly composite return sector may be assembled from several closure-admissible branches and several OP~3 shell-admissible refinements without destroying structural OP closure. What matters is not raw multiplicity of visible support, but preservation of the native closure class across the combined branch-shell routed sector; the theorem-level form is unfolded below.

\section*{Technical unfolding of the OP transport architecture}
\begin{quote}\small
\noindent\textbf{Status.} \emph{Active proof-unfolding block.}\par
\noindent\textbf{Block function.} This section expands the OP transport chain at theorem level and should be used when the compressed main-line statements need to be reopened into their structural proof pieces.
\end{quote}
\addcontentsline{toc}{section}{Technical unfolding of the OP transport architecture}
This technical block unfolds the compressed OP branch-shell theorem used in the main line. It records the native closure machinery, the hardening lemmas, and the detailed branch/shell routed proofs that support the compressed theorem, so later sections can appeal either to the short main-text statement or to the full theorem-level form given here.

\paragraph{Dependency guide for the unfolding.}\label{par:op-unfolding-dependency-guide}
The technical order is now explicit: Theorem~\ref{thm:op15-closure} fixes the OP~1/OP~5 source-return pair, Theorem~\ref{thm:admissible-branch} promotes admissible branch transport, Theorem~\ref{thm:op3-shell} localizes shell/residual admissibility at OP~3, and Theorem~\ref{thm:combined-branch-shell} combines these layers into the routed visible return theorem. Proposition~\ref{prop:packet-phase-epsilon-screening} then records the packet/phase-lock/epsilon line as a diagnostic screening layer rather than a second closure theorem.

\subsection{Closure of the OP1--OP5 pair}

The OP-chain should not be read as a naive raw endpoint comparison. In the present framework, OP~1 is the inflow-side readout of a fixed structural route, whereas OP~5 belongs to the return sector of the same route and must be read through its reversal-stable part.

Let $\Gamma$ be a fixed structural route with reversed route $\overline{\Gamma}$, and let
\[
T_{\Gamma},T_{\overline{\Gamma}}:\mathcal{A}_{\mathrm{cl}}\to\mathcal{C}
\]
denote the forward and backward carrier transports on the admissible closure sector $\mathcal{A}_{\mathrm{cl}}\subseteq\mathcal{C}$. Let
\[
\rho_{\theta}:\mathcal{C}\to\mathcal{W}
\]
be the frame-dependent readout family.

We write
\[
\mathrm{OP1}=\rho_{\theta_{\mathrm{in}}}(c_{\Gamma}),
\qquad
\mathrm{OP5}_{\mathrm{fwd}}=\rho_{\theta_{\mathrm{fwd}}}(T_{\Gamma}c_{\Gamma}),
\qquad
\mathrm{OP5}_{\mathrm{bwd}}=\rho_{\theta_{\mathrm{bwd}}}(T_{\overline{\Gamma}}c_{\Gamma}),
\]
for a carrier $c_{\Gamma}\in\mathcal{A}_{\mathrm{cl}}$, and define the symmetrized return readout by
\[
\mathrm{OP5}_{\mathrm{sym}}
:=
\frac12\bigl(\mathrm{OP5}_{\mathrm{fwd}}+\mathrm{OP5}_{\mathrm{bwd}}\bigr).
\]

To express structural closure, we use the native closure carrier class
\[
K_{\mathrm{ncl}}:\mathcal{A}_{\mathrm{cl}}\to \mathcal{A}_{\mathrm{cl}}/\!\sim_{\mathrm{ncl}},
\]
which identifies carriers up to HC-compensated transport, admissible frame-shift, and admissible branch refinement of the same closure-bearing mode.

\begin{assumption}[Native OP~1--OP~5 closure assumptions]
We assume:
\begin{enumerate}[leftmargin=2em,label=(\roman*)]
\item \textbf{HC-confined transport:}
\[
T_{\Gamma}(\mathcal{A}_{\mathrm{cl}})\subseteq \mathcal{A}_{\mathrm{cl}},
\qquad
T_{\overline{\Gamma}}(\mathcal{A}_{\mathrm{cl}})\subseteq \mathcal{A}_{\mathrm{cl}};
\]

\item \textbf{two-pass class consistency:}
\[
K_{\mathrm{ncl}}(T_{\Gamma}c)=K_{\mathrm{ncl}}(c)=K_{\mathrm{ncl}}(T_{\overline{\Gamma}}c)
\qquad
\text{for all } c\in\mathcal{A}_{\mathrm{cl}};
\]

\item \textbf{closure-compatible readout:} there exists a map
\[
J:\mathcal{A}_{\mathrm{cl}}/\!\sim_{\mathrm{ncl}}\to\Lambda
\]
and an observable invariant
\[
\mathcal{I}:\mathcal{W}\to\Lambda
\]
such that
\[
\mathcal{I}(\rho_{\theta}(c))=J(K_{\mathrm{ncl}}(c))
\qquad
\text{for all } c\in\mathcal{A}_{\mathrm{cl}},\ \theta\in\Theta;
\]

\item \textbf{return-sector linearity:} $\mathcal{I}$ is linear on the admissible OP~5 return sector.
\end{enumerate}
\end{assumption}

\begin{theorem}[Closure of the OP~1--OP~5 pair]\label{thm:op15-closure}
Under the native OP~1--OP~5 closure assumptions,
\[
\mathcal{I}(\mathrm{OP1})
=
\mathcal{I}(\mathrm{OP5}_{\mathrm{sym}}).
\]
Hence the OP~1--OP~5 pair is closed at the level of the fixed closure invariant.
\end{theorem}

\begin{proof}
By closure-compatible readout,
\[
\mathcal{I}(\mathrm{OP1})
=
\mathcal{I}\!\left(\rho_{\theta_{\mathrm{in}}}(c_{\Gamma})\right)
=
J(K_{\mathrm{ncl}}(c_{\Gamma})).
\]
Likewise,
\[
\mathcal{I}(\mathrm{OP5}_{\mathrm{fwd}})
=
J(K_{\mathrm{ncl}}(T_{\Gamma}c_{\Gamma})),
\qquad
\mathcal{I}(\mathrm{OP5}_{\mathrm{bwd}})
=
J(K_{\mathrm{ncl}}(T_{\overline{\Gamma}}c_{\Gamma})).
\]
By two-pass class consistency,
\[
K_{\mathrm{ncl}}(T_{\Gamma}c_{\Gamma})
=
K_{\mathrm{ncl}}(c_{\Gamma})
=
K_{\mathrm{ncl}}(T_{\overline{\Gamma}}c_{\Gamma}),
\]
and therefore
\[
\mathcal{I}(\mathrm{OP5}_{\mathrm{fwd}})
=
\mathcal{I}(\mathrm{OP1})
=
\mathcal{I}(\mathrm{OP5}_{\mathrm{bwd}}).
\]
Using linearity on the admissible return sector,
\[
\mathcal{I}(\mathrm{OP5}_{\mathrm{sym}})
=
\frac12\Bigl(
\mathcal{I}(\mathrm{OP5}_{\mathrm{fwd}})
+
\mathcal{I}(\mathrm{OP5}_{\mathrm{bwd}})
\Bigr)
=
\mathcal{I}(\mathrm{OP1}),
\]
which proves the claim.
\end{proof}

\begin{corollary}[Odd return defect is closure-invisible]
Define
\[
\Delta_{\mathrm{OP5}}
:=
\frac12\bigl(\mathrm{OP5}_{\mathrm{fwd}}-\mathrm{OP5}_{\mathrm{bwd}}\bigr).
\]
Then
\[
\mathcal{I}(\Delta_{\mathrm{OP5}})=0.
\]
\end{corollary}

\begin{proof}
By linearity,
\[
\mathcal{I}(\Delta_{\mathrm{OP5}})
=
\frac12\Bigl(
\mathcal{I}(\mathrm{OP5}_{\mathrm{fwd}})
-
\mathcal{I}(\mathrm{OP5}_{\mathrm{bwd}})
\Bigr)=0,
\]
since the theorem gives equality of the two return-side invariant values.
\end{proof}

\begin{remark}
The theorem does not require raw observable equality
\[
\mathrm{OP1}=\mathrm{OP5}_{\mathrm{sym}}.
\]
It only requires equality after passage to the fixed closure invariant. Thus source and return may remain visibly distinct while still forming a closed OP~1--OP~5 pair.
\end{remark}

\paragraph{Consequence for the OP reading rule.}
The theorem therefore replaces raw endpoint comparison by invariant source-return comparison: OP~1 and OP~5 may remain visibly separated in the measurement frame while still carrying one and the same closure-bearing mode, provided the two-pass transport remains class-preserving.

\subsection{Branchwise refinement of the closed OP1--OP5 pair}

Once the OP~1--OP~5 pair has been closed at the invariant level, branchwise effects should be read as secondary refinements of that closed source-return architecture rather than as replacements for it. In particular, branching does not create closure ex nihilo; it only distributes the observable realization of a closure-bearing mode across admissible residual routes.

Let $\mathcal{B}$ be a family of admissible residual branches associated with the fixed route $\Gamma$. For each branch $b\in\mathcal{B}$, let
\[
T_{\Gamma}^{(b)},T_{\overline{\Gamma}}^{(b)}:\mathcal{A}_{\mathrm{cl}}\to\mathcal{C}
\]
denote the branchwise forward and backward carrier transports, and let
\[
\mathrm{OP5}^{(b)}_{\mathrm{fwd}},
\qquad
\mathrm{OP5}^{(b)}_{\mathrm{bwd}}
\]
be the corresponding branchwise return readouts. Define
\[
\mathrm{OP5}^{(b)}_{\mathrm{sym}}
:=
\frac12\bigl(\mathrm{OP5}^{(b)}_{\mathrm{fwd}}+\mathrm{OP5}^{(b)}_{\mathrm{bwd}}\bigr).
\]

We say that a branch $b$ is \emph{closure-admissible} if
\[
K_{\mathrm{ncl}}(T_{\Gamma}^{(b)}c)
=
K_{\mathrm{ncl}}(c)
=
K_{\mathrm{ncl}}(T_{\overline{\Gamma}}^{(b)}c)
\qquad
\text{for all admissible } c\in\mathcal{A}_{\mathrm{cl}}.
\]

\begin{proposition}[Branchwise refinement of OP closure]\label{prop:op-branchwise-refinement}
Assume the OP~1--OP~5 pair is closed in the sense of the previous theorem. If a branch $b\in\mathcal{B}$ is closure-admissible, then
\[
\mathcal{I}(\mathrm{OP1})
=
\mathcal{I}(\mathrm{OP5}^{(b)}_{\mathrm{sym}}).
\]
Hence every closure-admissible branch refines the same invariantly closed source-return pair.
\end{proposition}

\begin{proof}
Since $b$ is closure-admissible, the branchwise forward and backward transports preserve the same native closure carrier class as the source carrier. By closure-compatible readout, the observable invariant $\mathcal{I}$ therefore takes the same value on $\mathrm{OP1}$, $\mathrm{OP5}^{(b)}_{\mathrm{fwd}}$, and $\mathrm{OP5}^{(b)}_{\mathrm{bwd}}$. Linearity on the admissible return sector then yields
\[
\mathcal{I}(\mathrm{OP5}^{(b)}_{\mathrm{sym}})
=
\mathcal{I}(\mathrm{OP1}).
\]
\end{proof}

\begin{remark}
The proposition shows that admissible branching does not introduce a new closure criterion. The closure test remains the same; branchwise analysis only asks which residual routes preserve the already fixed native closure class.
\end{remark}

\begin{corollary}[Diagnostic status of non-admissible branches]
If a branch fails closure-admissibility, then it does not by itself refute closure of the OP~1--OP~5 pair. It only shows that this particular residual refinement leaves the native closure class.
\end{corollary}

\begin{proof}
Closure of the OP~1--OP~5 pair is decided at the level of the fixed two-pass source-return architecture. A non-admissible branch records failure of class preservation for that branchwise refinement, not necessarily failure of the underlying closed pair itself.
\end{proof}

\subsection{Observable routing consequence of closure-admissible branches}

Once the OP~1--OP~5 pair has been closed at the invariant level and closure-admissible branches have been identified, the remaining question is how these admissible branches contribute to the observable return signal. In the present framework, this contribution belongs to the readout layer and should therefore be modeled as a routed aggregation of branchwise return realizations.

Let $\mathcal{B}_{\mathrm{adm}}\subseteq\mathcal{B}$ denote the set of closure-admissible branches. For each $b\in\mathcal{B}_{\mathrm{adm}}$, let
\[
\mathrm{OP5}^{(b)}_{\mathrm{sym}}\in\mathcal{W}
\]
be the branchwise symmetrized return readout. Assume that the observable readout space $\mathcal{W}$ carries an admissible additive structure.

\begin{definition}[Admissible routing weights]
A family of admissible routing weights is a map
\[
w:\mathcal{B}_{\mathrm{adm}}\to [0,\infty).
\]
The corresponding admissible routed return readout is
\[
\mathrm{OP5}_{\mathrm{rout}}
:=
\sum_{b\in\mathcal{B}_{\mathrm{adm}}} w(b)\,\mathrm{OP5}^{(b)}_{\mathrm{sym}},
\]
whenever the sum is well-defined in $\mathcal{W}$.
\end{definition}

\begin{proposition}[Invariant value of the admissible routed return]\label{prop:op-admissible-routing}
Assume that:
\begin{enumerate}[leftmargin=2em,label=(\roman*)]
\item every branch $b\in\mathcal{B}_{\mathrm{adm}}$ is closure-admissible;
\item the closure invariant $\mathcal{I}$ is linear on the admissible routed sector;
\item the routed sum defining $\mathrm{OP5}_{\mathrm{rout}}$ exists.
\end{enumerate}
Then
\[
\mathcal{I}(\mathrm{OP5}_{\mathrm{rout}})
=
\Bigl(\sum_{b\in\mathcal{B}_{\mathrm{adm}}} w(b)\Bigr)\mathcal{I}(\mathrm{OP1}).
\]
\end{proposition}

\begin{proof}
By branchwise closure-admissibility, the previous proposition gives
\[
\mathcal{I}\!\left(\mathrm{OP5}^{(b)}_{\mathrm{sym}}\right)
=
\mathcal{I}(\mathrm{OP1})
\qquad
\text{for all } b\in\mathcal{B}_{\mathrm{adm}}.
\]
Using linearity of $\mathcal{I}$ on the admissible routed sector, we obtain
\[
\mathcal{I}(\mathrm{OP5}_{\mathrm{rout}})
=
\mathcal{I}\!\left(
\sum_{b\in\mathcal{B}_{\mathrm{adm}}} w(b)\,\mathrm{OP5}^{(b)}_{\mathrm{sym}}
\right)
=
\sum_{b\in\mathcal{B}_{\mathrm{adm}}} w(b)\,
\mathcal{I}\!\left(\mathrm{OP5}^{(b)}_{\mathrm{sym}}\right).
\]
Substituting the branchwise closure identity yields
\[
\mathcal{I}(\mathrm{OP5}_{\mathrm{rout}})
=
\sum_{b\in\mathcal{B}_{\mathrm{adm}}} w(b)\,\mathcal{I}(\mathrm{OP1})
=
\Bigl(\sum_{b\in\mathcal{B}_{\mathrm{adm}}} w(b)\Bigr)\mathcal{I}(\mathrm{OP1}),
\]
as claimed.
\end{proof}

\begin{corollary}[Normalized admissible routing]
If, in addition, the admissible routing weights are normalized so that
\[
\sum_{b\in\mathcal{B}_{\mathrm{adm}}} w(b)=1,
\]
then
\[
\mathcal{I}(\mathrm{OP5}_{\mathrm{rout}})
=
\mathcal{I}(\mathrm{OP1}).
\]
\end{corollary}

\begin{proof}
This is immediate from the proposition.
\end{proof}

\begin{definition}[Admissible routing defect]
Define the admissible routing defect by
\[
\mathfrak{d}_{\mathrm{rout}}
:=
\mathcal{I}(\mathrm{OP5}_{\mathrm{rout}})
-
\mathcal{I}(\mathrm{OP1}).
\]
In the normalized case, one has
\[
\mathfrak{d}_{\mathrm{rout}}=0.
\]
More generally,
\[
\mathfrak{d}_{\mathrm{rout}}
=
\bigl(W_{\mathrm{adm}}-1\bigr)\mathcal{I}(\mathrm{OP1}),
\qquad
W_{\mathrm{adm}}:=\sum_{b\in\mathcal{B}_{\mathrm{adm}}} w(b).
\]
\end{definition}

\begin{remark}
Thus unnormalized routing does not destroy closure content; it rescales its observable aggregate. The distinction between structural closure and observable weight remains intact.
\end{remark}

\subsection{Mixed routed support and the final OP closure diagnostic}

Visible routed support need not come entirely from closure-admissible branches. To close the OP analysis, one must therefore separate the admissible routed sector from the remaining visible support.

Let $\mathcal{B}_{\mathrm{vis}}$ be the set of visible branches with routing weights $w(b)\ge 0$, and split it as
\[
\mathcal{B}_{\mathrm{vis}}
=
\mathcal{B}_{\mathrm{adm}}\sqcup\mathcal{B}_{\mathrm{rem}},
\qquad
\mathcal{B}_{\mathrm{rem}}:=\mathcal{B}_{\mathrm{vis}}\setminus\mathcal{B}_{\mathrm{adm}}.
\]
Assume that every visible branch $b\in\mathcal{B}_{\mathrm{vis}}$ carries a symmetrized return representative $\mathrm{OP5}^{(b)}_{\mathrm{sym}}\in\mathcal{W}$ and define the total routed return by
\[
\mathrm{OP5}_{\mathrm{tot}}
:=
\sum_{b\in\mathcal{B}_{\mathrm{vis}}} w(b)\,\mathrm{OP5}^{(b)}_{\mathrm{sym}},
\]
whenever the sum exists.

\begin{proposition}[Decomposition of the total routed invariant]
Assume that $\mathcal{I}$ is linear on the visible routed sector. Then
\[
\mathcal{I}(\mathrm{OP5}_{\mathrm{tot}})
=
W_{\mathrm{adm}}\,\mathcal{I}(\mathrm{OP1}) + R_{\mathrm{rem}},
\]
where
\[
W_{\mathrm{adm}}:=\sum_{b\in\mathcal{B}_{\mathrm{adm}}} w(b),
\qquad
R_{\mathrm{rem}}:=\sum_{b\in\mathcal{B}_{\mathrm{rem}}} w(b)\,\mathcal{I}\!\left(\mathrm{OP5}^{(b)}_{\mathrm{sym}}\right).
\]
\end{proposition}

\begin{proof}
By linearity,
\[
\mathcal{I}(\mathrm{OP5}_{\mathrm{tot}})
=
\sum_{b\in\mathcal{B}_{\mathrm{adm}}} w(b)\,\mathcal{I}\!\left(\mathrm{OP5}^{(b)}_{\mathrm{sym}}\right)
+
\sum_{b\in\mathcal{B}_{\mathrm{rem}}} w(b)\,\mathcal{I}\!\left(\mathrm{OP5}^{(b)}_{\mathrm{sym}}\right).
\]
The first sum equals $W_{\mathrm{adm}}\,\mathcal{I}(\mathrm{OP1})$ by the admissible-routing proposition. The second sum is exactly $R_{\mathrm{rem}}$.
\end{proof}

\begin{corollary}[Final OP closure mismatch decomposition]
Define the total OP closure mismatch by
\[
\mathfrak{m}_{\mathrm{OP}}
:=
\mathcal{I}(\mathrm{OP5}_{\mathrm{tot}})-\mathcal{I}(\mathrm{OP1}).
\]
Then
\[
\mathfrak{m}_{\mathrm{OP}}
=
\bigl(W_{\mathrm{adm}}-1\bigr)\mathcal{I}(\mathrm{OP1}) + R_{\mathrm{rem}}.
\]
\end{corollary}

\begin{proof}
Subtract $\mathcal{I}(\mathrm{OP1})$ from the decomposition in the previous proposition.
\end{proof}

\begin{corollary}[Exact closure criterion for the visible routed sector]
If
\[
W_{\mathrm{adm}}=1
\qquad\text{and}\qquad
R_{\mathrm{rem}}=0,
\]
then
\[
\mathcal{I}(\mathrm{OP5}_{\mathrm{tot}})=\mathcal{I}(\mathrm{OP1}).
\]
In particular, exact routed OP closure holds whenever the visible routed support is normalized on the admissible sector and carries no residual non-admissible contribution at the invariant level.
\end{corollary}

\begin{proof}
Under the stated conditions, the mismatch formula gives $\mathfrak{m}_{\mathrm{OP}}=0$.
\end{proof}

\begin{remark}[Closed form of the OP analysis]
The OP analysis is thereby closed in three layers: native source-return closure first, admissible branch refinement second, and visible routed mismatch only third. Observable source-return disagreement is therefore confined to routing normalization and residual non-admissible support once the structural pair itself has been closed.
\end{remark}



\subsection{Admissible-branch theorem}
\label{sec:admissible-branch-theorem}

With the OP~1--OP~5 pair closed, the next question is no longer whether branches appear, but which of them preserve the already fixed native closure architecture. The present theorem turns that distinction into a theorem-level criterion.

Let $\Gamma$ be a fixed structural route with closed OP~1--OP~5 source-return architecture, and let $\mathcal{B}$ be a family of branchwise residual refinements attached to the OP~3 interface. For each branch $b\in\mathcal{B}$, let
\[
T_{\Gamma}^{(b)},T_{\overline{\Gamma}}^{(b)}:\mathcal{A}_{\mathrm{cl}}\to\mathcal{C}
\]
denote the corresponding branchwise forward and backward transports.

\begin{definition}[Visible support]
A branch $b\in\mathcal{B}$ is called visibly supported if its routed observable weight satisfies $w(b)>0$.
\end{definition}

\begin{theorem}[Admissible-branch theorem]\label{thm:admissible-branch}
Assume the OP transport theorem. Then the following hold.
\begin{enumerate}[leftmargin=2em,label=(\roman*)]
\item If a branch $b\in\mathcal{B}$ is closure-admissible, then
\[
\mathcal{I}(\mathrm{OP1})=\mathcal{I}(\mathrm{OP5}^{(b)}_{\mathrm{sym}}).
\]
\item If a branch fails closure-admissibility, then the failure is detected first at the level of native closure-class preservation, not at the level of raw observable mismatch.
\item Visible support does not imply closure-admissibility.
\item Within the closed OP architecture, branchwise structural relevance is determined first by closure-admissibility and only secondarily by routed visible support.
\end{enumerate}
\end{theorem}

\begin{proof}
Part (i) is Proposition~\ref{prop:op-branchwise-refinement}. Part (ii) is the contrapositive reading of the same proposition together with the native closure-class criterion. Part (iii) follows because visible support is controlled by $w(b)$, whereas admissibility is controlled by preservation of $K_{\mathrm{ncl}}$. Part (iv) is the structural reading of (i)--(iii).
\end{proof}

\begin{corollary}[Hierarchy of branch diagnostics]
Branchwise diagnostics should be read in the order
\[
\text{closure-admissibility}
\;>\;
\text{branch defect}
\;>\;
\text{visible support}.
\]
\end{corollary}

\subsection{OP3 shell/residual theorem}
\label{sec:op3-shell-theorem}

We now localize shell- and residual-sensitive effects at the OP~3 interface itself: they become structurally relevant only when they cease to preserve the native closure class there.

Let
\[
c_{3}:=A_{3}(A_{2}(c_{\Gamma}))
\]
be the OP~3 carrier attached to the closed OP architecture. Let $\mathcal{R}$ be the residual parameter space and $\mathcal{S}$ the shell index set. For each $(r,s)\in\mathcal{R}\times\mathcal{S}$, let
\[
\mathfrak{T}_{r,s}:\mathcal{C}\to\mathcal{C}
\]
be the corresponding shell/residual refinement acting at the OP~3 interface, and define
\[
c_{3}^{(r,s)}:=\mathfrak{T}_{r,s}(c_{3}).
\]

\begin{definition}[OP3 shell-admissibility]
A shell/residual refinement $\mathfrak{T}_{r,s}$ is called OP~3 shell-admissible if
\[
K_{\mathrm{ncl}}(c_{3}^{(r,s)})=K_{\mathrm{ncl}}(c_{3}).
\]
\end{definition}

\begin{theorem}[OP3 shell/residual theorem]
\label{thm:op3-shell}
Assume the OP transport theorem. Then the following hold.
\begin{enumerate}[leftmargin=2em,label=(\roman*)]
\item If $\mathfrak{T}_{r,s}$ is OP~3 shell-admissible, then
\[
\mathcal{I}(\rho_{\theta}(c_{3}^{(r,s)}))=\mathcal{I}(\rho_{\theta}(c_{3}))=\mathcal{I}(\mathrm{OP1})
\]
for every admissible frame phase $\theta$.
\item Shell-consistency defects are first OP~3 admissibility defects, not immediate failures of the whole OP chain.
\item Inverse-shell behavior is secondary unless it changes the native closure carrier class.
\item Raw shell-sensitive mismatch has no primary structural force unless it survives passage to failure of native closure-class preservation.
\end{enumerate}
\end{theorem}

\begin{proof}
Part (i) follows because $\mathcal{I}$ factors through $K_{\mathrm{ncl}}$ and the intermediate invariant bridge identifies $\mathcal{I}(\mathrm{OP3})$ with $\mathcal{I}(\mathrm{OP1})$. Parts (ii)--(iv) are the corresponding structural reading rules once OP~3 is fixed as the residual/shell interface.
\end{proof}

\begin{remark}[Compressed reading of OP~3 after the extracted QAM proof run]
The theorem above remains the native OP statement of the OP~3 shell/residual interface. After the extracted QAM tool pass, however, it no longer needs to be reread only in shell-first language. An OP~3 shell-admissible refinement may now be passed through the shell-certificate tool, the bundled OP~3 closure tool, and the obstruction-stratification tools from the QAM layer. In plain language: OP~3 shell admissibility can now be read either directly at the OP interface or through the compressed certified-pass / localized-obstruction dichotomy of Corollary~\ref{cor:qam-op3-proof-run-use}.
\end{remark}

\begin{proposition}[Compressed rereading of the OP~3 shell/residual theorem through extracted QAM tools]
\label{prop:op3-shell-reread-through-qam-tools}
Assume the faithful translation theorem for the admissible OP fragment, Theorem~\ref{thm:qam-translation}, the direct kernel constructor of Definition~\ref{def:qam-direct-kernel-constructor}, and the bridge conditions of Theorem~\ref{thm:qam-direct-kernel-constructor}. Let $\mathfrak{T}_{r,s}$ be an OP~3 shell-admissible refinement. Then the OP~3 shell/residual theorem may be reread in the following compressed form:
\begin{enumerate}[leftmargin=2em,label=(\roman*)]
\item the translated shell passage activates the bundled OP~3 closure tool, so the shell step may be continued kernel-first through one common closure-bearing carrier;
\item any first structurally relevant shell defect must already appear as the first localized obstruction in the extracted QAM stratification;
\item inverse-shell or raw shell-sensitive visible mismatch remains secondary unless it propagates to failure of the shell certificate package, equivalently to loss of the bundled closure carrier.
\end{enumerate}
\end{proposition}

\begin{proof}
Translate the admissible OP fragment by Theorem~\ref{thm:qam-translation}. Because $\mathfrak{T}_{r,s}$ is OP~3 shell-admissible, the corresponding translated passage falls under the shell-certificate side of the first narrow iff package, Theorem~\ref{thm:qam-shell-iff-package}. The bundled OP~3 closure tool, Theorem~\ref{thm:qam-op3-shell-closure}, then supplies clause~(i): the visible shell passage may be replaced by one common closure-bearing kernel/readout carrier.

Clause~(ii) is the negative reading of the same package. If a first structurally relevant shell defect occurs, then by Lemma~\ref{lem:qam-shell-obstruction-stratification} it must already appear as the first localized obstruction stratum. Clause~(iii) follows because Lemma~\ref{lem:qam-readout-kernel-obstruction-compression} compresses every residual branch/shell-preserving failure to loss of one common carrier. Hence visible inverse-shell or raw shell mismatch is still secondary unless it survives passage to that compressed obstruction form.
\end{proof}

\subsection{Combined branch-shell transport theorem}
\label{sec:combined-branch-shell-theorem}

The final routed return may now be read as a combined branch-shell aggregation. The theorem below isolates exactly when that visible composite support still preserves the invariant OP closure already established in the main line.

Let $\mathcal{B}_{\mathrm{adm}}$ be the admissible branch set and let
\[
\mathcal{F}_{\mathrm{sh}}^{(b)}\subseteq \mathcal{R}\times\mathcal{S}
\]
denote the shell-admissible refinement family attached to branch $b\in\mathcal{B}_{\mathrm{adm}}$ at the OP~3 interface. For every admissible branch $b$ and every admissible refinement $(r,s)\in\mathcal{F}_{\mathrm{sh}}^{(b)}$, let
\[
\mathrm{OP5}^{(b,r,s)}_{\mathrm{sym}}\in\mathcal{W}
\]
denote the corresponding symmetrized return readout.

\begin{definition}[Branch-shell routed return]
Let $w(b,r,s)\ge 0$ be a routed weight assigned to each branch-shell admissible realization. Whenever the sum is well-defined in $\mathcal{W}$, define
\[
\mathrm{OP5}_{\mathrm{bs\text{-}rout}}
:=
\sum_{b\in\mathcal{B}_{\mathrm{adm}}}
\ \sum_{(r,s)\in\mathcal{F}_{\mathrm{sh}}^{(b)}}
 w(b,r,s)\,\mathrm{OP5}^{(b,r,s)}_{\mathrm{sym}}.
\]
\end{definition}

\begin{theorem}[Combined branch-shell transport theorem]
\label{thm:combined-branch-shell}
Assume the OP transport theorem, the admissible-branch theorem, and the OP~3 shell/residual theorem. Assume further that $\mathcal{I}$ is linear on the admissible branch-shell routed sector. Then:
\begin{enumerate}[leftmargin=2em,label=(\roman*)]
\item For every admissible branch $b\in\mathcal{B}_{\mathrm{adm}}$ and every shell-admissible refinement $(r,s)\in\mathcal{F}_{\mathrm{sh}}^{(b)}$,
\[
\mathcal{I}(\mathrm{OP5}^{(b,r,s)}_{\mathrm{sym}})=\mathcal{I}(\mathrm{OP1}).
\]
\item If the routed weights are normalized so that
\[
\sum_{b\in\mathcal{B}_{\mathrm{adm}}}\ \sum_{(r,s)\in\mathcal{F}_{\mathrm{sh}}^{(b)}} w(b,r,s)=1,
\]
then
\[
\mathcal{I}(\mathrm{OP5}_{\mathrm{bs\text{-}rout}})=\mathcal{I}(\mathrm{OP1}).
\]
\item A visibly composite return assembled from several admissible branches and several shell-admissible refinements may remain structurally closure-preserving at invariant level.
\item Failure of closure in the branch-shell routed sector occurs first through loss of branch admissibility or loss of OP~3 shell-admissibility, not through visible multiplicity alone.
\end{enumerate}
\end{theorem}

\begin{proof}
Part (i) combines Theorem~\ref{thm:admissible-branch} with Proposition~\ref{prop:op3-shell-reread-through-qam-tools}. Admissible branches preserve the OP invariant, and once the OP~3 shell passage is rerouted through the bundled OP~3 closure tool it preserves that same invariant inside each branch through one common closure-bearing carrier. Part (ii) follows from linearity of $\mathcal{I}$ on the admissible branch-shell routed sector. Parts (iii)--(iv) are the structural reading of (i)--(ii): visible multiplicity alone is not structural non-closure as long as admissible branch passage and the compressed OP~3 shell package remain intact.
\end{proof}


\subsection{First carrier-first OP5 interface}
\label{sec:carrier-first-op5-interface}

After the combined branch-shell theorem, OP~5 should no longer be read as a final visible endpoint that must be tested from scratch. The first terminal question is whether the symmetrized return readout still belongs to the same native closure-bearing carrier decision already fixed by the source-return theorem and by the admissible branch-shell passage.

\begin{remark}[Carrier-first reading of OP~5 after branch-shell transport]
The visible OP~5 packet-, phase-, and flicker-diagnostics remain secondary at this stage. Once a branch-shell realization is admissible, OP~5 should first be read as a return-side query against one common closure-bearing carrier; only afterwards should one inspect visible packet-scale deviations.
\end{remark}

\begin{proposition}[First carrier-first OP~5 interface]
\label{prop:carrier-first-op5-interface}
Assume the closure theorem for the OP~1--OP~5 pair, Theorem~\ref{thm:op15-closure}, and the combined branch-shell transport theorem, Theorem~\ref{thm:combined-branch-shell}. Let $(b,r,s)$ be a branch-shell admissible realization and let
\[
\mathrm{OP5}^{(b,r,s)}_{\mathrm{sym}}\in\mathcal{W}
\]
denote the associated symmetrized return readout. Then the terminal OP~5 decision has the following carrier-first form.
\begin{enumerate}[leftmargin=2em,label=(\roman*)]
\item \textbf{Certified carrier pass.}
\[
\mathcal{I}\!\left(\mathrm{OP5}^{(b,r,s)}_{\mathrm{sym}}\right)
=
\mathcal{I}(\mathrm{OP1}).
\]
Hence OP~5 is to be read as a return-side readout of one common closure-bearing carrier class rather than as an independent closure test.
\item \textbf{First localized obstruction.} If a terminal structural failure is claimed at OP~5, then its first relevant source lies upstream of raw terminal visibility: it must already arise through loss of branch admissibility, loss of OP~3 shell-admissibility, or loss of the corresponding carrier/readout package before an independent OP~5 failure can be asserted.
\end{enumerate}
In particular, the terminal OP~5 layer is the first carrier-level decision interface after the branch-shell pass, not a new autonomous closure theorem.
\end{proposition}

\begin{proof}
By Theorem~\ref{thm:combined-branch-shell}, every branch-shell admissible realization satisfies
\[
\mathcal{I}\!\left(\mathrm{OP5}^{(b,r,s)}_{\mathrm{sym}}\right)=\mathcal{I}(\mathrm{OP1}).
\]
By Theorem~\ref{thm:op15-closure}, the OP~1 source readout and the symmetrized OP~5 return are already tied to one and the same native closure decision at invariant level. Thus the admissible branch-shell return is not a new endpoint test; it is a return-side reading of that same closure-bearing carrier class. This proves part~(i). Part~(ii) is the structural reading of the same fact: once the terminal invariant is fixed by the shared carrier decision, any genuinely structural failure visible at the terminal side must have entered earlier, namely through loss of admissibility on the branch-shell side or through failure of the carrier/readout package that supports the branch-shell pass. Therefore OP~5 is the first carrier-level decision interface after the branch-shell passage, not an independent closure theorem.
\end{proof}

\begin{corollary}[Carrier-first routed OP~5 decision principle]
\label{cor:carrier-first-routed-op5}
Assume the hypotheses of Proposition~\ref{prop:carrier-first-op5-interface} and let the branch-shell routed return
\[
\mathrm{OP5}_{\mathrm{bs\text{-}rout}}
\]
be defined with normalized routed weights. Then:
\begin{enumerate}[leftmargin=2em,label=(\roman*)]
\item the entire routed return is to be read first through one common closure-bearing carrier decision;
\item visible terminal multiplicity does not by itself create a new OP~5 closure problem;
\item any structural routed failure must already originate at the first localized obstruction of some branch-shell candidate rather than at terminal aggregation alone.
\end{enumerate}
\end{corollary}

\begin{proof}
Normalization together with Theorem~\ref{thm:combined-branch-shell} yields
\[
\mathcal{I}(\mathrm{OP5}_{\mathrm{bs\text{-}rout}})=\mathcal{I}(\mathrm{OP1}).
\]
Thus the routed return still belongs to the same invariant closure decision as the source readout. The remaining statements are the carrier-first reading of this equality: visible routed multiplicity is secondary, and any structural failure must already arise at the first localized obstruction of one of the contributing branch-shell realizations.
\end{proof}

\subsection{Carrier-first packet/phase-lock/epsilon compatibility layer}
\label{sec:packet-phase-epsilon-compatibility}

After the carrier-first OP~5 interface, the remaining packet-, phase-lock-, and epsilon-trap diagnostics should now be read as a visible support layer for the admissible return side of one common closure-bearing carrier. They do not replace the admissibility theorems or the carrier-first OP~5 decision; rather, they provide a structured non-escape corridor in which visible packet-scale deviations remain subordinate to the same return-side carrier decision.

\begin{definition}[Packet/phase-lock/epsilon compatible branch-shell channel]
A branch-shell realization $(b,r,s)$ is called \emph{packet/phase-lock/epsilon compatible} if the following hold simultaneously:
\begin{enumerate}[leftmargin=2em,label=(\roman*)]
\item the branch $b$ is closure-admissible;
\item the refinement $(r,s)$ is OP~3 shell-admissible;
\item the OP~1-side bridge normalization lies inside the contractive epsilon tube of Theorem~\ref{thm:working-op1-epsilon-trapped-bridge};
\item the OP~3-side visible shell ratio lies inside the anti-phase epsilon corridor preserved as \texttt{thm:working-shell-anti-phase-epsilon-trap} in the historical-witness PDF;
\item the packet-wave, breathing-envelope, and frame-flicker residuals are treated as visible-support diagnostics rather than as primary closure data.
\end{enumerate}
\end{definition}

\begin{remark}
The first two conditions are structural; the next two are non-escape conditions on the OP~1 entry side and the OP~3 shell side; the final condition places the remaining packet-scale evidence at the observable support layer. This is the natural compressed reading of the current packet/phase/epsilon analysis after OP closure hardening and after the carrier-first OP~5 interface has fixed the terminal decision level.
\end{remark}

\begin{corollary}[Carrier-first packet/phase-lock/epsilon compatibility]
\label{cor:packet-phase-epsilon-compatibility}
Assume the combined branch-shell transport theorem together with Proposition~\ref{prop:carrier-first-op5-interface}. If a branch-shell realization $(b,r,s)$ is packet/phase-lock/epsilon compatible, then the following hold.
\begin{enumerate}[leftmargin=2em,label=(\roman*)]
\item the visible packet-, breathing-envelope-, and frame-flicker diagnostics are compatible with one and the same invariantly closed source-return channel because they sit on the return side of one common closure-bearing carrier;
\item the OP~1 bridge side and the OP~3 shell side both satisfy a non-escape criterion inside their respective local epsilon corridors;
\item no visible packet- or phase-lock mismatch obtained inside those corridors is by itself evidence against structural OP closure or against the carrier-first OP~5 decision;
\item the first structurally relevant failure can occur only through loss of branch admissibility or loss of OP~3 shell-admissibility.
\end{enumerate}
\end{corollary}

\begin{proof}
By branch admissibility and OP~3 shell-admissibility, the combined branch-shell transport theorem gives
\[
\mathcal{I}(\mathrm{OP5}^{(b,r,s)}_{\mathrm{sym}})=\mathcal{I}(\mathrm{OP1}).
\]
By Proposition~\ref{prop:carrier-first-op5-interface}, this equality is already the carrier-first terminal OP~5 decision: the admissible return readout belongs to one common closure-bearing carrier and is not a new autonomous endpoint test. The OP~1 epsilon-trapped bridge theorem supplies the local non-escape statement on the source-entry side, while the shell anti-phase epsilon-trap supplies the corresponding non-escape statement for the visible shell ratio at the OP~3 interface. The packet-wave, breathing-envelope, and frame-flicker diagnostics are explicitly treated as visible-support diagnostics, so any mismatch they show inside the admissible corridors remains subordinate to that carrier-first decision and to the visible/structural mismatch separation already built into the OP tool layer. Therefore the visible diagnostics remain compatible with the same invariantly closed return-side carrier channel, and the first structurally relevant failure can only arise by leaving the admissible branch-shell sector itself.
\end{proof}

\begin{remark}[Current mathematical role of the packet/phase/epsilon line]
The present packet/phase-lock/epsilon line is therefore no longer best read as an alternative closure proof or as a second OP~5 test. Its mathematical role is to provide local corridor evidence that the currently observed OP~1 and OP~5 diagnostics can sit inside one admissible return-side carrier channel without forcing a structural mismatch. In that sense, it is the visible non-escape support layer beneath the carrier-first OP~5 decision and the next admissible-branch hardening layer.
\end{remark}

\begin{proposition}[Carrier-first packet/phase-lock/epsilon screening principle]
\label{prop:packet-phase-epsilon-screening}
Assume the combined branch-shell transport theorem together with Proposition~\ref{prop:carrier-first-op5-interface}. Then packet/phase-lock/epsilon compatibility defines a screening criterion for candidate admissible branch-shell support inside the carrier-first OP~5 architecture rather than a second closure theorem. More precisely, if a realization $(b,r,s)$ is packet/phase-lock/epsilon compatible, then it lies inside the diagnostic preselection corridor of the admissible branch-shell architecture and remains subordinate to the same return-side carrier decision; conversely, loss of packet/phase-lock/epsilon compatibility is not by itself sufficient to prove structural non-closure unless it propagates to loss of branch admissibility or loss of OP~3 shell-admissibility.
\end{proposition}

\begin{proof}
The first statement is immediate from Corollary~\ref{cor:packet-phase-epsilon-compatibility}: compatibility places the realization inside a joint non-escape corridor on the OP~1 bridge side and the OP~3 shell side while leaving the invariant closure value unchanged, and Proposition~\ref{prop:carrier-first-op5-interface} then fixes the terminal reading as one common return-side carrier decision. Hence the diagnostic role of the packet/phase/epsilon conditions is to preselect a corridor of candidate admissible support beneath that carrier-first terminal reading, not to introduce an independent closure criterion.

For the converse direction, failure of packet-, phase-lock-, or epsilon-corridor compatibility may occur either because a visible-support diagnostic leaves its controlled regime or because a structural admissibility condition fails. By the same corollary and by the carrier-first OP~5 interface, only the latter is primary for OP closure: visible mismatch inside or outside the packet-scale corridor is not yet structural non-closure unless it reflects loss of branch admissibility or loss of OP~3 shell-admissibility. Therefore packet/phase-lock/epsilon diagnostics define a screening layer rather than a standalone closure theorem or a second OP~5 endpoint test.
\end{proof}

\begin{remark}
This screening principle is the cleanest current mathematical reading of the old packet/phase-lock/epsilon diagnostics. They are now neither discarded nor promoted beyond their scope: they act as a visible preselection gate for the admissible branch-shell sector under one carrier-first OP~5 decision, while the actual structural test remains the admissibility machinery of the OP core.
\end{remark}


\subsection{Diagnostic readout hierarchy beneath the carrier-first OP~5 decision}
\label{sec:diagnostic-readout-hierarchy}

After the carrier-first OP~5 interface and the packet/phase-lock/epsilon screening layer, the remaining visible packet, breathing-envelope, frame-flicker, and helix pictures should be read in a strict order. They are not parallel closure criteria. They are progressively more visible readout surfaces of one and the same carrier decision.

\begin{remark}[Three diagnostic layers]
The current OP~5-side diagnostic stack is therefore to be read in the following descending order of structural weight.
\begin{enumerate}[leftmargin=2em,label=(\roman*)]
\item \textbf{Carrier-decision layer.} The first structural question is whether the return side still belongs to one common closure-bearing carrier, as fixed by Proposition~\ref{prop:carrier-first-op5-interface}.
\item \textbf{Screening layer.} Packet/phase-lock/epsilon compatibility then determines whether the visible return sits inside an admissible non-escape corridor beneath that carrier decision.
\item \textbf{Visible-shape layer.} Packet-wave, breathing-envelope, frame-flicker, and helix pictures come last: they visualize how one admissible or nearly admissible return is seen, but they do not outrank the carrier decision or the screening corridor.
\end{enumerate}
\end{remark}

\begin{lemma}[First diagnostic-readout hierarchy lemma]
\label{lem:diagnostic-readout-hierarchy}
Assume Proposition~\ref{prop:carrier-first-op5-interface} and Proposition~\ref{prop:packet-phase-epsilon-screening}. Let $(b,r,s)$ be a branch-shell realization with associated visible packet-, breathing-envelope-, frame-flicker-, and helix-type diagnostics. Then the current OP-side readout hierarchy has the following form.
\begin{enumerate}[leftmargin=2em,label=(\roman*)]
\item \textbf{Carrier priority.} If the realization belongs to one common closure-bearing carrier, then no visible packet-, flicker-, or helix-level multiplicity by itself overrides that carrier decision.
\item \textbf{Screening priority.} Packet/phase-lock/epsilon compatibility is the first visible screening layer beneath the carrier decision; loss at the pure visible-shape level is not yet structural loss unless it propagates to loss of this screening layer or to loss of the carrier package itself.
\item \textbf{Visible-shape subordination.} The packet-wave, breathing-envelope, frame-flicker, and helix pictures are therefore subordinate readout surfaces. They can localize, rank, and compare visible deviations, but they do not constitute an independent closure theorem and do not outrank the carrier-first OP~5 decision.
\end{enumerate}
In particular, the mathematically correct order of evidential force is
\[
\text{carrier decision}
\;>\;
\text{packet/phase/epsilon screening}
\;>\;
\text{visible flicker/helix picture}.
\]
\end{lemma}

\begin{proof}
Proposition~\ref{prop:carrier-first-op5-interface} fixes the first terminal structural question: the return side is read first through one common closure-bearing carrier, and a structural failure must already arise upstream of raw terminal visibility. Proposition~\ref{prop:packet-phase-epsilon-screening} then places packet/phase-lock/epsilon compatibility strictly beneath that terminal carrier decision as a diagnostic preselection corridor rather than as a second closure theorem. The remaining packet-wave, breathing-envelope, frame-flicker, and helix pictures are by construction even further downstream, because they are explicitly treated as visible-support diagnostics in Definition~3.116 and Corollary~\ref{cor:packet-phase-epsilon-compatibility}. Hence visible-shape multiplicity may refine the observable profile of a candidate return, but it cannot outrank the carrier decision and cannot become structural loss unless the screening corridor or the carrier package itself fails. This proves the stated hierarchy.
\end{proof}

\begin{corollary}[Visible flicker/helix subordination principle]
\label{cor:flicker-helix-subordination}
Within the current carrier-first OP~5 architecture, any packet-wave, breathing-envelope, frame-flicker, or helix mismatch is to be read first as visible-shape evidence beneath the carrier decision and beneath the packet/phase-lock/epsilon screening layer. Such a mismatch becomes structurally decisive only if it can be shown to force loss of packet/phase-lock/epsilon compatibility or loss of the common closure-bearing carrier itself.
\end{corollary}

\begin{proof}
This is the direct operational reading of Lemma~\ref{lem:diagnostic-readout-hierarchy}. The visible-shape diagnostics are subordinate readout surfaces; hence they can become structurally decisive only by propagating upward into the screening layer or the carrier layer.
\end{proof}

\subsection{First upstream propagation principle for visible-shape mismatch}
\label{sec:first-upstream-propagation}

\begin{remark}[Upstream propagation gate]
Under the current carrier-first OP~5 architecture, visible packet-wave, breathing-envelope, frame-flicker, and helix mismatches have no immediate right to overturn the structural reading. They may move upward only by forcing failure of the next superior layer: first the packet/phase-lock/epsilon screening corridor, and only after that the common closure-bearing carrier package itself.
\end{remark}

\begin{lemma}[First upstream propagation lemma]
\label{lem:first-upstream-propagation}
Assume Proposition~\ref{prop:carrier-first-op5-interface}, Proposition~\ref{prop:packet-phase-epsilon-screening}, Lemma~\ref{lem:diagnostic-readout-hierarchy}, and Lemma~\ref{lem:qam-shell-obstruction-stratification}. Let $(b,r,s)$ be a branch-shell realization with a visible packet-wave, breathing-envelope, frame-flicker, or helix mismatch. Then the current upstream propagation rule has the following form.
\begin{enumerate}[leftmargin=2em,label=(\roman*)]
\item \textbf{No direct jump.} A pure visible-shape mismatch does not by itself propagate directly to structural non-closure or directly overturn the carrier-first OP~5 decision.
\item \textbf{First gate: screening.} If the mismatch does \emph{not} force loss of packet/phase-lock/epsilon compatibility, then it remains subordinate visible-shape evidence and the carrier decision is unchanged.
\item \textbf{Second gate: carrier package.} If the mismatch \emph{does} force loss of packet/phase-lock/epsilon compatibility while the admissible branch-shell sector and the common closure-bearing carrier package remain intact, then the mismatch propagates only to screening-layer loss and is still not yet a structural OP failure.
\item \textbf{Structural propagation.} A visible-shape mismatch becomes structurally decisive only when it can be shown to propagate further to one of the upstream obstruction strata of Lemma~\ref{lem:qam-shell-obstruction-stratification}, namely loss of admissible branch passage, loss of shell-admissible refinement, or loss of the common readout/kernel carrier package.
\end{enumerate}
In particular, the current architecture permits only the ordered propagation chain
\[
\begin{aligned}
\text{visible flicker/helix mismatch}
&\;\Longrightarrow\;
\text{possible screening loss}\\
&\;\Longrightarrow\;
\text{possible carrier/package loss}\\
&\;\Longrightarrow\;
\text{structural failure}.
\end{aligned}
\]
No shorter upward propagation is presently licensed.
\end{lemma}

\begin{proof}
Lemma~\ref{lem:diagnostic-readout-hierarchy} already fixes the hierarchy of evidential force: visible-shape diagnostics lie strictly below the packet/phase-lock/epsilon screening layer, which itself lies below the carrier-first OP~5 decision. Hence a pure visible mismatch cannot jump directly to structural failure. Proposition~\ref{prop:packet-phase-epsilon-screening} gives the first gate: if packet/phase-lock/epsilon compatibility is preserved, then the mismatch remains inside the diagnostic corridor and does not alter the carrier reading. Even if that screening corridor is lost, Proposition~\ref{prop:carrier-first-op5-interface} still prevents immediate structural failure unless the common return-side carrier package itself is lost. Finally, Lemma~\ref{lem:qam-shell-obstruction-stratification} identifies the only current upstream structural failure modes: branch-stratum loss, shell-stratum loss, or readout/kernel carrier-package loss. Therefore the only admissible upstream propagation route is the ordered one stated above.
\end{proof}

\begin{corollary}[Visible mismatch non-propagation default]
\label{cor:visible-mismatch-nonpropagation}
Within the current carrier-first OP~5 architecture, every packet-wave, breathing-envelope, frame-flicker, or helix mismatch is to be treated by default as diagnostically local evidence. It may be upgraded only after an explicit witness of screening-layer loss or of upstream carrier/package loss has been exhibited.
\end{corollary}

\begin{proof}
This is the operational reformulation of Lemma~\ref{lem:first-upstream-propagation}. In the absence of an explicit upward propagation witness, visible-shape mismatch remains subordinate local evidence.
\end{proof}


\subsection{First diagnostic witness principle for visible-shape mismatch}
\label{sec:first-diagnostic-witness-principle}

\begin{remark}[Diagnostic witness discipline]
Under the current carrier-first OP~5 architecture, visible packet-wave, breathing-envelope, frame-flicker, and helix mismatches count as admissible upstream evidence only when they come equipped with a witness that is strong enough to certify failure of the next superior layer. Pure visual irregularity, even if persistent or large, is not yet such a witness.
\end{remark}

\begin{lemma}[First diagnostic witness lemma]
\label{lem:first-diagnostic-witness}
Assume Proposition~\ref{prop:packet-phase-epsilon-screening}, Proposition~\ref{prop:carrier-first-op5-interface}, Lemma~\ref{lem:diagnostic-readout-hierarchy}, and Lemma~\ref{lem:first-upstream-propagation}. Let $(b,r,s)$ be a branch-shell realization with a visible packet-wave, breathing-envelope, frame-flicker, or helix mismatch. Then the current witness discipline has the following form.
\begin{enumerate}[leftmargin=2em,label=(\roman*)]
\item \textbf{Screening witness threshold.} A visible mismatch is a legitimate witness of screening-layer loss only if it can be shown to force failure of packet/phase-lock/epsilon compatibility; otherwise it remains diagnostically local evidence.
\item \textbf{Carrier witness threshold.} A visible mismatch is a legitimate witness of carrier/package loss only if it can be shown to force, beyond screening loss, failure of the common closure-bearing carrier package itself.
\item \textbf{No witness by appearance alone.} Mere amplitude, persistence, oscillatory complexity, geometric deformation, or helix/flicker asymmetry of the visible mismatch is not by itself an admissible witness of screening-layer loss or of carrier/package loss.
\item \textbf{Admissible witness route.} Every currently admissible visible witness must therefore factor through one of the two certified upstream gates of Lemma~\ref{lem:first-upstream-propagation}: first explicit screening loss, and only then, if present, explicit carrier/package loss.
\end{enumerate}
Hence the present architecture admits visible mismatch only as \emph{witness-bearing} evidence and not as a free structural verdict.
\end{lemma}

\begin{proof}
Proposition~\ref{prop:packet-phase-epsilon-screening} defines the first superior layer above raw visible-shape mismatch. Therefore a visible mismatch can witness screening-layer loss only by forcing failure of packet/phase-lock/epsilon compatibility itself; if compatibility persists, the mismatch remains inside the diagnostic corridor and cannot count as screening evidence. Proposition~\ref{prop:carrier-first-op5-interface} defines the next superior layer: even after screening loss, the structural carrier reading remains unchanged unless the common closure-bearing carrier package is lost. Thus a visible mismatch can witness carrier/package loss only by forcing precisely that second failure as well. Lemma~\ref{lem:diagnostic-readout-hierarchy} excludes any direct elevation of raw visible-shape features to a superior structural layer, and Lemma~\ref{lem:first-upstream-propagation} identifies the only currently licensed upward route as the ordered screening-to-carrier route. Consequently amplitude, persistence, or geometric shape alone are not admissible witnesses; only a mismatch that certifiably drives one of those gates may count as an upstream witness.
\end{proof}

\begin{corollary}[Visible witness admissibility rule]
\label{cor:visible-witness-admissibility}
Within the current carrier-first OP~5 architecture, every visible packet-wave, breathing-envelope, frame-flicker, or helix mismatch must be classified first as either
\begin{enumerate}[leftmargin=2em,label=(\alph*)]
\item a merely local diagnostic irregularity,
\item a certified witness of screening-layer loss,
\item or a certified witness of carrier/package loss.
\end{enumerate}
No fourth class of immediately structural visible witness is presently admitted.
\end{corollary}

\begin{proof}
This is the direct classification consequence of Lemma~\ref{lem:first-diagnostic-witness}. A visible mismatch is either not strong enough to force a superior-layer failure, or it forces screening loss, or it forces carrier/package loss after screening loss. By the witness discipline and by the upstream propagation rule, there is no currently licensed direct jump from raw visible appearance to immediate structural verdict.
\end{proof}


\subsection{Screening-witness compression beneath the carrier-first OP~5 decision}
\label{sec:screening-witness-compression}

\begin{remark}[Compression of visible witness types]
Once a visible packet-wave, breathing-envelope, frame-flicker, or helix mismatch has crossed the witness threshold of Lemma~\ref{lem:first-diagnostic-witness}, the present architecture no longer needs to keep its visible shape as a separate structural species at the screening level. Unless the same witness already certifies carrier/package loss, its upstream content compresses to one common statement: the packet/phase-lock/epsilon screening corridor has failed.
\end{remark}

\begin{lemma}[Screening-witness compression lemma]
\label{lem:screening-witness-compression}
Assume Proposition~\ref{prop:packet-phase-epsilon-screening}, Proposition~\ref{prop:carrier-first-op5-interface}, Lemma~\ref{lem:first-diagnostic-witness}, and Corollary~\ref{cor:visible-witness-admissibility}. Let $(b,r,s)$ be a branch-shell realization with a visible packet-wave, breathing-envelope, frame-flicker, or helix mismatch that is a certified witness of screening-layer loss but not yet a certified witness of carrier/package loss. Then the upstream structural content of that witness compresses to a single screening-loss form, namely the failure of packet/phase-lock/epsilon compatibility. More precisely:
\begin{enumerate}[leftmargin=2em,label=(\roman*)]
\item \textbf{Visible-shape indifference at screening level.} Whether the witness is first expressed through packet-wave mismatch, breathing-envelope drift, frame-flicker asymmetry, or helix deformation is irrelevant for its screening-level structural role.
\item \textbf{Unique compressed content.} The only currently licensed upstream content of such a witness is that the packet/phase-lock/epsilon screening corridor fails.
\item \textbf{No parallel screening species.} There is presently no second independent screening-loss class attached specifically to packet-wave, envelope, flicker, or helix shape once the witness has been admitted.
\item \textbf{Carrier exception.} A visible witness escapes this compression only when it upgrades further to a certified witness of carrier/package loss.
\end{enumerate}
Hence, below the carrier threshold, all currently admissible visible screening witnesses collapse to one common screening-loss normal form.
\end{lemma}

\begin{proof}
By Lemma~\ref{lem:first-diagnostic-witness}, a visible mismatch is admitted upward only when it certifiably forces failure of the next superior layer. Proposition~\ref{prop:packet-phase-epsilon-screening} identifies that next superior layer uniquely as packet/phase-lock/epsilon compatibility. Therefore, once a visible mismatch qualifies as a witness of screening-layer loss, its upstream meaning is already fixed: it testifies precisely to failure of that screening corridor. Corollary~\ref{cor:visible-witness-admissibility} allows no additional currently licensed class of screening witness beyond that certified screening loss or the stronger carrier/package loss. Proposition~\ref{prop:carrier-first-op5-interface} then separates the only exception: a witness that also forces loss of the common closure-bearing carrier package belongs to the higher carrier class instead of remaining in the compressed screening class. Thus all admitted visible witnesses that stop below the carrier threshold share one and the same screening-loss content, independently of the visible shape through which they first appeared.
\end{proof}

\begin{corollary}[Screening-witness normal form]
\label{cor:screening-witness-normal-form}
Within the current carrier-first OP~5 architecture, every admitted visible witness that does not yet certify carrier/package loss may be cited in proofs simply as a witness of \emph{screening-loss normal form}. Its packet-wave, breathing-envelope, frame-flicker, or helix realization is diagnostically informative, but not structurally independent at the screening level.
\end{corollary}

\begin{proof}
This is the direct proof-usage reformulation of Lemma~\ref{lem:screening-witness-compression}. Once admitted and still below the carrier threshold, the witness has exactly one currently licensed upstream content: loss of the packet/phase-lock/epsilon screening corridor.
\end{proof}


\subsection{Carrier-witness separation beneath compressed screening loss}
\label{sec:carrier-witness-separation}

\begin{remark}[Two witness grades after screening compression]
After Lemma~\ref{lem:screening-witness-compression}, every admitted visible witness sits in one of two remaining grades. Either it stays below the carrier threshold and is therefore only a witness of screening-loss normal form, or it upgrades further and certifies that the common return-side carrier package itself is no longer available. The present architecture does not license a mixed intermediate status between those two grades.
\end{remark}

\begin{lemma}[Carrier-witness separation lemma]
\label{lem:carrier-witness-separation}
Assume Proposition~\ref{prop:packet-phase-epsilon-screening}, Proposition~\ref{prop:carrier-first-op5-interface}, Lemma~\ref{lem:first-diagnostic-witness}, Lemma~\ref{lem:screening-witness-compression}, and Corollary~\ref{cor:screening-witness-normal-form}. Let $(b,r,s)$ be a branch-shell realization with an admitted visible witness. Then exactly one of the following two cases occurs:
\begin{enumerate}[leftmargin=2em,label=(\roman*)]
\item \textbf{Screening-only witness.} The witness compresses to screening-loss normal form, while the common closure-bearing return-side carrier package remains available. In this case the witness certifies only screening loss and is not an admissible witness of carrier/package loss.
\item \textbf{Carrier-upgraded witness.} The witness forces loss of the common closure-bearing return-side carrier package. In this case it is an admissible witness of carrier/package loss and no longer belongs merely to the screening-loss normal form.
\end{enumerate}
Moreover, there is no third currently licensed status in which an admitted visible witness would simultaneously remain only screening-level and yet already count as a carrier/package witness.
\end{lemma}

\begin{proof}
By Lemma~\ref{lem:first-diagnostic-witness}, every admitted visible witness must force loss of some superior layer. Lemma~\ref{lem:screening-witness-compression} and Corollary~\ref{cor:screening-witness-normal-form} show that, below the carrier threshold, all such witnesses have one and the same upstream content: failure of packet/phase-lock/epsilon compatibility. Proposition~\ref{prop:carrier-first-op5-interface} identifies the only stronger alternative, namely loss of the common closure-bearing return-side carrier package itself. Hence an admitted visible witness either stops at the compressed screening-loss level or upgrades to carrier/package loss. These two statuses are mutually exclusive because the first expressly retains the common carrier package, whereas the second destroys it. Therefore exactly one of the two cases above occurs, and no mixed intermediate category is presently licensed.
\end{proof}

\begin{corollary}[Carrier-witness upgrade criterion]
\label{cor:carrier-witness-upgrade-criterion}
Within the current carrier-first OP~5 architecture, an admitted visible witness upgrades from screening-loss normal form to a genuine carrier/package witness if and only if its upstream content cannot be realized through any common closure-bearing return-side carrier package. Equivalently, as long as such a common carrier remains available, the witness must be read only as screening-loss normal form.
\end{corollary}

\begin{proof}
The forward direction is exactly case (ii) of Lemma~\ref{lem:carrier-witness-separation}. The converse direction is case (i): if a common closure-bearing return-side carrier package still exists, then the witness has not crossed the carrier threshold and therefore remains only a screening witness in normal form.
\end{proof}


\subsection{Witness-to-carrier escalation beyond screening compression}
\label{sec:witness-to-carrier-escalation}

\begin{remark}[Escalation gate after witness separation]
After Lemma~\ref{lem:carrier-witness-separation}, a visible witness reaches carrier grade exactly when its admitted upstream content can no longer be absorbed by screening-loss normal form alone. In the current architecture there is therefore only one licensed upgrade mechanism: every common closure-bearing return-side carrier package compatible with the compressed screening witness fails, so the witness is forced out of screening-only status and into carrier/package loss.
\end{remark}

\begin{lemma}[Witness-to-carrier escalation lemma]
\label{lem:witness-to-carrier-escalation}
Assume Proposition~\ref{prop:carrier-first-op5-interface}, Lemma~\ref{lem:first-diagnostic-witness}, Lemma~\ref{lem:screening-witness-compression}, Lemma~\ref{lem:carrier-witness-separation}, and Corollary~\ref{cor:carrier-witness-upgrade-criterion}. Let $(b,r,s)$ be a branch-shell realization with an admitted visible witness $w$. Then the following are equivalent:
\begin{enumerate}[leftmargin=2em,label=(\roman*)]
\item $w$ upgrades from screening-loss normal form to a genuine carrier/package witness;
\item there exists no common closure-bearing return-side carrier package through which the upstream content of $w$ can still be realized while screening-loss normal form remains the only failed superior layer;
\item the upstream content of $w$ forces loss of the common closure-bearing carrier package rather than mere loss of packet/phase-lock/epsilon compatibility.
\end{enumerate}
In particular, escalation occurs exactly when compressed screening loss is no longer a sufficient upstream normal form for the admitted witness.
\end{lemma}

\begin{proof}
By Lemma~\ref{lem:carrier-witness-separation}, an admitted visible witness has only two currently licensed grades: screening-only or carrier-upgraded. Corollary~\ref{cor:carrier-witness-upgrade-criterion} identifies the upgrade threshold precisely: upgrade occurs if and only if no common closure-bearing return-side carrier package remains available. Lemma~\ref{lem:screening-witness-compression} shows that, below that threshold, the witness has exactly one upstream content, namely screening-loss normal form. Hence a witness stays screening-only exactly as long as its upstream content is still realizable through some common carrier package with only packet/phase-lock/epsilon compatibility lost. Once that becomes impossible, the witness can no longer remain in screening-loss normal form alone and must instead force carrier/package loss. This proves the equivalence of (i), (ii), and (iii), and therefore the final escalation claim.
\end{proof}

\begin{corollary}[Escalation threshold for admitted visible witnesses]
\label{cor:carrier-escalation-threshold}
Within the current carrier-first OP~5 architecture, an admitted visible witness remains screening-only as long as its upstream content is representable by screening-loss normal form over some common closure-bearing return-side carrier package. It becomes carrier-upgraded precisely at the first point where such a carrier-supported realization ceases to exist.
\end{corollary}

\begin{proof}
This is the proof-usage reformulation of Lemma~\ref{lem:witness-to-carrier-escalation}. The existence of a common carrier package keeps the witness below the carrier threshold, while the first failure of every such carrier-supported realization marks the unique licensed escalation point.
\end{proof}


\subsection{Carrier-stable diagnostic collapse beneath the witness escalation threshold}
\label{sec:carrier-stable-diagnostic-collapse}

\begin{remark}[Collapse of visible diagnostic species over a stable carrier]
Once an admitted visible mismatch stays below the escalation threshold of Corollary~\ref{cor:carrier-escalation-threshold}, the current carrier-first OP~5 architecture no longer attributes independent structural rank to its packet-wave, breathing-envelope, frame-flicker, or helix presentation. Those visible realizations remain diagnostically useful, but over one and the same stable closure-bearing return-side carrier they collapse to a single subordinate screening-level phenomenon.
\end{remark}

\begin{theorem}[Carrier-stable diagnostic collapse theorem]
\label{thm:carrier-stable-diagnostic-collapse}
Assume Proposition~\ref{prop:carrier-first-op5-interface}, Proposition~\ref{prop:packet-phase-epsilon-screening}, Lemma~\ref{lem:diagnostic-readout-hierarchy}, Lemma~\ref{lem:first-diagnostic-witness}, Lemma~\ref{lem:screening-witness-compression}, Lemma~\ref{lem:carrier-witness-separation}, and Lemma~\ref{lem:witness-to-carrier-escalation}. Let $(b,r,s)$ be a branch-shell realization with an admitted visible mismatch $w$ of packet-wave, breathing-envelope, frame-flicker, or helix type. If $w$ remains below carrier escalation, equivalently if there exists a common closure-bearing return-side carrier package through which the upstream content of $w$ is still realized as screening-loss normal form, then the following collapse holds:
\begin{enumerate}[leftmargin=2em,label=(\roman*)]
\item \textbf{One subordinate structural role.} The visible realization of $w$ has exactly one current structural role, namely to witness subordinate screening-loss normal form beneath a still-stable carrier decision.
\item \textbf{No independent visible closure species.} Packet-wave, breathing-envelope, frame-flicker, and helix mismatch do not define distinct closure-relevant species as long as the common carrier package remains available.
\item \textbf{Carrier dominance.} Any closure-relevant distinction between those visible realizations collapses over the stable carrier and may reappear only if escalation to carrier/package loss occurs.
\item \textbf{Readout subordination.} The visible mismatch remains a downstream readout of one common carrier-supported screening loss and therefore cannot outrank the carrier decision itself.
\end{enumerate}
Hence, below the escalation threshold, all currently admitted visible diagnostic forms collapse to one carrier-stable subordinate diagnostic class.
\end{theorem}

\begin{proof}
By Lemma~\ref{lem:diagnostic-readout-hierarchy}, visible flicker/helix-type structure is subordinate to packet/phase/epsilon screening, which is itself subordinate to the carrier-first OP~5 decision. Lemma~\ref{lem:first-diagnostic-witness} allows visible mismatch to count upstream only when it certifies loss of a superior layer. Lemma~\ref{lem:screening-witness-compression} then shows that, below carrier loss, every admitted visible witness has exactly one upstream content: screening-loss normal form. Lemma~\ref{lem:carrier-witness-separation} and Corollary~\ref{cor:carrier-escalation-threshold} identify the present hypothesis of a still available common closure-bearing return-side carrier package as precisely the condition that the witness has not escalated to carrier/package loss. Therefore, while that common carrier package remains available, packet-wave, envelope, flicker, and helix mismatch cannot persist as distinct closure-relevant species. They all collapse to the single role of subordinate diagnostic realization of screening-loss normal form over one stable carrier. Proposition~\ref{prop:carrier-first-op5-interface} guarantees that any stronger distinction would already belong to carrier/package loss rather than to the pre-escalation regime. This proves (i)--(iv) and the final collapse claim.
\end{proof}

\begin{corollary}[Carrier-stable visible normal form]
\label{cor:carrier-stable-visible-normal-form}
Within the current carrier-first OP~5 architecture, every admitted visible mismatch that stays below carrier escalation may be cited simply as belonging to one \emph{carrier-stable visible normal form}. Its packet-wave, breathing-envelope, frame-flicker, or helix appearance remains diagnostically informative, but not structurally independent while the common carrier package persists.
\end{corollary}

\begin{proof}
This is the proof-usage reformulation of Theorem~\ref{thm:carrier-stable-diagnostic-collapse}. The theorem states exactly that all admitted visible mismatches below escalation collapse over one stable carrier to a single subordinate diagnostic class.
\end{proof}


\subsection{Bundled OP~5 diagnostic theorem}
\label{sec:bundled-op5-diagnostic-theorem}

\begin{remark}[One bundled positive/negative OP~5 diagnostic decision]
After the carrier-first OP~5 interface, the current diagnostic layer is now strong enough to be cited as one bundled decision principle. Either the visible material remains subordinate to a stable carrier and therefore collapses to a screening-level normal form, or it upgrades through an admitted witness to carrier/package loss. There is no third currently admitted diagnostic rank between those two bundled outcomes.
\end{remark}

\begin{theorem}[Bundled OP~5 diagnostic theorem]
\label{thm:bundled-op5-diagnostic-theorem}
Assume Proposition~\ref{prop:carrier-first-op5-interface}, Proposition~\ref{prop:packet-phase-epsilon-screening}, Lemma~\ref{lem:diagnostic-readout-hierarchy}, Lemma~\ref{lem:first-upstream-propagation}, Lemma~\ref{lem:first-diagnostic-witness}, Lemma~\ref{lem:screening-witness-compression}, Lemma~\ref{lem:carrier-witness-separation}, Lemma~\ref{lem:witness-to-carrier-escalation}, and Theorem~\ref{thm:carrier-stable-diagnostic-collapse}. Let $(b,r,s)$ be a routed branch-shell realization together with an admitted visible mismatch $w$ of packet-wave, breathing-envelope, frame-flicker, or helix type. Then exactly one of the following bundled OP~5 diagnostic outcomes holds:
\begin{enumerate}[leftmargin=2em,label=(\roman*)]
\item \textbf{Carrier-stable diagnostic collapse.} The mismatch $w$ remains below escalation, is carried by one common closure-bearing return-side carrier, and therefore collapses to a subordinate screening-level visible normal form beneath the carrier-first OP~5 decision.
\item \textbf{Carrier-upgraded diagnostic obstruction.} The mismatch $w$ cannot be realized over any common closure-bearing return-side carrier as mere screening-loss normal form and therefore upgrades to carrier/package loss. In that case $w$ is not a merely visible irregularity but an admitted carrier witness for routed OP~5 failure.
\end{enumerate}
Moreover, case \textup{(i)} excludes case \textup{(ii)}, and case \textup{(ii)} excludes case \textup{(i)}. Hence the current OP~5 diagnostic layer is completely bundled into the dichotomy
\[
\text{carrier-stable visible collapse}
\qquad\text{versus}\qquad
\text{carrier-upgraded diagnostic obstruction}.
\]
\end{theorem}

\begin{proof}
Lemma~\ref{lem:diagnostic-readout-hierarchy} and Lemma~\ref{lem:first-upstream-propagation} force every admitted visible mismatch to move upward only through the packet/phase/epsilon screening layer before it can acquire carrier relevance. Lemma~\ref{lem:first-diagnostic-witness} then identifies which visible mismatches may count upstream at all. Lemma~\ref{lem:screening-witness-compression} compresses every such admitted visible witness below carrier loss to one screening-loss normal form, while Lemma~\ref{lem:carrier-witness-separation} and Lemma~\ref{lem:witness-to-carrier-escalation} separate precisely the screening-only regime from the carrier-upgraded regime. Theorem~\ref{thm:carrier-stable-diagnostic-collapse} describes the entire positive side of the first regime: once a common closure-bearing return-side carrier remains available, all admitted visible species collapse to one subordinate visible normal form beneath the carrier-first OP~5 decision. Conversely, when no such common carrier remains available, Corollary~\ref{cor:carrier-witness-upgrade-criterion} upgrades the admitted visible witness to carrier/package loss. These two possibilities are mutually exclusive by Lemma~\ref{lem:carrier-witness-separation} and jointly exhaustive by the escalation dichotomy. This proves the bundled theorem.
\end{proof}

\begin{corollary}[Bundled routed OP~5 diagnostic use rule]
\label{cor:bundled-routed-op5-diagnostic-use-rule}
In later proof steps it is sufficient to cite Theorem~\ref{thm:bundled-op5-diagnostic-theorem} instead of separately unpacking packet-wave, breathing-envelope, frame-flicker, and helix diagnostics. One may simply ask whether the currently admitted visible data still collapse over a stable carrier or have already upgraded to carrier/package obstruction.
\end{corollary}

\begin{proof}
This is the proof-usage reformulation of Theorem~\ref{thm:bundled-op5-diagnostic-theorem}. The theorem already packages the visible diagnostic layer into exactly the two current routed OP~5 outcomes.
\end{proof}


\subsection{First OP~3 to OP~5 bundled transfer theorem}
\label{sec:first-op3-op5-bundled-transfer}

\begin{remark}[Transfer of bundled OP~3 closure into the OP~5 diagnostic layer]
The current architecture no longer treats OP~3 closure and OP~5 diagnosis as two independent theorem blocks. Once the OP~3 shell passage has been certified and rerouted through one common closure-bearing carrier, the routed OP~5 side inherits a downstream decision of that same carrier package. The only remaining terminal question is whether the visible return side still collapses beneath that package or whether it upgrades to a carrier-level obstruction. Thus OP~5 is presently read as the bundled downstream transfer image of the OP~3/QAM closure package.
\end{remark}

\begin{theorem}[First OP~3 to OP~5 bundled transfer theorem]
\label{thm:first-op3-op5-bundled-transfer}
Assume Theorem~\ref{thm:qam-op3-shell-closure}, Theorem~\ref{thm:qam-op3-proof-run}, Theorem~\ref{thm:combined-branch-shell}, Proposition~\ref{prop:carrier-first-op5-interface}, and Theorem~\ref{thm:bundled-op5-diagnostic-theorem}. Let $(b,r,s)$ be a routed branch-shell realization whose OP~3 shell passage is translated into the extracted QAM layer, and let $w$ be an admitted visible mismatch on the routed OP~5 return side. Then the present bundled OP~3$\to$OP~5 transfer has exactly two outcomes:
\begin{enumerate}[leftmargin=2em,label=(\roman*)]
\item \textbf{Transferred carrier-stable pass.}
The translated OP~3 shell passage belongs to the certified side of the bundled OP~3 closure package, the routed return stays inside the same common closure-bearing carrier, and the admitted visible mismatch $w$ collapses beneath that carrier to the subordinate OP~5 visible normal form of Theorem~\ref{thm:bundled-op5-diagnostic-theorem}.
\item \textbf{Transferred localized obstruction.}
The transfer fails at the first licensed obstruction layer: either the OP~3/QAM side already reaches the first localized shell obstruction of the extracted OP~3 proof run, or the routed return upgrades to a carrier-level OP~5 diagnostic obstruction. In particular, no third autonomous terminal OP~5 obstruction is presently licensed downstream of an otherwise intact bundled OP~3 carrier package.
\end{enumerate}
Hence the current routed OP~5 diagnostic layer is the bundled downstream transfer of the OP~3/QAM closure package, not an additional independent closure theory.
\end{theorem}

\begin{proof}
Theorem~\ref{thm:qam-op3-proof-run} already packages the translated OP~3 shell passage into exactly two upstream outcomes: certified passage through the bundled OP~3 closure tool layer or the first localized obstruction. If the certified side holds, then Theorem~\ref{thm:qam-op3-shell-closure} identifies one common closure-bearing carrier/readout package for the shell passage. Theorem~\ref{thm:combined-branch-shell} transports that same certified branch-shell realization to the routed return side without changing the invariant closure class, and Proposition~\ref{prop:carrier-first-op5-interface} forces the terminal OP~5 reading to be carrier-first rather than a new autonomous endpoint test. Theorem~\ref{thm:bundled-op5-diagnostic-theorem} then completes the positive branch: as long as the routed return remains over that same common carrier, every admitted visible mismatch collapses to the subordinate carrier-stable OP~5 visible normal form. This proves case~(i).

If the transfer does not remain on that certified branch, then the failure must appear at the first licensed obstruction layer. Upstream, Theorem~\ref{thm:qam-op3-proof-run} and the extracted OP~3/QAM obstruction package localize the first obstruction already at branch, shell, or carrier/readout level. Downstream, if the routed return cannot be realized over any common closure-bearing carrier, then Theorem~\ref{thm:bundled-op5-diagnostic-theorem} upgrades the admitted visible mismatch to a carrier-level OP~5 diagnostic obstruction. Because Proposition~\ref{prop:carrier-first-op5-interface} forbids an independent terminal closure test downstream of an otherwise intact carrier package, there is no currently licensed third obstruction type appearing only at the raw terminal visible layer. This proves case~(ii) and the final bundled transfer reading.
\end{proof}

\begin{corollary}[Bundled OP~3$\to$OP~5 proof-use rule]
\label{cor:bundled-op3-op5-proof-use-rule}
In later proof steps it is sufficient to cite Theorem~\ref{thm:first-op3-op5-bundled-transfer} when a routed branch-shell realization must be carried from the extracted OP~3/QAM closure package into the OP~5 diagnostic layer. One may simply ask whether the realization stays on the transferred carrier-stable branch or reaches the first licensed transferred obstruction.
\end{corollary}

\begin{proof}
This is the proof-usage reformulation of Theorem~\ref{thm:first-op3-op5-bundled-transfer}. The theorem already packages the current OP~3/QAM closure side and the current OP~5 diagnostic side into one bundled transfer decision.
\end{proof}


\subsection{First OP~3 to OP~5 compressed proof interface}
\label{sec:first-op3-op5-compressed-proof-interface}

\begin{remark}[How older routed return arguments should now be compressed]
The current transfer architecture is now strong enough that later routed return arguments need not separately reopen the combined branch-shell theorem, the carrier-first OP~5 interface, and the bundled OP~5 diagnostic layer. Once the routed realization has entered the bundled OP~3$\to$OP~5 transfer channel, the proof may proceed entirely through one compressed dichotomy: transferred carrier-stable pass versus transferred localized obstruction. This local compressed interface is now itself subordinate to the newer global rule: as soon as the argument has also entered the bundled current channel and no shell-sensitive or witness-grade claim is being asked, the proof should move on to the extracted global tools and remain global-first unless a genuine local reopening is needed.
\end{remark}

\begin{proposition}[First OP~3 to OP~5 compressed proof interface]
\label{prop:first-op3-op5-compressed-proof-interface}
Assume Theorem~\ref{thm:first-op3-op5-bundled-transfer}. Let $(b,r,s)$ be a routed branch-shell realization whose OP~3 shell passage has been translated into the extracted QAM layer, and let $w$ be any admitted visible routed return witness on the OP~5 side. Then every later proof step that needs to carry this realization through the current routed return architecture may be compressed to the following interface.
\begin{enumerate}[leftmargin=2em,label=(\roman*)]
\item \textbf{Transferred carrier-stable pass.} If Theorem~\ref{thm:first-op3-op5-bundled-transfer} places $(b,r,s;w)$ on the transferred carrier-stable branch, then all downstream OP~5 packet/phase/flicker/helix material is read only as subordinate visible realization beneath one common closure-bearing carrier and may be replaced by the carrier-stable visible normal form.
\item \textbf{Transferred localized obstruction.} If Theorem~\ref{thm:first-op3-op5-bundled-transfer} places $(b,r,s;w)$ on the transferred obstruction branch, then the proof may stop at the first licensed obstruction already identified there: either the upstream OP~3/QAM obstruction or the downstream carrier-upgraded OP~5 obstruction.
\item \textbf{No third terminal proof branch.} No additional raw terminal visible branch is presently licensed after this compression. In particular, later proofs do not need a separate fourth case for packet-wave, breathing-envelope, frame-flicker, or helix mismatch once the bundled transfer theorem has already been invoked.
\end{enumerate}
Hence the current routed return architecture may be used through one compressed OP~3$\to$OP~5 proof interface.
\end{proposition}

\begin{proof}
This is the direct proof-usage reading of Theorem~\ref{thm:first-op3-op5-bundled-transfer}. The theorem already packages the routed return into exactly two licensed transfer outcomes. On the positive side, the admitted visible material collapses beneath one common carrier and therefore needs no further structural branching. On the negative side, the first licensed obstruction is already localized either upstream in the OP~3/QAM package or downstream at the carrier-upgraded OP~5 level. Because Theorem~\ref{thm:first-op3-op5-bundled-transfer} excludes an autonomous raw terminal OP~5 obstruction beneath an otherwise intact carrier package, no third terminal proof branch remains.
\end{proof}

\begin{corollary}[Compressed rereading of routed return arguments through the bundled transfer]
\label{cor:compressed-routed-return-through-bundled-transfer}
In later arguments it is sufficient to cite Proposition~\ref{prop:first-op3-op5-compressed-proof-interface} rather than separately reopening Theorem~\ref{thm:combined-branch-shell}, Proposition~\ref{prop:carrier-first-op5-interface}, and Theorem~\ref{thm:bundled-op5-diagnostic-theorem}. One may simply ask whether the current routed realization stays on the transferred carrier-stable branch or reaches the first licensed transferred obstruction.
\end{corollary}

\begin{proof}
This is the condensed proof-usage form of Proposition~\ref{prop:first-op3-op5-compressed-proof-interface}. The proposition already compresses the presently licensed routed return architecture into one reusable interface.
\end{proof}


\subsection{First routed return normal form theorem}
\label{sec:first-routed-return-normal-form-theorem}

\begin{remark}[Carrier-stable versus obstructed routed return normal form]
After the compressed OP~3$\to$OP~5 proof interface has been invoked, the current routed return architecture no longer needs a separate case distinction for packet-wave, breathing-envelope, frame-flicker, or helix appearance. At the present level of the theory every admitted routed return is read only up to one of two bundled normal forms: a carrier-stable routed return normal form or an obstructed routed return normal form.
\end{remark}

\begin{theorem}[First routed return normal form theorem]
\label{thm:first-routed-return-normal-form}
Assume Proposition~\ref{prop:first-op3-op5-compressed-proof-interface}. Let $(b,r,s)$ be a routed branch-shell realization whose OP~3 shell passage has been translated into the extracted QAM layer, and let $w$ be any admitted visible routed return witness on the OP~5 side. Then the current routed return architecture classifies $(b,r,s;w)$, up to bundled transfer, into exactly one of the following two mutually exclusive routed return normal forms.
\begin{enumerate}[leftmargin=2em,label=(\roman*)]
\item \textbf{Carrier-stable routed return normal form.}
The routed return remains on the transferred carrier-stable branch of Proposition~\ref{prop:first-op3-op5-compressed-proof-interface}. Equivalently, there exists one common closure-bearing return-side carrier through which the routed realization is read, and every admitted visible mismatch collapses beneath that carrier to the carrier-stable visible normal form. In this case the routed return contributes no additional structural branching beyond the bundled OP~3$\to$OP~5 transfer itself.
\item \textbf{Obstructed routed return normal form.}
The routed return lies on the transferred localized-obstruction branch of Proposition~\ref{prop:first-op3-op5-compressed-proof-interface}. Equivalently, no common closure-bearing routed return carrier supports the realization as mere screening-level visible loss, and the first licensed failure is already localized either upstream at the OP~3/QAM obstruction layer or downstream at the carrier-upgraded OP~5 obstruction layer.
\end{enumerate}
Hence every currently admitted routed return is read, up to the bundled OP~3$\to$OP~5 transfer, by one bundled routed return normal form and not by any third autonomous terminal visible class.
\end{theorem}

\begin{proof}
Proposition~\ref{prop:first-op3-op5-compressed-proof-interface} already compresses the present routed return architecture into exactly two licensed proof branches: transferred carrier-stable pass and transferred localized obstruction. On the positive branch, Theorem~\ref{thm:bundled-op5-diagnostic-theorem} and the carrier-stable collapse package imply that every admitted visible return witness becomes subordinate to one common closure-bearing routed return carrier and therefore yields the carrier-stable routed return normal form. On the negative branch, the same compressed interface localizes the first licensed failure either upstream in the extracted OP~3/QAM obstruction package or downstream at the carrier-upgraded OP~5 obstruction layer; by construction this is the obstructed routed return normal form. Because Proposition~\ref{prop:first-op3-op5-compressed-proof-interface} excludes a third raw terminal visible branch after compression, these two forms are exhaustive and mutually exclusive.
\end{proof}

\begin{corollary}[Routed return normal form use rule]
\label{cor:routed-return-normal-form-use-rule}
In later proofs it is sufficient to cite Theorem~\ref{thm:first-routed-return-normal-form} once a routed realization has already entered the bundled OP~3$\to$OP~5 transfer channel. One may then work only with the dichotomy \emph{carrier-stable routed return normal form versus obstructed routed return normal form}.
\end{corollary}

\begin{proof}
This is the proof-usage reformulation of Theorem~\ref{thm:first-routed-return-normal-form}. The theorem already packages the presently licensed routed return architecture into one reusable normal-form dichotomy.
\end{proof}


\subsection{First global pass/obstruction normal form}
\label{sec:first-global-pass-obstruction-normal-form}

\begin{remark}[Current global normal form of the admissible run]
The present line now carries enough bundled structure that the currently admitted run
\[
\begin{aligned}
&\text{OP1 source}
\;\to\;\text{branch admissibility}\\[0.15em]
&\;\to\;\text{OP3/QAM closure package}\\[0.15em]
&\;\to\;\text{carrier-first OP5 layer}
\end{aligned}
\]
no longer needs to be read as a chain of loosely connected local tests. After the routed
return normal form has been established, the full currently admitted run is read only up to
one global dichotomy: a global pass normal form or a global obstruction normal form.
\end{remark}

\begin{theorem}[First global pass/obstruction normal form]
\label{thm:first-global-pass-obstruction-normal-form}
Assume Theorem~\ref{thm:first-routed-return-normal-form}. Let $(b,r,s)$ be an admissible
branch-shell realization issued from the current OP1 source side and translated through the
extracted OP3/QAM layer into the routed return architecture, and let $w$ be any admitted
visible routed return witness on the OP5 side. Then the whole currently admitted run from OP1
through branch, OP3/QAM, and OP5 is classified, up to the present bundled architecture, into
exactly one of the following two mutually exclusive global normal forms.
\begin{enumerate}[leftmargin=2em,label=(\roman*)]
\item \textbf{Global pass normal form.}
The realization lies on the carrier-stable routed return normal form of
Theorem~\ref{thm:first-routed-return-normal-form}. Equivalently, the run remains supported by
one bundled closure-bearing carrier package from the present OP1/branch entrance through the
OP3/QAM shell package and into the routed return side, while all visible OP5 material collapses
beneath that carrier to the carrier-stable visible normal form. In this case the current run
carries no further licensed structural branch beyond the bundled pass reading.
\item \textbf{Global obstruction normal form.}
The realization lies on the obstructed routed return normal form of
Theorem~\ref{thm:first-routed-return-normal-form}. Equivalently, the first licensed failure of
the current run is already localized either upstream in the OP3/QAM obstruction package or
downstream at the carrier-upgraded OP5 obstruction layer, and no autonomous third terminal
visible class is currently licensed outside that bundled obstruction reading.
\end{enumerate}
Hence the whole currently admitted run is read, at the present stage of the theory, by one
global pass/obstruction normal form and not by any third independent terminal morphology.
\end{theorem}

\begin{proof}
Theorem~\ref{thm:first-routed-return-normal-form} already classifies every currently admitted
routed return into exactly two bundled normal forms. On the positive side, the carrier-stable
routed return normal form means that the run has remained on the transferred carrier-stable
branch from the OP3/QAM closure package into the routed return side; consequently the whole
present run remains supported by one bundled closure-bearing carrier package, and all visible
OP5 diagnostics are subordinate to the carrier-stable visible normal form. This is precisely the
global pass normal form. On the negative side, the obstructed routed return normal form already
states that the first licensed failure is localized either upstream in the OP3/QAM obstruction
package or downstream at the carrier-upgraded OP5 obstruction layer; this is exactly the global
obstruction normal form. Because Theorem~\ref{thm:first-routed-return-normal-form} excludes a
third autonomous terminal visible class, the two global forms are exhaustive and mutually
exclusive.
\end{proof}

\begin{corollary}[Global pass/obstruction use rule]
\label{cor:global-pass-obstruction-use-rule}
In later proofs it is sufficient to cite Theorem~\ref{thm:first-global-pass-obstruction-normal-form}
once the current admissible run has entered the bundled OP3/QAM-to-OP5 channel. One may then
work only with the dichotomy \emph{global pass normal form versus global obstruction normal
form}.
\end{corollary}

\begin{proof}
This is the proof-usage reformulation of
Theorem~\ref{thm:first-global-pass-obstruction-normal-form}. The theorem already compresses the
currently admitted run from OP1 through OP5 into one reusable bundled normal-form dichotomy.
\end{proof}


\begin{remark}[Status of the bundled OP~3$\to$OP~5/global chain]
\label{rem:status-bundled-op3-op5-global-chain}
The present local-to-global chain should now also be read in two levels of proof strength. The \emph{active bundled interface} consists of the first OP~3$\to$OP~5 bundled transfer theorem, the compressed OP~3$\to$OP~5 proof interface, the routed return normal form, and the global pass/obstruction normal form together with their direct proof-use corollaries. These statements already give one self-contained current channel from translated OP~3 certification to a bundled global pass/obstruction decision. By contrast, the finer witness-grade hierarchy beneath the carrier-first OP~5 layer --- screening-loss normal form, carrier-witness upgrade, escalation thresholds, and the related witness-compression remarks --- remains a \emph{local refinement layer}: it must be reopened only when a later argument genuinely depends on witness species, witness grade, escalation, or screening-versus-carrier distinctions that are invisible at the bundled transfer/global level.
\end{remark}

\begin{proposition}[Citation and reopening rule for the bundled OP~3$\to$OP~5/global chain]
\label{prop:citation-reopening-rule-bundled-op3-op5-global}
Let a later argument concern a run already translated from the OP~3/QAM shell package into the routed OP~5 side. Then the following citation discipline is currently licensed.
\begin{enumerate}[leftmargin=2em,label=(\roman*)]
  \item If the argument needs only the bundled transfer/global dichotomy, it may cite Theorem~\ref{thm:first-op3-op5-bundled-transfer}, Proposition~\ref{prop:first-op3-op5-compressed-proof-interface}, Theorem~\ref{thm:first-routed-return-normal-form}, and Theorem~\ref{thm:first-global-pass-obstruction-normal-form} without reopening the finer witness hierarchy.
  \item If the argument further needs shell-level certificate or shell-level obstruction data, it must reopen the completed OP~3/QAM block rather than the entire older OP~1--OP~5 line.
  \item If the argument further needs witness grade, witness escalation, screening-witness compression, or carrier-witness separation, it must reopen the carrier-first OP~5 witness layer rather than treating the bundled global dichotomy as sufficient.
  \item No later proof may insert a third autonomous raw terminal visible branch after the bundled transfer/global interface has already been invoked.
\end{enumerate}
\end{proposition}

\begin{proof}
Clause~(i) is exactly the compressed reading already supplied by Theorem~\ref{thm:first-op3-op5-bundled-transfer}, Proposition~\ref{prop:first-op3-op5-compressed-proof-interface}, Theorem~\ref{thm:first-routed-return-normal-form}, and Theorem~\ref{thm:first-global-pass-obstruction-normal-form}. Clause~(ii) is the existing shell-sensitive fallback discipline of the completed OP~3/QAM block. Clause~(iii) is the corresponding witness-sensitive fallback discipline of the carrier-first OP~5 layer, where screening-loss normal form, upgrade, escalation, and separation are defined. Clause~(iv) is the bundled exclusion already built into the transfer theorem, the compressed proof interface, the routed return normal form, and the global pass/obstruction theorem: once those interfaces are active, no additional autonomous raw terminal visible branch is currently licensed.
\end{proof}

\begin{corollary}[Current sufficiency of the bundled OP~3$\to$OP~5/global chain]
\label{cor:current-sufficiency-bundled-op3-op5-global}
Assume a current argument starts from a translated OP~3 shell passage and asks only whether the run reaches a bundled pass reading or a bundled first licensed obstruction. Then the present document already contains a self-contained proof interface for that question: one may pass from the completed OP~3/QAM block through the bundled OP~3$\to$OP~5 transfer, the routed return normal form, and finally the global pass/obstruction normal form without introducing any additional intermediate theorem layer. In particular, the current line is already strong enough for the next public-shape step whenever no stronger claim about witness grade or deeper kernel-first activation is being asserted.
\end{corollary}

\begin{proof}
By Theorem~\ref{thm:qam-op3-proof-run} and Corollary~\ref{cor:qam-op3-proof-run-use}, the completed OP~3/QAM block already reduces translated shell data to certified passage versus first localized obstruction. Theorem~\ref{thm:first-op3-op5-bundled-transfer} and Proposition~\ref{prop:first-op3-op5-compressed-proof-interface} then transport that same dichotomy into the routed OP~5 side. Theorem~\ref{thm:first-routed-return-normal-form} rewrites the routed return as carrier-stable versus obstructed normal form, and Theorem~\ref{thm:first-global-pass-obstruction-normal-form} finally packages the whole current run into global pass versus global obstruction. Thus no extra intermediate theorem layer is needed for that bundled question. The final sentence is only the proof-theoretic reading of Remark~\ref{rem:status-bundled-op3-op5-global-chain}: stronger witness-grade or deeper kernel-first claims still require reopening the finer local layers.
\end{proof}


\subsection{Current extracted global pass/obstruction tool layer}
\label{sec:current-extracted-global-pass-obstruction-tool-layer}

\begin{remark}[Purpose of the current global tool layer]
The global pass/obstruction normal form is not meant to remain only as a terminal summary theorem. It is now extracted as a reusable tool layer so that later proofs can cite a small global interface instead of reopening the whole chain
\[
\text{OP1 source}
\to
\text{branch admissibility}
\to
\text{OP3/QAM closure}
\to
\text{carrier-first OP5 layer}.
\]
The extracted tools below are therefore not new theory beyond Theorem~\ref{thm:first-global-pass-obstruction-normal-form}; they are proof-usable projections of that bundled global dichotomy.
\end{remark}

\begin{proposition}[Global-pass tool]
\label{prop:global-pass-tool}
Assume Theorem~\ref{thm:first-global-pass-obstruction-normal-form}. If a current admissible run lies on the global pass normal form, then one may cite this fact directly as a reusable proof tool: the run is carried by one bundled closure-bearing carrier package from the OP1/branch entrance through the OP3/QAM shell package into the routed return side, and every admitted OP5 visible diagnostic remains subordinate to the carrier-stable visible normal form.
\end{proposition}

\begin{proof}
This is precisely the positive branch of Theorem~\ref{thm:first-global-pass-obstruction-normal-form}, rewritten as a direct reusable proof tool.
\end{proof}

\begin{proposition}[Global-obstruction tool]
\label{prop:global-obstruction-tool}
Assume Theorem~\ref{thm:first-global-pass-obstruction-normal-form}. If a current admissible run lies on the global obstruction normal form, then one may cite this fact directly as a reusable proof tool: the first licensed failure is already localized either upstream in the OP3/QAM obstruction package or downstream at the carrier-upgraded OP5 obstruction layer, and no third autonomous terminal visible class is currently licensed outside that bundled obstruction reading.
\end{proposition}

\begin{proof}
This is precisely the negative branch of Theorem~\ref{thm:first-global-pass-obstruction-normal-form}, rewritten as a direct reusable proof tool.
\end{proof}

\begin{proposition}[Global transfer tool]
\label{prop:global-transfer-tool}
Assume Theorem~\ref{thm:first-global-pass-obstruction-normal-form}. Any later proof that has already routed the current admissible run into the bundled OP3/QAM-to-OP5 channel may work only with the bundled transfer dichotomy
\[
\text{global pass}
\qquad\text{versus}\qquad
\text{global obstruction},
\]
without reopening the intermediate local decomposition into separate branch, shell, witness, escalation, and routed-return subcases.
\end{proposition}

\begin{proof}
By Corollary~\ref{cor:global-pass-obstruction-use-rule}, the whole currently admitted run has already been compressed into one reusable global dichotomy. The omitted intermediate subcases are already absorbed into the proof of Theorem~\ref{thm:first-global-pass-obstruction-normal-form}.
\end{proof}

\begin{proposition}[Global diagnostic subordination tool]
\label{prop:global-diagnostic-subordination-tool}
Assume Theorem~\ref{thm:first-global-pass-obstruction-normal-form}. On the global pass branch, every currently admitted visible routed-return diagnostic is globally subordinate to the same bundled closure-bearing carrier package. In particular, packet-wave, breathing-envelope, frame-flicker, and helix pictures cannot reintroduce a new structural branch once the run has already entered the global pass normal form.
\end{proposition}

\begin{proof}
Theorem~\ref{thm:first-global-pass-obstruction-normal-form} sends the positive branch through the carrier-stable routed return normal form and therefore through the carrier-stable visible normal form. By construction, all visible routed-return diagnostics are subordinate to that same bundled carrier package and cannot reopen the global dichotomy.
\end{proof}

\begin{proposition}[Global proof-interface tool]
\label{prop:global-proof-interface-tool}
Assume Theorem~\ref{thm:first-global-pass-obstruction-normal-form}. In later arguments it is sufficient to cite one of the four extracted global tools
\[
\text{global pass},\quad
\text{global obstruction},\quad
\text{global transfer},\quad
\text{global diagnostic subordination},
\]
according to the local need of the argument. No additional independent proof interface for the currently admitted run is required at this stage.
\end{proposition}

\begin{proof}
Theorem~\ref{thm:first-global-pass-obstruction-normal-form} already supplies the full current bundled dichotomy, while Propositions~\ref{prop:global-pass-tool}--\ref{prop:global-diagnostic-subordination-tool} expose its presently relevant projections. Hence these four extracted global tools form the current reusable proof interface for the admissible run.
\end{proof}

\begin{remark}[How to cite the extracted global tools]
The current global tool layer should be cited only after the run has already been shown to lie inside the bundled OP1$\to$branch$\to$OP3/QAM$\to$OP5 channel. Once that has been established, the present extracted tools allow later arguments to cite a compact global interface rather than reopening the entire local transport and diagnostic stack.
\end{remark}


\subsection{First global compressed proof run through extracted global tools}
\label{sec:first-global-compressed-proof-run-through-extracted-global-tools}

\begin{remark}[Reading key for the first global compressed proof run]
The present block is the first place where the current admissible run
\[
\text{OP1 source}
\to
\text{branch admissibility}
\to
\text{OP3/QAM closure}
\to
\text{carrier-first OP5 layer}
\]
is no longer reopened stage by stage. Instead, once the run has already been shown to lie inside the bundled global channel, the extracted global tools are used as one compressed proof interface. The proof run therefore has only two licensed bundled outcomes: a \emph{global pass} and a \emph{global obstruction}. The purpose of the theorem below is to state this compressed use rule explicitly, so that later arguments may pass through the whole current architecture in one step.
\end{remark}

\begin{theorem}[First global compressed proof run through extracted global tools]
\label{thm:first-global-compressed-proof-run-through-extracted-global-tools}
Assume that a current admissible run issued from the OP1 source side has already been routed into the bundled OP1$\to$branch$\to$OP3/QAM$\to$OP5 channel of the present architecture. Then the run may be processed through the extracted global tool layer without reopening the intermediate local stack. More precisely, exactly one of the following two bundled proof outputs is currently licensed.
\begin{enumerate}[leftmargin=2em,label=(\roman*)]
\item \textbf{Global compressed pass.}
By Proposition~\ref{prop:global-pass-tool}, the run lies on the global pass normal form and is supported by one bundled closure-bearing carrier package from the OP1/branch entrance through the OP3/QAM shell package into the routed return side. By Proposition~\ref{prop:global-diagnostic-subordination-tool}, all admitted visible routed-return diagnostics remain subordinate to that same bundled carrier package.
\item \textbf{Global compressed obstruction.}
By Proposition~\ref{prop:global-obstruction-tool}, the run lies on the global obstruction normal form and its first licensed failure is already localized either upstream in the OP3/QAM obstruction package or downstream at the carrier-upgraded OP5 obstruction layer. By Proposition~\ref{prop:global-transfer-tool}, no third autonomous terminal visible branch is currently licensed outside that bundled obstruction reading.
\end{enumerate}
Hence the present global proof run is now compressed to the dichotomy
\[
\text{global compressed pass}
\qquad\text{versus}\qquad
\text{global compressed obstruction},
\]
and no further independent proof interface for the currently admitted run is required once Proposition~\ref{prop:global-proof-interface-tool} applies.
\end{theorem}

\begin{proof}
This is the first direct use of the extracted global tool layer as a compressed proof interface. Once the current run has entered the bundled
\[
\text{OP1}\to\text{branch}\to\text{OP3/QAM}\to\text{OP5}
\]
channel, Proposition~\ref{prop:global-proof-interface-tool} states that no additional independent proof interface is currently required.

If the run lies on the positive side of Theorem~\ref{thm:first-global-pass-obstruction-normal-form}, Proposition~\ref{prop:global-pass-tool} supplies the bundled closure-bearing carrier package and Proposition~\ref{prop:global-diagnostic-subordination-tool} shows that all admitted visible OP5 diagnostics remain subordinate to that same carrier; this is exactly the global compressed pass.

If the run lies on the negative side, Proposition~\ref{prop:global-obstruction-tool} localizes the first licensed failure and Proposition~\ref{prop:global-transfer-tool} excludes any third independent terminal visible branch; this is exactly the global compressed obstruction.

Exhaustiveness and mutual exclusivity therefore come directly from the extracted global tools and ultimately from Theorem~\ref{thm:first-global-pass-obstruction-normal-form}.
\end{proof}

\begin{corollary}[First operational use of the extracted global tools]
\label{cor:first-operational-use-of-the-extracted-global-tools}
In later arguments it is sufficient to cite Theorem~\ref{thm:first-global-compressed-proof-run-through-extracted-global-tools} once the current admissible run has entered the bundled global channel. One may then work only with the dichotomy \emph{global compressed pass versus global compressed obstruction}, together with the induced carrier-stable subordination or obstruction localization, without reopening the entire local OP3/QAM-to-OP5 stack.
\end{corollary}

\begin{proof}
This is the operational reformulation of Theorem~\ref{thm:first-global-compressed-proof-run-through-extracted-global-tools}. The theorem already packages the extracted global tools into one reusable compressed proof run for the currently admitted architecture.
\end{proof}


\subsection{Global proof compression pass for the older OP1--OP5 line}
\label{sec:global-proof-compression-pass-for-the-older-op1-op5-line}

\begin{remark}[How the older OP1--OP5 line should now be reread]
The older connected OP1--OP5 transport line should no longer be read as a sequence of largely separate source-return, branch, shell, routed, and diagnostic arguments that must each be reopened in full. After the present global compressed proof run, the mathematically correct reading is instead this: once an argument has entered the currently admitted bundled channel, the older line should be read global-first through the dichotomy
\[
\text{global compressed pass}
\qquad\text{versus}\qquad
\text{global compressed obstruction}.
\]
Only if a later step genuinely depends on shell certificates, shell-obstruction resolution, witness grades, witness escalation, or other local carrier-diagnostic data should one reopen the OP3/QAM or OP5 layers. The purpose of the proposition below is not to replace the older theorems, but to state explicitly that they may now be reused through one compressed global interface.
\end{remark}

\begin{proposition}[Compressed rereading of the older OP1--OP5 line through the extracted global tools]
\label{prop:compressed-rereading-of-the-older-op1-op5-line-through-the-extracted-global-tools}
Assume that an argument along the older OP1--OP5 transport line has already been shown to lie inside the bundled channel of Theorem~\ref{thm:first-global-compressed-proof-run-through-extracted-global-tools}. Then the older line may be reread through the extracted global tools in the following compressed form.
\begin{enumerate}[leftmargin=2em,label=(\roman*)]
\item \textbf{Compressed pass rereading.}
If the argument lies on the positive side of Theorem~\ref{thm:first-global-compressed-proof-run-through-extracted-global-tools}, then the older OP1--OP5 line is to be read as one globally compressed pass: the source-return closure theorem, the admissible branch layer, the OP3/QAM closure package, the carrier-first OP5 interface, and the subordinated diagnostic stack are all already absorbed into one bundled carrier-stable run.
\item \textbf{Compressed obstruction rereading.}
If the argument lies on the negative side of Theorem~\ref{thm:first-global-compressed-proof-run-through-extracted-global-tools}, then the older OP1--OP5 line is to be read as one globally compressed obstruction: the first licensed failure is already localized in the bundled obstruction package, and no extra autonomous terminal visible branch may be added afterwards.
\end{enumerate}
In particular, later reuse of the older OP1--OP5 text should proceed through the extracted global proof interface rather than by reopening the whole local stack unless a genuinely new layer is being introduced.
\end{proposition}

\begin{proof}
This is the direct compression reading supplied by Theorem~\ref{thm:first-global-compressed-proof-run-through-extracted-global-tools}. On the positive side, the theorem already packages the full currently admitted run into the global compressed pass, and the extracted global tools show that the visible OP5 diagnostics remain subordinate to the same bundled carrier package. On the negative side, the same theorem packages the run into the global compressed obstruction and excludes any third autonomous terminal visible branch. Hence the older OP1--OP5 line can now be reused through one global proof interface rather than by reopening each intermediate theorem separately.
\end{proof}

\begin{corollary}[Global compressed reading of older routed return arguments]
\label{cor:global-compressed-reading-of-older-routed-return-arguments}
Whenever an older routed return argument is already known to lie inside the bundled current channel, it should by default be cited through Proposition~\ref{prop:compressed-rereading-of-the-older-op1-op5-line-through-the-extracted-global-tools}, Theorem~\ref{thm:first-global-compressed-proof-run-through-extracted-global-tools}, and the global citation discipline of Proposition~\ref{prop:first-global-tool-citation-guide}. In that situation, it is sufficient to work only with the dichotomy
\[
\text{global compressed pass}
\qquad\text{versus}\qquad
\text{global compressed obstruction},
\]
and with the induced bundled carrier-stable or obstruction-localized reading, unless a genuinely shell-sensitive or witness-sensitive local claim must still be reopened.
\end{corollary}

\begin{proof}
This is the operational reformulation of Proposition~\ref{prop:compressed-rereading-of-the-older-op1-op5-line-through-the-extracted-global-tools}. Once the argument is inside the bundled current channel, the extracted global proof interface supplies the compressed rereading rule for the older routed return line.
\end{proof}


\subsection{Witness-grade hardening and the first local kernel-first rewrite window}
\label{sec:witness-grade-hardening-and-first-local-kernel-first-rewrite-window}

\begin{remark}[Purpose of the v3.19.0 witness-grade hardening step]
The public v3.18.0 release already exposed the bundled OP~3/QAM closure package and the bundled OP~3$\to$OP~5/global channel in compact public shape. The next task is therefore not a new visible closure notion, but a stricter proof-use discipline for that same channel. In particular, the present step isolates three things: first, a bundled witness-grade reopening package for the OP~5 side; second, a sharper local/global reopening rule explaining exactly when the bundled global interface is sufficient; and third, one first narrow kernel-first rewrite window whose use is licensed only after the bundled channel has already been fixed and only at a local witness-grade transition where the deferred bridge actually reduces proof load.
\end{remark}

\begin{theorem}[Bundled witness-grade reopening package]
\label{thm:bundled-witness-grade-reopening-package}
Assume Theorem~\ref{thm:bundled-op5-diagnostic-theorem}, Proposition~\ref{prop:first-op3-op5-compressed-proof-interface}, Theorem~\ref{thm:first-global-pass-obstruction-normal-form}, Proposition~\ref{prop:first-global-tool-citation-guide}, and Proposition~\ref{prop:first-global-usage-checklist}. Let a later argument concern a run already known to lie inside the bundled current channel. Then exactly one of the following two proof-use regimes is presently licensed.
\begin{enumerate}[leftmargin=2em,label=(\roman*)]
\item \textbf{Bundled global-use regime.} If the argument needs only the bundled outputs global compressed pass/global compressed obstruction together with their carrier-stable or obstruction-localized consequences, then the extracted global interface is sufficient and no local witness-grade reopening is licensed or needed.
\item \textbf{Witness-grade reopening regime.} If the argument needs the witness grade itself, witness escalation, screening-witness compression, carrier-witness separation, diagnostic collapse below the escalation threshold, or a comparison between shell-grade and witness-grade data, then the bundled global interface must be reopened locally at the OP~5 witness layer, and the proof must work with the finer witness hierarchy rather than with the global dichotomy alone.
\end{enumerate}
Hence the current architecture already supports one bundled witness-grade reopening package: bundled-global use by default, local witness-grade reopening only for genuinely witness-sensitive claims.
\end{theorem}

\begin{proof}
The bundled global-use regime is exactly the content of Proposition~\ref{prop:first-global-tool-citation-guide} together with Proposition~\ref{prop:first-global-usage-checklist}: once the argument is already inside the bundled current channel and only the global pass/obstruction outputs are needed, the extracted global layer is the default proof interface. By contrast, Theorem~\ref{thm:bundled-op5-diagnostic-theorem} and Proposition~\ref{prop:first-op3-op5-compressed-proof-interface} show that witness grade, witness escalation, screening-witness compression, carrier-witness separation, and diagnostic collapse all live strictly below that bundled global interface. Therefore any genuinely witness-sensitive claim forces a local reopening of the OP~5 witness layer. The two regimes are exhaustive because every later proof step either remains at bundled-global level or asks for one of the finer witness-sensitive distinctions just listed.
\end{proof}

\begin{proposition}[Witness-grade local/global reopening discipline]
\label{prop:witness-grade-local-global-reopening-discipline}
Assume Theorem~\ref{thm:bundled-witness-grade-reopening-package}. Then the current bundled channel should be used with the following reopening discipline.
\begin{enumerate}[leftmargin=2em,label=(\roman*)]
\item work \emph{global-first} once the argument is already inside the bundled current channel;
\item reopen only the \emph{OP~3/QAM shell layer} when the step depends on shell certificates, shell-defect strata, shell commuting behavior, or OP~3-side carrier compression;
\item reopen only the \emph{OP~5 witness layer} when the step depends on witness grade, witness escalation, screening-witness compression, carrier-witness separation, or diagnostic collapse below escalation;
\item reopen \emph{both local layers} only when the step genuinely compares shell-grade and witness-grade data or transfers a shell-sensitive obstruction into a witness-sensitive one.
\end{enumerate}
In particular, visible OP~5 packet, envelope, flicker, and helix data do not by themselves justify a local reopening once the argument is already licensed to remain global-first.
\end{proposition}

\begin{proof}
Clause~(i) is exactly the bundled-global default from Theorem~\ref{thm:bundled-witness-grade-reopening-package}. Clause~(ii) restates the OP~3/QAM reopening part already fixed by Proposition~\ref{prop:first-global-tool-citation-guide}. Clause~(iii) is the witness-grade reopening regime from Theorem~\ref{thm:bundled-witness-grade-reopening-package}. Clause~(iv) records the only mixed local situation still licensed at the present stage: a step that must compare shell-grade and witness-grade structure rather than merely cite one local layer. The final sentence follows because, below escalation, visible OP~5 diagnostics remain subordinate to the carrier-stable visible normal form and therefore do not force reopening by appearance alone.
\end{proof}

\begin{corollary}[Witness-grade reopening checklist]
\label{cor:witness-grade-reopening-checklist}
For later proof use, the current bundled channel may be tested by the following checklist.
\begin{enumerate}[leftmargin=2em,label=(\roman*)]
\item Has the run already entered the bundled current channel?
\item Is only the bundled global pass/obstruction output needed?
\item If not, is the missing information shell-grade, witness-grade, or genuinely mixed shell/witness data?
\item Reopen only the corresponding local layer, and return to the bundled global interface immediately after the local claim has been discharged.
\end{enumerate}
This is the default proof-use checklist for the v3.19.0 witness-grade hardening step.
\end{corollary}

\begin{proof}
This is the operational reformulation of Proposition~\ref{prop:witness-grade-local-global-reopening-discipline}. The checklist simply turns the local/global reopening rule into a direct proof-use procedure.
\end{proof}

\begin{definition}[First licensed local kernel-first rewrite window]
\label{def:first-licensed-local-kernel-first-rewrite-window}
Assume the direct kernel constructor of Definition~\ref{def:qam-direct-kernel-constructor}, the bridge conditions of Theorem~\ref{thm:qam-direct-kernel-constructor}, and the bundled witness-grade reopening package of Theorem~\ref{thm:bundled-witness-grade-reopening-package}. The \emph{first licensed local kernel-first rewrite window} is the following strictly local situation: a proof step has already entered the bundled current channel, has already been forced to reopen at witness grade, and now needs only to compare source-side and return-side witness-bearing states through one common carrier-stable normal form without promoting any broader kernel-first migration on the surrounding global channel.
\end{definition}

\begin{proposition}[First licensed local kernel-first rewrite window]
\label{prop:first-licensed-local-kernel-first-rewrite-window}
Assume Definition~\ref{def:first-licensed-local-kernel-first-rewrite-window}, Proposition~\ref{prop:carrier-first-op5-interface}, and Theorem~\ref{thm:bundled-op5-diagnostic-theorem}. Let a witness-sensitive proof step concern a routed return realization already known to remain on the carrier-stable branch of the bundled OP~3$\to$OP~5 transfer channel. Then that step may be rewritten kernel-first in the following narrow sense.
\begin{enumerate}[leftmargin=2em,label=(\roman*)]
\item the visible witness-bearing source/return surfaces may be replaced by their exposed payload kernels and compared only through the common kernel normal form forced by the carrier-stable branch;
\item the local readout comparison may be performed at the level of the resulting common kernel/readout carrier class instead of on the full visible witness-bearing surfaces;
\item after that local comparison has been discharged, the proof must return immediately to the bundled carrier-stable/global interface and may not infer a broader kernel-first rewrite rule for the full OP~3$\to$OP~5/global channel.
\end{enumerate}
Hence the present line now contains one first licensed kernel-first rewrite window: local witness-grade carrier comparison under an already fixed carrier-stable branch, but no general kernel-first migration of the whole bundled channel.
\end{proposition}

\begin{proof}
Once the proof step is already known to lie on the carrier-stable branch, the witness-bearing visible data are subordinate to one common closure-bearing carrier package. Under the direct kernel constructor and the bridge conditions, this is exactly the situation in which the local comparison load can be shifted from full visible surfaces to their exposed payload kernels and then to the corresponding common kernel normal form. The local readout comparison may therefore be discharged at the level of the common kernel/readout carrier class. However, Theorem~\ref{thm:bundled-witness-grade-reopening-package} and Definition~\ref{def:first-licensed-local-kernel-first-rewrite-window} keep this move strictly local: the argument was already forced into witness-grade reopening, and nothing in the present hypotheses licenses promotion of that local kernel-first step to a global rewrite rule for the surrounding bundled channel. The final sentence is exactly this restriction written as a proof-use statement.
\end{proof}

\begin{corollary}[Use rule for the first local kernel-first rewrite window]
\label{cor:use-rule-for-the-first-local-kernel-first-rewrite-window}
Whenever a later proof step is already inside the bundled current channel, already forced to reopen at witness grade, and already known to remain on the carrier-stable branch, it is sufficient to cite Proposition~\ref{prop:first-licensed-local-kernel-first-rewrite-window} for the local source/return comparison. After that local comparison, the proof should return to the bundled global interface immediately.
\end{corollary}

\begin{proof}
This is the operational reformulation of Proposition~\ref{prop:first-licensed-local-kernel-first-rewrite-window}. The proposition already states both the local license and the immediate return rule.
\end{proof}


\begin{remark}[Status of the first corridor-level activation step]
\label{rem:status-of-the-first-corridor-level-activation-step}
The present release line still does \emph{not} activate a kernel-first rewrite rule for the whole bundled OP~3$\to$OP~5/global channel. Nevertheless, the combination of the bundled witness-grade reopening package and the first licensed local kernel-first rewrite window already supports one slightly stronger activation step: a finite witness-grade corridor on one fixed carrier-stable branch may be discharged kernel-first through one common carrier-stable kernel/readout class. This is the first point at which the current release line moves beyond a single isolated local rewrite without yet promoting the deferred kernel architecture to an active global migration layer.
\end{remark}

\begin{definition}[Carrier-stable witness corridor]
\label{def:carrier-stable-witness-corridor}
Assume Theorem~\ref{thm:bundled-witness-grade-reopening-package} and Proposition~\ref{prop:first-licensed-local-kernel-first-rewrite-window}. A \emph{carrier-stable witness corridor} is a finite family of witness-sensitive proof steps
\[
\Sigma_1,\ldots,\Sigma_N
\]
along one routed return run such that:
\begin{enumerate}[leftmargin=2em,label=(\roman*)]
\item the run is already known to lie inside the bundled current channel;
\item every step $\Sigma_j$ is already forced to reopen at witness grade and remains on the same carrier-stable branch;
\item for every $j$, the local source/return comparison occurring in $\Sigma_j$ satisfies the hypotheses of Proposition~\ref{prop:first-licensed-local-kernel-first-rewrite-window};
\item no step in the family changes the fixed carrier-stable normal-form class or requires a new shell/witness transfer beyond that same carrier-stable branch.
\end{enumerate}
Thus a carrier-stable witness corridor is the smallest finite witness-grade subrun on which the same carrier-stable kernel/readout class may be reused without reopening the whole bundled channel anew at each local comparison.
\end{definition}

\begin{lemma}[First corridor-level kernel-first activation lemma]
\label{lem:first-corridor-level-kernel-first-activation}
Assume Definition~\ref{def:carrier-stable-witness-corridor}. Then the whole witness corridor may be discharged kernel-first through one common carrier-stable kernel/readout class in the following precise sense.
\begin{enumerate}[leftmargin=2em,label=(\roman*)]
\item for every local step $\Sigma_j$, the visible witness-bearing source/return surfaces may be replaced by their exposed payload kernels and compared only through the same carrier-stable kernel normal-form class;
\item the local readout comparisons occurring across the full finite family $\Sigma_1,\ldots,\Sigma_N$ may therefore be performed at class level rather than re-established on the full visible witness-bearing surfaces step by step;
\item after the corridor has been discharged, the proof must return to the bundled global interface immediately and may not infer a general kernel-first rewrite rule for branch-changing, shell-changing, or globally mixed witness-grade arguments.
\end{enumerate}
Equivalently, the present line already licenses one first corridor-level activation step: a finite bundled family of witness-grade local comparisons along one fixed carrier-stable branch, but still no general kernel-first migration of the whole bundled channel.
\end{lemma}

\begin{proof}
By Definition~\ref{def:carrier-stable-witness-corridor}, every local step $\Sigma_j$ already satisfies the hypotheses of Proposition~\ref{prop:first-licensed-local-kernel-first-rewrite-window}. Hence each such step may be rewritten kernel-first at the local witness-sensitive level. Because all steps lie on the same carrier-stable branch and no step changes the fixed carrier-stable normal-form class, the corresponding local kernel/readout comparisons all factor through one and the same carrier-stable kernel/readout class rather than through a new class at every step. Therefore the whole finite family may be discharged kernel-first corridor by corridor instead of reopening the full visible witness-bearing surfaces separately at each local comparison.

The return rule is unchanged: Theorem~\ref{thm:bundled-witness-grade-reopening-package} keeps witness-grade reopening subordinate to the bundled global interface, and Definition~\ref{def:carrier-stable-witness-corridor} excludes branch-changing, shell-changing, or globally mixed witness-grade transitions. Consequently, once the corridor has been discharged, the proof must return immediately to the bundled global interface, and no broader global kernel-first migration follows from this corridor-level argument.
\end{proof}

\begin{corollary}[Use rule for the first corridor-level activation lemma]
\label{cor:use-rule-for-the-first-corridor-level-activation-lemma}
Whenever a later proof contains a finite witness-grade subrun on one fixed carrier-stable branch, it is sufficient to cite Lemma~\ref{lem:first-corridor-level-kernel-first-activation} instead of reopening each witness-bearing source/return comparison separately. After that corridor-level comparison has been discharged, the proof should return to the bundled global interface immediately.
\end{corollary}

\begin{proof}
This is the operational reformulation of Lemma~\ref{lem:first-corridor-level-kernel-first-activation}. The lemma already states both the corridor-level license and the immediate return rule.
\end{proof}




\begin{remark}[Status of the first corridor-to-bundle propagation step]
\label{rem:status-of-the-first-corridor-to-bundle-propagation-step}
The present release line still does \emph{not} activate a kernel-first rewrite rule for the whole bundled OP~3$\to$OP~5/global channel. Nevertheless, the combination of the bundled witness-grade reopening package, the first licensed local kernel-first rewrite window, and the first corridor-level activation lemma already supports one slightly stronger propagation step: a bundled run may now alternate between globally readable stretches and finitely many carrier-stable witness corridors without forcing each witness-bearing local comparison to be reopened from scratch. This is the first point at which the current release line moves beyond one isolated corridor-level activation and begins to control a whole bundled run by a finite alternation of global citations and corridor-level kernel-first discharges.
\end{remark}

\begin{definition}[Bundle-compatible witness-corridor decomposition]
\label{def:bundle-compatible-witness-corridor-decomposition}
Assume Theorem~\ref{thm:bundled-witness-grade-reopening-package}, Lemma~\ref{lem:first-corridor-level-kernel-first-activation}, and Proposition~\ref{prop:first-global-tool-citation-guide}. A \emph{bundle-compatible witness-corridor decomposition} of a later proof run is a finite alternating decomposition
\[
\Gamma_0,\ \mathcal C_1,\ \Gamma_1,\ \ldots,\ \mathcal C_M,\ \Gamma_M
\]
such that:
\begin{enumerate}[leftmargin=2em,label=(\roman*)]
\item the entire run is already known to lie inside the bundled current channel;
\item each global segment $\Gamma_i$ needs only the bundled outputs global compressed pass/global compressed obstruction together with their induced carrier-stable or obstruction-localized consequences;
\item each local segment $\mathcal C_k$ is a carrier-stable witness corridor in the sense of Definition~\ref{def:carrier-stable-witness-corridor};
\item after each local corridor $\mathcal C_k$, the proof returns immediately to the bundled global interface before any later corridor begins;
\item no corridor changes the fixed bundled run into a new shell-changing, branch-changing, or globally mixed witness-grade run between two global segments.
\end{enumerate}
Thus a bundle-compatible witness-corridor decomposition is the smallest finite alternation in which the bundled global interface remains the default proof surface and only finitely many carrier-stable witness corridors are discharged locally.
\end{definition}

\begin{proposition}[First corridor-to-bundle propagation principle]
\label{prop:first-corridor-to-bundle-propagation-principle}
Assume Definition~\ref{def:bundle-compatible-witness-corridor-decomposition}. Then the whole decomposed run may be discharged by the following finite proof pattern.
\begin{enumerate}[leftmargin=2em,label=(\roman*)]
\item on each global segment $\Gamma_i$, it is sufficient to cite the bundled global interface and work only with the dichotomy
\[
\text{global compressed pass}
\quad\text{versus}\quad
\text{global compressed obstruction};
\]
\item on each local corridor $\mathcal C_k$, it is sufficient to cite Lemma~\ref{lem:first-corridor-level-kernel-first-activation} and discharge the whole finite witness-grade corridor through one common carrier-stable kernel/readout class;
\item after each local corridor $\mathcal C_k$, the proof returns immediately to the bundled global interface and resumes global-first reasoning on the next segment $\Gamma_k$;
\item consequently, no individual witness-bearing source/return comparison inside the corridors needs to be reopened separately, and no general kernel-first rewrite rule for the full bundled channel is inferred.
\end{enumerate}
Equivalently, the present line already licenses one first corridor-to-bundle propagation rule: a bundled run may be handled by finitely many global-interface citations plus finitely many corridor-level kernel-first discharges, while the active proof surface remains globally bundled outside those local corridors.
\end{proposition}

\begin{proof}
By Definition~\ref{def:bundle-compatible-witness-corridor-decomposition}, every global segment $\Gamma_i$ satisfies exactly the hypotheses of the bundled global citation discipline, so Proposition~\ref{prop:first-global-tool-citation-guide} supplies the global compressed pass/global compressed obstruction interface there. Likewise, every local segment $\mathcal C_k$ is a carrier-stable witness corridor, so Lemma~\ref{lem:first-corridor-level-kernel-first-activation} discharges that whole finite corridor kernel-first through one common carrier-stable kernel/readout class.

The return rule is part of both ingredients: the global interface is the default outside local reopenings, while the corridor lemma requires immediate return to the bundled global interface after each local corridor has been discharged. Hence the full run is controlled by a finite alternation of global citations and corridor-level kernel-first discharges. Because the corridor lemma already absorbs the local witness-grade source/return comparisons inside each $\mathcal C_k$, no individual witness-bearing comparison has to be reopened separately. Finally, nothing in the hypotheses licenses promotion of this finite alternation to a general kernel-first rewrite rule for the whole bundled channel; outside the listed corridors the active proof surface remains the bundled global interface.
\end{proof}

\begin{corollary}[Use rule for the first corridor-to-bundle propagation principle]
\label{cor:use-rule-for-the-first-corridor-to-bundle-propagation-principle}
Whenever a later proof run already inside the bundled current channel decomposes into globally readable stretches separated by finitely many carrier-stable witness corridors, it is sufficient to cite Proposition~\ref{prop:first-corridor-to-bundle-propagation-principle} together with Lemma~\ref{lem:first-corridor-level-kernel-first-activation}, rather than reopening each witness-bearing local comparison separately. Only the corridor boundaries, not the internal witness steps of each corridor, remain visible at proof-use level.
\end{corollary}

\begin{proof}
This is the operational reformulation of Proposition~\ref{prop:first-corridor-to-bundle-propagation-principle}. The proposition already states both the finite propagation rule and the immediate return to the bundled global interface after each corridor.
\end{proof}


\begin{remark}[Status of the first bundle-level reopening compression step]
\label{rem:status-of-the-first-bundle-level-reopening-compression-step}
The present release line still does \emph{not} collapse every local reopening beneath the bundled channel to one uniform active theorem surface. What it now does support, however, is one further compression of proof-use data: once a bundled run has already been decomposed into global stretches and finitely many carrier-stable witness corridors, the later proof no longer needs to describe each corridor by all of its internal shell-grade and witness-grade subclaims. It is enough to record, corridor by corridor, which local reopening profile is actually being used. This is the first point at which the current release line compresses not only the internal witness comparisons of each corridor, but also the reopening metadata attached to the whole bundled run.
\end{remark}

\begin{definition}[Profile-stable bundle-level reopening data]
\label{def:profile-stable-bundle-level-reopening-data}
Assume Definition~\ref{def:bundle-compatible-witness-corridor-decomposition}. A \emph{profile-stable bundle-level reopening datum} on such a decomposition is a choice of one profile label
\[
\pi_k\in\{\mathrm{Q},\mathrm{W},\mathrm{QW}\}
\]
for every corridor $\mathcal C_k$ such that:
\begin{enumerate}[leftmargin=2em,label=(\roman*)]
\item $\pi_k=\mathrm{Q}$ means that the corridor is reopened only through the OP~3/QAM shell side: shell certificates, shell-obstruction stratification, shell commuting behavior, or OP~3-side carrier compression;
\item $\pi_k=\mathrm{W}$ means that the corridor is reopened only through the carrier-first OP~5 witness side: witness grade, witness escalation, screening-witness compression, carrier-witness separation, or diagnostic collapse beneath the escalation threshold;
\item $\pi_k=\mathrm{QW}$ means that the corridor genuinely compares shell-grade and witness-grade data inside one common carrier-stable witness corridor;
\item the profile label of a corridor remains fixed throughout that corridor and changes, if at all, only at corridor boundaries.
\end{enumerate}
Thus the profile sequence $(\pi_1,\ldots,\pi_M)$ records the smallest local reopening data needed at proof-use level beyond the bundled global interface.
\end{definition}

\begin{proposition}[First bundle-level reopening compression principle]
\label{prop:first-bundle-level-reopening-compression-principle}
Assume a bundle-compatible witness-corridor decomposition together with profile-stable bundle-level reopening data in the sense of Definition~\ref{def:profile-stable-bundle-level-reopening-data}. Then the whole bundled run may be handled at proof-use level by the finite data
\[
\Gamma_0,\ \pi_1,\ \Gamma_1,\ \ldots,\ \pi_M,\ \Gamma_M
\]
with the following force.
\begin{enumerate}[leftmargin=2em,label=(\roman*)]
\item Each global segment $\Gamma_i$ is still treated only through the bundled global interface and the dichotomy global compressed pass versus global compressed obstruction.
\item Each corridor $\mathcal C_k$ is reopened only through its single profile label $\pi_k$, together with Lemma~\ref{lem:first-corridor-level-kernel-first-activation} whenever the corridor is discharged kernel-first as one carrier-stable witness corridor.
\item No later proof needs to list the internal shell-grade or witness-grade subclaims of $\mathcal C_k$ separately as long as they remain subordinate to the same fixed corridor profile $\pi_k$.
\item If a corridor would force a profile change inside itself, then the present compression rule stops at that point and the corridor must be split into smaller profile-stable corridors before the bundled proof-use compression is applied again.
\end{enumerate}
Equivalently, the current release line already licenses one first bundle-level reopening compression step: after corridor-level kernel-first discharge has been fixed, the remaining local proof-use burden of a bundled run is reduced to a finite sequence of profile labels attached to finitely many carrier-stable corridors.
\end{proposition}

\begin{proof}
By Proposition~\ref{prop:first-corridor-to-bundle-propagation-principle}, the bundled run is already controlled by finitely many global citations together with finitely many corridor-level kernel-first discharges. The only remaining local proof-use burden is therefore to specify which finer local reopening regime is actually needed on each corridor. Definition~\ref{def:profile-stable-bundle-level-reopening-data} makes exactly that finite burden explicit: every corridor carries one fixed profile label $\pi_k$, and that label determines whether the local reopening is OP~3/QAM-only, OP~5-witness-only, or genuinely mixed.

Because the profile remains fixed throughout each corridor, all internal shell-grade and witness-grade subclaims of that corridor stay subordinate to the same local reopening regime and need not be listed one by one at proof-use level. If a profile change occurred inside a corridor, then the corridor would no longer be profile-stable, and the current compression would no longer be legitimate without first splitting that corridor into smaller profile-stable pieces. This proves both the finite proof-use description and the stated stopping rule.
\end{proof}

\begin{corollary}[Use rule for the first bundle-level reopening compression principle]
\label{cor:use-rule-for-the-first-bundle-level-reopening-compression-principle}
Whenever a later proof run already inside the bundled current channel admits both a bundle-compatible witness-corridor decomposition and profile-stable bundle-level reopening data, it is sufficient to cite Proposition~\ref{prop:first-bundle-level-reopening-compression-principle} and record only the resulting finite profile sequence $(\pi_1,\ldots,\pi_M)$, rather than restating all local shell-grade and witness-grade subclaims corridor by corridor. A finer reopening is required only when one corridor changes profile internally or leaves the carrier-stable witness-corridor regime.
\end{corollary}

\begin{proof}
This is the operational reformulation of Proposition~\ref{prop:first-bundle-level-reopening-compression-principle}. The proposition already states both the compressed proof-use description and the precise point where that compression stops.
\end{proof}




\begin{remark}[Status of the first profile-to-run compression step]
\label{rem:status-of-the-first-profile-to-run-compression-step}
The present release line still does \emph{not} identify all bundled runs by one fully kernel-activated normal form. What it now does support, however, is one further compression of the proof-use surface: once a bundled run has already been decomposed into global stretches and finitely many profile-stable carrier corridors, the later proof no longer needs to remember the names of the individual global segments or the internal witness subclaims of those corridors. It is enough to remember the alternating \emph{run profile word} that records only when the proof is global and which local reopening profile occurs at each corridor. This is the first point at which the current release line compresses a whole bundled run to one finite symbolic proof-use word.
\end{remark}

\begin{definition}[Run-stable bundled proof word]
\label{def:run-stable-bundled-proof-word}
Assume a bundle-compatible witness-corridor decomposition together with profile-stable bundle-level reopening data in the sense of Definition~\ref{def:profile-stable-bundle-level-reopening-data}. The associated \emph{run-stable bundled proof word} is the finite alternating word
\[
\omega(\Gamma):=\mathrm{G}\,\pi_1\,\mathrm{G}\,\pi_2\,\mathrm{G}\cdots \pi_M\,\mathrm{G}
\]
over the alphabet
\[
\{\mathrm{G},\mathrm{Q},\mathrm{W},\mathrm{QW}\},
\]
where:
\begin{enumerate}[leftmargin=2em,label=(\roman*)]
\item each letter $\mathrm{G}$ records one globally readable bundled segment handled only through the dichotomy global compressed pass versus global compressed obstruction;
\item each letter $\pi_k\in\{\mathrm{Q},\mathrm{W},\mathrm{QW}\}$ records the fixed reopening profile of the corresponding carrier-stable witness corridor $\mathcal C_k$;
\item the word begins and ends with $\mathrm{G}$, and the letters alternate globally/local in the same order as the bundled run;
\item no letter change is allowed inside a corridor; if such a change occurs, the current word is not defined until the corridor has been split into smaller profile-stable corridors.
\end{enumerate}
Thus $\omega(\Gamma)$ is the smallest finite symbolic record that still remembers where the proof is global and which local reopening profile occurs whenever the bundled channel is reopened.
\end{definition}

\begin{proposition}[First profile-to-run compression principle]
\label{prop:first-profile-to-run-compression-principle}
Assume Definition~\ref{def:run-stable-bundled-proof-word}. Then the proof-use burden of the whole bundled run is reduced to its finite run-stable bundled proof word
\[
\omega(\Gamma)\in\{\mathrm{G},\mathrm{Q},\mathrm{W},\mathrm{QW}\}^{*}
\]
with the following force.
\begin{enumerate}[leftmargin=2em,label=(\roman*)]
\item Every occurrence of $\mathrm{G}$ means that the corresponding segment is handled only through the bundled global interface and its compressed pass/obstruction dichotomy.
\item Every occurrence of $\mathrm{Q}$, $\mathrm{W}$, or $\mathrm{QW}$ means that the next corridor is reopened only in that single local profile, together with Lemma~\ref{lem:first-corridor-level-kernel-first-activation} whenever the corridor is licensed for carrier-stable kernel-first discharge.
\item No later proof needs to retain the individual names of the global segments $\Gamma_i$ or the separate internal shell-grade/witness-grade clauses of the corridors once the alternating word $\omega(\Gamma)$ has been fixed.
\item If one wants to refine a corridor internally, change its local reopening profile midstream, or distinguish two globally readable segments by more than the bundled pass/obstruction interface, then the present compression stops and the run must first be rewritten at the preceding bundle-level stage.
\end{enumerate}
Equivalently, the current release line already licenses one first profile-to-run compression step: after bundle-compatible corridor decomposition and profile stabilization have been fixed, the whole bundled proof surface may be cited by one finite alternating word over $\{\mathrm{G},\mathrm{Q},\mathrm{W},\mathrm{QW}\}$.
\end{proposition}

\begin{proof}
By Proposition~\ref{prop:first-bundle-level-reopening-compression-principle}, the bundled run is already controlled by finitely many global segments together with finitely many profile-stable carrier corridors. The only remaining proof-use information is therefore the order in which those two types of pieces occur and, on each corridor, which fixed reopening profile is active. Definition~\ref{def:run-stable-bundled-proof-word} records exactly that data as the alternating word $\omega(\Gamma)$.

The force of each letter is already supplied by the previous stages: $\mathrm{G}$ stands for the bundled global interface, while $\mathrm{Q}$, $\mathrm{W}$, and $\mathrm{QW}$ stand for the corresponding local reopening profiles and, where licensed, for the corridor-level kernel-first discharge of Lemma~\ref{lem:first-corridor-level-kernel-first-activation}. Since the previous bundle-level compression already removes the need to list individual internal witness clauses corridor by corridor, fixing the alternating word removes the need to retain the names of the global segments as separate proof-use data. If one wants finer information than the word provides, then one has already left the current compressed regime and must reopen the run at the preceding bundle-level stage. This proves the stated stopping rule.
\end{proof}

\begin{corollary}[Use rule for the first profile-to-run compression principle]
\label{cor:use-rule-for-the-first-profile-to-run-compression-principle}
Whenever a later proof run already inside the bundled current channel admits a run-stable bundled proof word $\omega(\Gamma)$, it is sufficient to cite Proposition~\ref{prop:first-profile-to-run-compression-principle} and record only that finite word, rather than listing the separate global segments and corridor-internal local claims again. A finer reopening is required only when a corridor changes profile internally, leaves the carrier-stable witness-corridor regime, or one global segment must be distinguished from another by more than the bundled pass/obstruction interface.
\end{corollary}

\begin{proof}
This is the operational reformulation of Proposition~\ref{prop:first-profile-to-run-compression-principle}. The proposition already states both the finite run-word description and the precise point where that description stops being legitimate.
\end{proof}


\begin{remark}[Status of the first word-determined run-class step]
\label{rem:status-of-the-first-word-determined-run-class-step}
The present release line still does \emph{not} identify all bundled runs by one globally activated kernel normal form, nor does it erase the distinction between different geometric realizations of the same bundled channel. What it now does support, however, is one further proof-use quotient: once a bundled run has already been compressed to a run-stable bundled proof word, the later proof may forget the individual realization of that run and retain only the \emph{run class determined by that word}. This is the first point at which the current release line passes from one finite symbolic proof-use word to one finite symbolic proof-use class.
\end{remark}

\begin{definition}[Word-determined bundled run class]
\label{def:word-determined-bundled-run-class}
Assume Definition~\ref{def:run-stable-bundled-proof-word}. Two bundled runs $\Gamma$ and $\Gamma'$ are called \emph{word-determined proof-use equivalent} if both run-stable bundled proof words are defined and satisfy
\[
\omega(\Gamma)=\omega(\Gamma').
\]
The corresponding equivalence class is written
\[
[\Gamma]_{\omega}:=\{\Gamma' : \omega(\Gamma')=\omega(\Gamma)\}.
\]
It is called the \emph{word-determined bundled run class} of $\Gamma$. In this reading, the class $[\Gamma]_{\omega}$ records exactly the proof-use information retained by the current compressed line and forgets the names of individual global segments, the particular realization of each carrier-stable witness corridor, and all corridor-internal shell-grade or witness-grade subclaims beyond what is already encoded by the letters of $\omega(\Gamma)$.
\end{definition}

\begin{proposition}[First word-to-run-class compression principle]
\label{prop:first-word-to-run-class-compression-principle}
Assume Definition~\ref{def:word-determined-bundled-run-class}. Then the current proof-use burden of a bundled run depends only on its word-determined bundled run class $[\Gamma]_{\omega}$ and no longer on the individual run $\Gamma$ itself, with the following force.
\begin{enumerate}[leftmargin=2em,label=(\roman*)]
\item If two bundled runs $\Gamma$ and $\Gamma'$ satisfy $\omega(\Gamma)=\omega(\Gamma')$, then every later argument that uses only the bundled global interface, the associated local reopening profiles, and the corresponding corridor-level kernel-first licenses may treat $\Gamma$ and $\Gamma'$ as the same proof-use object.
\item Consequently, once the run-stable bundled proof word has been fixed, the later proof may cite the class $[\Gamma]_{\omega}$ instead of the particular run realization $\Gamma$.
\item The present compression stops exactly when one needs information not encoded by the word: for example, a distinction between two globally readable segments carrying the same letter $\mathrm G$, a distinction between two carrier-stable witness corridors carrying the same profile letter, or a geometric/property-level comparison internal to one corridor that is invisible at word level.
\end{enumerate}
Equivalently, the current release line already licenses one first word-to-run-class compression step: after passage from bundled runs to run-stable bundled proof words, proof use depends only on the induced equivalence class under equality of those words.
\end{proposition}

\begin{proof}
By Proposition~\ref{prop:first-profile-to-run-compression-principle}, the current proof-use surface of a bundled run is already reduced to the finite alternating word $\omega(\Gamma)$. Hence every later argument that cites only the force carried by the letters $\mathrm G$, $\mathrm Q$, $\mathrm W$, and $\mathrm{QW}$, together with the associated corridor-level kernel-first license where available, uses no further data from the particular realization of $\Gamma$.

Therefore, if $\omega(\Gamma)=\omega(\Gamma')$, then both runs carry exactly the same compressed proof-use data and may be treated as the same proof-use object at the current level of abstraction. This proves (i) and the resulting class-level citation claim in (ii). Clause (iii) is simply the contrapositive reading of the same compression: whenever a later argument needs data not already encoded by the word, the quotient by equality of words is too coarse, so the proof must reopen at the preceding run-level or bundle-level stage.
\end{proof}

\begin{corollary}[Use rule for the first word-to-run-class compression principle]
\label{cor:use-rule-for-the-first-word-to-run-class-compression-principle}
Whenever a later proof already inside the bundled current channel depends only on the run-stable bundled proof word and the force of its letters, it is sufficient to cite Proposition~\ref{prop:first-word-to-run-class-compression-principle} and work with the corresponding class $[\Gamma]_{\omega}$, rather than with one chosen realization of the bundled run. A finer reopening is required only when the argument needs information not encoded by the word itself.
\end{corollary}

\begin{proof}
This is the operational reformulation of Proposition~\ref{prop:first-word-to-run-class-compression-principle}. The proposition already states both the induced class-level proof-use rule and the precise point where that compression stops.
\end{proof}




\begin{remark}[Status of the bundled proof-use compression stack in the present v\docversion\ release line]
\label{rem:status-of-the-first-run-class-concatenation-step}
The present release line still does \emph{not} equip bundled current-channel reasoning with one unrestricted global kernel algebra. What it now supports, however, is a much shorter proof-use stack: bundled runs may be compressed successively from run words to run classes, from finite concatenation chains to one reduced normal-form word, from that word to one proof-use class, from binary and finite class composition to one associative class word, and finally to one contraction-stable terminal citation surface. The next items record only the exact proof-use interfaces needed later; the longer stage-by-stage derivations are no longer repeated here.
\end{remark}

\begin{definition}[Concatenable word-determined bundled run classes]
\label{def:concatenable-word-determined-bundled-run-classes}
Assume Definition~\ref{def:word-determined-bundled-run-class}. Two word-determined bundled run classes $[\Gamma]_{\omega}$ and $[\Gamma']_{\omega}$ are \emph{concatenable at bundled level} if their join is handled entirely by the bundled global interface and no new local reopening is created at the splice. In that case define the reduced concatenated run word by
\[
\omega(\Gamma)\star_{\mathrm G}\omega(\Gamma')
\]
obtained by identifying the terminal global letter $\mathrm G$ of the first word with the initial global letter $\mathrm G$ of the second, and write
\[
[\Gamma]_{\omega}\star[\Gamma']_{\omega}
\]
for the induced concatenated run-class citation object.
\end{definition}

\begin{proposition}[First run-class concatenation principle]
\label{prop:first-run-class-concatenation-principle}
Under Definition~\ref{def:concatenable-word-determined-bundled-run-classes}, bundled-level concatenation depends only on the two input run classes and on the reduced concatenated run word. In particular, later proof use may cite $[\Gamma]_{\omega}\star[\Gamma']_{\omega}$ instead of chosen representatives, and this compression stops exactly when the splice creates local data not already encoded by the adjacent global-interface letters.
\end{proposition}

\begin{proof}
Immediate from the word-level quotient already fixed by Proposition~\ref{prop:first-word-to-run-class-compression-principle}: once representatives are forgotten, the only remaining bundled splice datum is the reduced global-interface identification. Any genuinely new local burden at the join lies beyond this quotient and forces reopening.
\end{proof}

\begin{corollary}[Use rule for the first run-class concatenation principle]
\label{cor:use-rule-for-the-first-run-class-concatenation-principle}
Whenever two bundled current-channel runs are already compressed to word-determined bundled run classes and their join is handled entirely by the bundled global interface, it is sufficient to cite Proposition~\ref{prop:first-run-class-concatenation-principle} and work with the concatenated class $[\Gamma]_{\omega}\star[\Gamma']_{\omega}$.
\end{corollary}

\begin{proof}
This is the operational reading of Proposition~\ref{prop:first-run-class-concatenation-principle}.
\end{proof}

\begin{remark}[Status of the first finite run-class normal-form step]
\label{rem:status-of-the-first-finite-run-class-normal-form-step}
The next compression replaces a finite ordered chain of concatenable bundled run classes by one reduced normal-form word. Only the ordered local reopening profile and the bundled global separators survive.
\end{remark}

\begin{definition}[Finite concatenable bundled run-class chain and reduced normal-form word]
\label{def:finite-concatenable-bundled-run-class-chain-and-reduced-normal-form-word}
Let
\[
[\Gamma_1]_{\omega},\dots,[\Gamma_n]_{\omega}
\]
be bundled run classes such that each adjacent pair is concatenable at bundled level. Define the reduced normal-form concatenation word
\[
\operatorname{NF}_{\star}([\Gamma_1]_{\omega},\dots,[\Gamma_n]_{\omega})
\]
by concatenating the words $\omega(\Gamma_k)$ and identifying every adjacent terminal/initial global letter $\mathrm G$ as one common bundled global-interface letter.
\end{definition}

\begin{proposition}[First finite run-class normal-form principle]
\label{prop:first-finite-run-class-normal-form-principle}
For every finite bundled concatenation chain covered by Definition~\ref{def:finite-concatenable-bundled-run-class-chain-and-reduced-normal-form-word}, the proof-use burden depends only on the ordered list of run classes and on the reduced normal-form word $\operatorname{NF}_{\star}$. Representative choices and parenthesization of successive concatenations do not matter.
\end{proposition}

\begin{proof}
Repeated use of Proposition~\ref{prop:first-run-class-concatenation-principle} shows that each bundled-level splice is governed only by the reduced global-interface identification. Iterating that splice therefore yields one parenthesization-independent reduced word.
\end{proof}

\begin{corollary}[Use rule for the first finite run-class normal-form principle]
\label{cor:use-rule-for-the-first-finite-run-class-normal-form-principle}
Whenever a later argument uses only the ordered local profile carried by a finite concatenable run-class chain, it is sufficient to cite Proposition~\ref{prop:first-finite-run-class-normal-form-principle} and work with $\operatorname{NF}_{\star}$.
\end{corollary}

\begin{proof}
This is the operational reading of Proposition~\ref{prop:first-finite-run-class-normal-form-principle}.
\end{proof}

\begin{remark}[Status of the first normal-form-determined bundled proof-use class step]
\label{rem:status-of-the-first-normal-form-determined-bundled-proof-use-class-step}
Once a finite concatenable bundled run-class chain has been reduced to one normal-form word, later proof use may forget the chosen chain realization and retain only the class determined by that word.
\end{remark}

\begin{definition}[Normal-form-determined bundled proof-use class]
\label{def:normal-form-determined-bundled-proof-use-class}
Two finite concatenable bundled run-class chains are called normal-form proof-use equivalent if their reduced normal-form words agree. The corresponding equivalence class is the \emph{normal-form-determined bundled proof-use class} and is written $[\mathfrak C]_{\operatorname{NF}}$.
\end{definition}

\begin{proposition}[First normal-form-to-proof-use-class compression principle]
\label{prop:first-normal-form-to-proof-use-class-compression-principle}
The current proof-use burden of a finite concatenable bundled run-class chain depends only on its normal-form-determined bundled proof-use class. Equivalently, once only the reduced normal-form word and the force of its letters matter, one may replace the chain by $[\mathfrak C]_{\operatorname{NF}}$.
\end{proposition}

\begin{proof}
This is precisely the quotient induced by equality of reduced normal-form words from Definition~\ref{def:normal-form-determined-bundled-proof-use-class}.
\end{proof}

\begin{corollary}[Use rule for the first normal-form-to-proof-use-class compression principle]
\label{cor:use-rule-for-the-first-normal-form-to-proof-use-class-compression-principle}
Whenever a later proof already inside the bundled current channel depends only on the reduced normal-form concatenation word and the force of its letters, it is sufficient to cite Proposition~\ref{prop:first-normal-form-to-proof-use-class-compression-principle} and work with the corresponding class $[\mathfrak C]_{\operatorname{NF}}$.
\end{corollary}

\begin{proof}
This is the operational reading of Proposition~\ref{prop:first-normal-form-to-proof-use-class-compression-principle}.
\end{proof}

\begin{remark}[Status of the first proof-use-class composition step]
\label{rem:status-of-the-first-proof-use-class-composition-step}
After passage to normal-form-determined classes, bundled composition can be handled at class level whenever the splice itself adds no new local burden.
\end{remark}

\begin{definition}[Composable normal-form-determined bundled proof-use classes]
\label{def:composable-normal-form-determined-bundled-proof-use-classes}
Two normal-form-determined bundled proof-use classes are \emph{composable} if realizations may be joined entirely through the bundled global interface with no new local reopening at the splice. Their reduced composed normal-form word is obtained by the same $\star_{\mathrm G}$ splice on the two reduced normal-form words.
\end{definition}

\begin{proposition}[First proof-use-class composition principle]
\label{prop:first-proof-use-class-composition-principle}
Whenever two normal-form-determined bundled proof-use classes are composable, their bundled composition depends only on the two classes and on the reduced composed normal-form word. Representative-level choices are invisible at this class-compression stage.
\end{proposition}

\begin{proof}
Combine Proposition~\ref{prop:first-normal-form-to-proof-use-class-compression-principle} with the same bundled-global splice rule already used at run-class level.
\end{proof}

\begin{corollary}[Use rule for the first proof-use-class composition principle]
\label{cor:use-rule-for-the-first-proof-use-class-composition-principle}
Whenever two bundled current-channel proof packages are already compressed to normal-form-determined bundled proof-use classes and their splice is handled entirely by the bundled global interface, it is sufficient to cite Proposition~\ref{prop:first-proof-use-class-composition-principle} and work with their composed class.
\end{corollary}

\begin{proof}
This is the operational reading of Proposition~\ref{prop:first-proof-use-class-composition-principle}.
\end{proof}

\begin{remark}[Status of the first associative proof-use-class normal-form step]
\label{rem:status-of-the-first-associative-proof-use-class-normal-form-step}
Iterated class composition can now be read through one parenthesization-stable reduced class word.
\end{remark}

\begin{definition}[Associative reduced class-composition normal form]
\label{def:associative-reduced-class-composition-normal-form}
For a finite ordered composable family of normal-form-determined bundled proof-use classes, define the associative reduced class-composition normal form by concatenating the reduced class words and identifying every interior terminal/initial global letter $\mathrm G$ as one common bundled global-interface letter.
\end{definition}

\begin{proposition}[First associative proof-use-class normal-form principle]
\label{prop:first-associative-proof-use-class-normal-form-principle}
The proof-use burden of any finite ordered composable family of normal-form-determined bundled proof-use classes depends only on its associative reduced class-composition normal form and is independent of parenthesization.
\end{proposition}

\begin{proof}
At word level the bundled-global splice is associative; Proposition~\ref{prop:first-proof-use-class-composition-principle} removes dependence on representatives, so only the resulting reduced class word remains.
\end{proof}

\begin{corollary}[Use rule for the first associative proof-use-class normal-form principle]
\label{cor:use-rule-for-the-first-associative-proof-use-class-normal-form-principle}
Whenever a later bundled current-channel proof uses only the force of a finite ordered composable family of normal-form-determined bundled proof-use classes and all intermediate joins are handled entirely by the bundled global interface, it is sufficient to cite Proposition~\ref{prop:first-associative-proof-use-class-normal-form-principle} and work with the single associative class word.
\end{corollary}

\begin{proof}
This is the operational reading of Proposition~\ref{prop:first-associative-proof-use-class-normal-form-principle}.
\end{proof}

\begin{remark}[Status of the first contraction-stable class-normal-form step]
\label{rem:status-of-the-first-contraction-stable-class-normal-form-step}
The final local compression removes only redundant consecutive global separators and leaves the genuinely reopened local content unchanged.
\end{remark}

\begin{definition}[Contraction-stable reduced class-normal-form word]
\label{def:contraction-stable-reduced-class-normal-form-word}
Given a finite composable family of normal-form-determined bundled proof-use classes with associative reduced class word $\omega_{\mathrm{assoc}}$, define
\[
\operatorname{NF}_{\mathrm{ctr}}
\]
by replacing each maximal consecutive block of the symbol $G$ in $\omega_{\mathrm{assoc}}$ by a single occurrence of $G$.
\end{definition}

\begin{proposition}[First contraction-stable class-normal-form principle]
\label{prop:first-contraction-stable-class-normal-form-principle}
The contraction-stable class-normal-form word $\operatorname{NF}_{\mathrm{ctr}}$ is independent of representatives and parenthesization and determines the same proof-use surface as the associative class-normal-form word, except that redundant purely global separators are contracted.
\end{proposition}

\begin{proof}
Representative and parenthesization independence are inherited from Proposition~\ref{prop:first-associative-proof-use-class-normal-form-principle}; the contraction step touches only maximal $G$-blocks and therefore does not alter the order of the genuinely local letters $Q$, $W$, and $QW$.
\end{proof}

\begin{corollary}[Use rule for the first contraction-stable class-normal-form principle]
\label{cor:use-rule-for-the-first-contraction-stable-class-normal-form-principle}
Once a finite bundled proof run has already been compressed to composable proof-use classes and the only residual redundancy lies in repeated consecutive global separators, the run may be cited through the single contraction-stable class-normal-form word $\operatorname{NF}_{\mathrm{ctr}}$.
\end{corollary}

\begin{proof}
This is the operational reading of Proposition~\ref{prop:first-contraction-stable-class-normal-form-principle}.
\end{proof}

\begin{remark}[Status of the terminal proof-use surface in the present v\docversion\ release line]
\label{rem:status-of-the-terminal-proof-use-surface-in-the-present-v319-line}
For the present v\docversion\ release line, the bundled proof-use compression stack ends at the contraction-stable class-normal-form surface: after this point later proof use needs only the surviving letters $Q$, $W$, or $QW$ as possible local reopening sites.
\end{remark}

\begin{proposition}[Terminal citation rule for the present compressed proof-use line]
\label{prop:terminal-citation-rule-for-the-present-v319-compression-line}
Let a bundled current-channel argument already be compressed to a finite composable family of normal-form-determined bundled proof-use classes, with associative class normal form formed and residual global redundancy removed by the contraction-stable rule. Then the argument may be cited through the single object $\operatorname{NF}_{\mathrm{ctr}}$ alone. A finer reopening is licensed only when the later step depends on local data surviving at one of the symbols $Q$, $W$, or $QW$.
\end{proposition}

\begin{proof}
This is the terminal reading of Proposition~\ref{prop:first-contraction-stable-class-normal-form-principle}: once redundant global separators are removed, no further class-level proof-use data remain beyond the surviving local letters.
\end{proof}

\begin{corollary}[Operational mini-checklist for the terminal proof-use surface]
\label{cor:operational-mini-checklist-for-the-terminal-proof-use-surface}
For the present v\docversion\ release line, the compressed workflow is
\[
\text{bundled run}
\Longrightarrow \text{proof-use classes}
\Longrightarrow \text{associative class normal form}
\Longrightarrow \operatorname{NF}_{\mathrm{ctr}}.
\]
If this terminal object already suffices for the claim at hand, no earlier layer should be reopened.
\end{corollary}

\begin{proof}
This is the direct operational reading of Proposition~\ref{prop:terminal-citation-rule-for-the-present-v319-compression-line}.
\end{proof}

\subsection{First global tool citation guide}
\label{sec:first-global-tool-citation-guide}

\begin{center}
\fbox{\begin{minipage}{0.93\linewidth}
\textbf{Global-first mini-box.} Use the bundled global interface by default once the argument is already inside the current OP1$\to$branch$\to$OP3/QAM$\to$OP5 channel. Work only with \emph{global compressed pass} versus \emph{global compressed obstruction} unless a step explicitly needs finer local data.

\textbf{Reopen OP3/QAM only if} the step depends on shell certificates, shell-level obstruction strata, shell commuting behavior, or OP3-side carrier compression.

\textbf{Reopen OP5 only if} the step depends on witness grade, witness escalation, screening-witness compression, carrier-witness separation, or diagnostic collapse below the escalation threshold.

\textbf{Reopen both local layers only if} the step genuinely compares shell-grade and witness-grade data.
\end{minipage}}
\end{center}

\begin{remark}[When the global tool layer may be cited directly]
The extracted global tool layer is intended to be the default citation interface for later arguments only after three conditions have already been established: first, the run has already entered the bundled OP1$\to$branch$\to$OP3/QAM$\to$OP5 channel; second, the argument no longer depends on a specific local shell, witness, packet, or helix species; and third, only the bundled dichotomy
\[
\text{global compressed pass}
\qquad\text{versus}\qquad
\text{global compressed obstruction}
\]
needs to be used. The purpose of the guide below is to state explicitly when the present global interface is sufficient and when one must reopen the local OP3/QAM or OP5 layers.
\end{remark}

\begin{proposition}[First global tool citation guide]
\label{prop:first-global-tool-citation-guide}
Assume that a later argument concerns a run already known to lie inside the bundled current channel of Theorem~\ref{thm:first-global-compressed-proof-run-through-extracted-global-tools}. Then the following citation discipline is currently licensed.
\begin{enumerate}[leftmargin=2em,label=(\roman*)]
\item \textbf{Global citation regime.}
It is sufficient to cite the extracted global tool layer alone whenever the argument needs only the bundled global outputs
\[
\text{global compressed pass}
\qquad\text{or}\qquad
\text{global compressed obstruction},
\]
together with their induced transfer and diagnostic-subordination consequences. In that situation one may work only through Proposition~\ref{prop:global-proof-interface-tool}, Theorem~\ref{thm:first-global-compressed-proof-run-through-extracted-global-tools}, Proposition~\ref{prop:compressed-rereading-of-the-older-op1-op5-line-through-the-extracted-global-tools}, and the extracted global tools cited there.
\item \textbf{OP3/QAM reopening regime.}
One must reopen the local OP3/QAM layer whenever the argument depends on a specifically shell-sensitive claim, including shell certificates, shell defect witnesses, the shell commuting-square, shell-obstruction stratification, carrier compression on the OP3 side, or any need to identify the first local obstruction already before the bundled global dichotomy is applied.
\item \textbf{OP5 diagnostic reopening regime.}
One must reopen the local carrier-first OP5 layer whenever the argument depends on the witness grade itself, on witness escalation, on screening-witness compression, on carrier-witness separation, on diagnostic collapse below the escalation threshold, or on any distinction between screening-only and carrier-upgraded visible mismatches.
\end{enumerate}
Hence the extracted global tool layer is the default citation interface for bundled arguments, but it is not meant to erase local OP3/QAM or OP5 structure whenever a later proof genuinely needs the finer shell or witness resolution.
\end{proposition}

\begin{proof}
The global citation regime follows from Proposition~\ref{prop:global-proof-interface-tool}, Theorem~\ref{thm:first-global-compressed-proof-run-through-extracted-global-tools}, and Proposition~\ref{prop:compressed-rereading-of-the-older-op1-op5-line-through-the-extracted-global-tools}: once the argument is already inside the bundled current channel and only the global pass/obstruction dichotomy is needed, the extracted global tools already supply the full current proof interface. By contrast, every claim listed in the OP3/QAM reopening regime depends on shell certificates, shell-level obstruction localization, or shell-level carrier compression, and therefore cannot be recovered from the global compressed dichotomy alone. Likewise, every claim listed in the OP5 reopening regime depends on the internal witness hierarchy beneath the carrier-first OP5 interface and therefore requires reopening that finer diagnostic layer. The stated citation discipline is therefore exactly the current boundary between the bundled global interface and the local shell/witness interfaces.
\end{proof}

\begin{corollary}[Practical citation rule for the current global tools]
\label{cor:practical-citation-rule-for-the-current-global-tools}
In later proofs one should cite the extracted global tool layer by default once the argument is already inside the bundled current channel. The local OP3/QAM or OP5 blocks need to be reopened only if the proof truly depends on shell-level certificates or obstruction localization, or on witness-grade and escalation data that are invisible at the global compressed level.
\end{corollary}

\begin{proof}
This is the operational restatement of Proposition~\ref{prop:first-global-tool-citation-guide}. The proposition already separates the current bundled default interface from the finer local shell and witness interfaces.
\end{proof}

\subsection{First global usage checklist}
\label{sec:first-global-usage-checklist}

\begin{remark}[Three practical questions before reopening the local stack]
The global-first discipline becomes easiest to use if later proofs ask three short questions in the same order. First: is the argument already inside the bundled current channel? Second: is the proof using only the dichotomy \emph{global compressed pass versus global compressed obstruction}? Third: does any step depend on shell-grade or witness-grade data that the bundled global interface does not retain? The checklist below turns those questions into a practical default rule.
\end{remark}

\begin{proposition}[First global usage checklist]
\label{prop:first-global-usage-checklist}
Assume a later proof step concerns a run in the current OP1$\to$branch$\to$OP3/QAM$\to$OP5 architecture. Then the following checklist is currently licensed.
\begin{enumerate}[leftmargin=2em,label=(\arabic*)]
\item \textbf{Enter the bundled channel first.}
Before using the global interface, verify that the proof step already lies in the bundled channel of Theorem~\ref{thm:first-global-compressed-proof-run-through-extracted-global-tools}. If not, the global checklist is not yet available.
\item \textbf{Stay global if the dichotomy is enough.}
If the proof step needs only the outputs \emph{global compressed pass} or \emph{global compressed obstruction}, together with bundled transfer or diagnostic-subordination consequences, cite Proposition~\ref{prop:first-global-tool-citation-guide}, Proposition~\ref{prop:global-proof-interface-tool}, and Theorem~\ref{thm:first-global-compressed-proof-run-through-extracted-global-tools}, and do not reopen the local stack.
\item \textbf{Reopen OP3/QAM only for shell-sensitive data.}
If the proof must identify shell certificates, shell-level obstruction strata, shell commuting behavior, or OP3-side carrier compression before the global dichotomy is applied, reopen the OP3/QAM layer and not the entire OP1--OP5 chain.
\item \textbf{Reopen OP5 only for witness-sensitive data.}
If the proof must distinguish screening-only from carrier-upgraded visible witnesses, or must use witness escalation, witness compression, or diagnostic collapse below the escalation threshold, reopen the carrier-first OP5 layer and not the entire older transport line.
\item \textbf{Reopen both layers only when the proof genuinely compares them.}
Only when a step explicitly couples shell-grade and witness-grade information should both local layers be reopened together. Otherwise, the bundled global interface remains the correct default proof surface.
\end{enumerate}
\end{proposition}

\begin{proof}
Item~(1) is the entry condition from Theorem~\ref{thm:first-global-compressed-proof-run-through-extracted-global-tools}. Item~(2) is exactly the bundled citation regime of Proposition~\ref{prop:first-global-tool-citation-guide}. Items~(3) and~(4) unpack the two licensed local reopening regimes from that same proposition into a practical proof checklist. Item~(5) follows because the global interface was introduced precisely to avoid reopening both local stacks unless a proof step truly depends on their finer interaction.
\end{proof}

\begin{corollary}[Practical default rule: global first, local only if needed]
\label{cor:practical-default-rule-global-first-local-only-if-needed}
In the current architecture, later proofs should proceed by default through the bundled global interface. Local reopening is exceptional and should be limited to shell-sensitive OP3/QAM claims, witness-sensitive OP5 claims, or genuinely mixed shell/witness comparisons.
\end{corollary}

\begin{proof}
This is the direct operational summary of Proposition~\ref{prop:first-global-usage-checklist}. The checklist turns the present global-first discipline into a short reusable proof rule.
\end{proof}


\subsection[The neutrino k-value]{The neutrino \texorpdfstring{$k$}{k}-value}

\begin{proposition}[$k(\nu) = 4$ as hypothesis, hypothesis-grade]
\label{prop:k-neutrino}
\emph{(Hypothesis-grade; requires derivation from the fermionic sector.)}
The unusually small neutrino masses may be explained by $k(\nu) = 4$:
\[
  m_\nu^{(\mathrm{matter})} = \gL^4 \cdot m_\nu^{(\mathrm{photon})},
  \qquad \gL^4 \approx 12.1.
\]
In the photon frame, neutrino masses would be smaller by the factor
$\gL^{-3} \approx 1/6.5$ than one would expect from the charged-lepton to
neutrino-mass ratio measured in the matter frame.
Structurally, $k(\nu) = 4$ corresponds to
$k(m_e)+k(G_N)+k(\ell_P^{-1}) = 1+2+3/2$ --- a lepton--graviton--Planck
vertex, consistent with the LQG coupling in the $\Hplusminus$ sector
(Theorem~\ref{thm:op2-selfadj}).
A structural derivation from the Dirac operator $D_-$ on $\Hplusminus$
is still outstanding.
\end{proposition}

\subsection{Hubble tension as cosmological evolution of $\bsym$}

\begin{proposition}[H$_0$ tension as $\bsym(t)$ evolution, hypothesis-grade]
\label{prop:H0-tension}
\emph{(Hypothesis-grade; requires cosmological model.)}
The Hubble tension ($H_0^{\mathrm{CMB}} \approx 67.4$ vs.
$H_0^{\mathrm{local}} \approx 73.0$ km/s/Mpc, factor $\approx 1.083$)
may reflect a cosmological evolution of the frame offset $\bsym(t)$:
\begin{align*}
  \bsym(z=0) &\approx 0.5493 \quad \text{(today)}, \\
  \bsym(z=1100) &\approx 0.527 \quad \text{(CMB epoch)}.
\end{align*}
The required change $\Delta\bsym \approx +0.023$ over 13.8 Gyr
corresponds to a rate of $\approx 1.6\times 10^{-12}$ per year.
Physically, as matter structures in the universe grew ($z: 1100\to 0$),
the frame offset $\bsym$ also increased --- the universe moves
continuously away from the photon frame.
In this picture, the $H_0$ discrepancy would require no new physics,
but would instead be an artefact of the time variation of $\bsym(t)$.
\end{proposition}

\begin{remark}[Evaluation of the hypotheses]
\label{rem:inflow-evaluation}
Among the three hypotheses, the structurally most robust one is the
$\bsym(t)$ evolution (Prop.~\ref{prop:H0-tension}): it explains the
$H_0$ tension without free parameters (only
$\beta_{\mathrm{sym},0} = 0.5493$ and
$\beta_{\mathrm{sym},\mathrm{CMB}}$ determined from the tension).
The $k(\nu) = 4$ hypothesis (Prop.~\ref{prop:k-neutrino}) is
structurally compatible with the GOE class of $D_-$ on $\Hplusminus$
(Theorem~\ref{thm:op2-selfadj}), but still needs an explicit derivation.
The inflow picture (Prop.~\ref{prop:inflow-cascade}) is already contained
in the existing framework.
\end{remark}

\subsection{Muon anomaly and the $k$-value of the strong coupling}
\label{sec:muon-anomaly}

\begin{proposition}[Muon $g-2$ as hypothesis, hypothesis-grade]
\label{prop:muon-g2}
\emph{(Hypothesis-grade; the source value $k=1.2$ is treated here as a provisional phenomenological coefficient rather than a structurally derived constant.)}

The anomalous magnetic moment of the muon shows a discrepancy of
\[
  \delta a_\mu
  := a_\mu^{(\mathrm{exp})} - a_\mu^{(\mathrm{SM})}
  \approx 2.51\times 10^{-9},
\]
which corresponds to about $3.6\%$ of the hadronic vacuum-polarization
contribution $a_\mu^{(\mathrm{hvp})} \approx 6.9\times 10^{-8}$.

The framework offers a possible explanation through the unknown
$k$-values of the strong coupling $\alpha_s$ and of the QCD scale
$\Lambda_{\mathrm{QCD}}$:
\begin{itemize}
  \item If $k(\alpha_s) = 1$ (like the electromagnetic coupling $\alpha$),
    then $\alpha_s^{(\mathrm{photon})} = \gL^{-1}\cdot\alpha_s^{(\mathrm{matter})}$,
    which scales the hadronic contribution by $\gL^{-1}\approx 0.537$.
  \item If $k(\Lambda_{\mathrm{QCD}}) \in [0,1]$, then the correction factor
    for $a_\mu^{(\mathrm{hvp})}$ lies in the range $[\gL^0,\gL^{-2}]$.
\end{itemize}
The $k$-value pair $(k(\alpha_s), k(\Lambda_{\mathrm{QCD}}))$ that
reproduces $\delta a_\mu$ is an open question in the framework's QCD
sector. The value $k=1.2$ from the source document is a phenomenological
placeholder, not a derived value.
\end{proposition}

\begin{proposition}[Saturation model, hypothesis-grade]
\label{prop:saturation}
\emph{(Hypothesis-grade; cosmological end-state.)}

The saturation model connects $k(\nu)=4$
(Proposition~\ref{prop:k-neutrino}) with long-term cosmology:
\begin{enumerate}[label=(\arabic*)]
  \item As expansion increases, neutrino masses increase in the matter frame
    ($m_\nu^{(\mathrm{matter})} = \gL^4(t)\cdot m_\nu^{(\mathrm{photon})}$
    when $\gL(t)$ grows, i.e. $\bsym(t)\nearrow 1$).
  \item When $\bsym(t)\to 1$ (tachyon limit), one has
    $\gL\to\infty$, $m_\nu^{(\mathrm{matter})}\to\infty$, and the cascade
    probability $P(S^8\to S^4) = \bsym^2\to 1$ --- the universe becomes
    fully ``matter dominated.''
  \item Shortly before the limit, $H_0\to 0$ and the expansion stops ---
    the universe reaches the ``tachyon horizon'' $\zminus = -\bsym$
    (framework definition limits; cf.
    Proposition~\ref{prop:horizon-cascade-singularity}).
\end{enumerate}
This model is not a new assumption, but the consistent continuation of
the $\bsym(t)$ evolution (Proposition~\ref{prop:H0-tension}) into the
distant future.
\end{proposition}

\begin{remark}[$k$-value table for the open QCD questions]
\label{rem:k-qcd}
The following table summarizes the $k$-value status of the QCD sector:
\begin{quote}\small
\noindent\textbf{Reader cue.} Read the next table as a compact status ledger. Its purpose is to separate established assignments from genuinely open QCD entries, so the discussion that follows can orient itself without adding a new proof step.
\end{quote}
\begin{center}
\renewcommand{\arraystretch}{1.3}
\begin{tabular}{llll}
\toprule
Constant & Value & $k$ & Status \\
\midrule
$\alpha$ (EM)    & $1/137.036$ & $1$ & established \\
$G_N$ (gravitation) & $6.674\times10^{-11}$ & $2$ & established \\
$\gBI$ (LQG)     & $0.2375$    & $0$ & established \\
$\alpha_s$ (QCD strong) & $0.118$ & $?$ & \textbf{open} \\
$\Lambda_{\mathrm{QCD}}$ & $\approx 200$ MeV & $?$ & \textbf{open} \\
$m_\nu$ (neutrino)  & $< 0.04$ eV & $4$? & hypothesis \\
\bottomrule
\end{tabular}
\end{center}
Determining $k(\alpha_s)$ and $k(\Lambda_{\mathrm{QCD}})$ would explain
the muon $g$-2 anomaly, the running coupling in the QCD sector, and
possibly CP violation within the framework.
\end{remark}





\subsection{OP 4: CKM/PMNS from \texorpdfstring{$G_2$}{G2}
  Euler angles — closing the factor of 2}
\label{sec:op4}
%% ============================================================

\subsubsection{What is already established}

From the Dixon algebra structure (Theorem~11.3, Companion):
\begin{itemize}
\item The three fermion generations are $\mathrm{Gen}_i \cong
  (\mathbb{C}\otimes\mathbb{H})\otimes e_{i+3}$ for $i = 1,2,3$.
\item Generation-changing transitions are $G_2$-rotations in
  $\mathrm{Im}(\mathbb{O})$: $e_4 \cdot e_5 = +e_7$,
  $e_5 \cdot e_6 = +e_1$, $e_6 \cdot e_4 = -e_3$.
\item A natural candidate for the Cabibbo angle is
  \[
    \theta_{12} \approx \arcsin(\gamma_{\mathrm{string}})
    \approx \arcsin\!\Bigl(\frac{\ln 2}{\pi\sqrt{3}}\Bigr)
    \approx 7.3^\circ.
  \]
\item Observed value: $\theta_C \approx 13.04^\circ$.
  Ratio: $13.04/7.3 \approx 1.79 \approx 2$.
\end{itemize}

\subsubsection{The source of the factor of 2}

\begin{proposition}[The factor 2 is a Hopf-projection doubling]
\label{prop:factor2}
The $G_2$-Euler-angle derivation of the CKM angles uses the
\emph{bare} $G_2$ rotation angle on $S^7 \to S^{15} \to S^8$.
However, the physical mixing angle is the Hopf-projected angle
onto the $SU(3) \subset G_2$ subgroup. The Hopf projection
$S^7 \to S^4$ introduces a factor of $2$ in the relation between
the $G_2$ Euler angle and the physical CKM entry:
\[
  \theta_{\mathrm{CKM}}
  = 2 \cdot \theta_{G_2}^{\mathrm{bare}}
  + \mathcal{O}(\delta_{\mathrm{CL}}).
\]
\end{proposition}

\begin{proof}
On $S^3 \cong SU(2)$, the Hopf map $\pi: S^3 \to S^2$ has the
property that a rotation by angle $\varphi$ in the fibre (i.e.\ in
the $U(1)$ direction) produces a rotation by $2\varphi$ in the base
$S^2$. The same doubling occurs in the higher Hopf fibration
$S^7 \to S^{15} \to S^8$: a $G_2$-rotation by $\theta$ in the
total space $S^{15}$ (which carries the full spinor structure)
projects to an $SU(3)$ rotation by $2\theta$ in the base $S^8$.
The CKM matrix entries measure angles in the $SU(3)$ base space
(quark flavor space), not in the $S^{15}$ total space.
Hence $\theta_{12}^{\mathrm{CKM}} = 2\arcsin(\gamma_{\mathrm{string}})
\approx 2 \times 7.3^\circ = 14.6^\circ$, within $12\%$ of the
observed $13.04^\circ$.
\end{proof}

\begin{remark}[Residual 12\% after the factor-2 correction]
After applying the factor-of-2 correction:
$2\arcsin(\gamma_{\mathrm{string}}) \approx 14.6^\circ$
vs.\ observed $13.04^\circ$. The residual gap is $12\%$,
which is consistent with the $\delta_{\mathrm{CL}}$ correction
from OP 1:
\[
  \theta_{12}^{\mathrm{pred}}
  = 2\arcsin(\gamma_{\mathrm{string}}) \cdot \delta_{\mathrm{CL}}
  \approx 14.6^\circ \times 0.9515
  \approx 13.9^\circ.
\]
This is now within $6.5\%$ of the observed value.
The remaining gap requires the full lattice-to-continuum correction
(same $\delta_{\mathrm{CL}}$ as in OP 1).
\end{remark}

\subsubsection{The PMNS matrix and the triolic angle}

\begin{proposition}[PMNS angles from $\mathbb{Z}/3\mathbb{Z}$ symmetry]
\label{prop:pmns-z3}
The large mixing angles of the PMNS matrix (neutrino mixing) differ
from CKM precisely because the lepton sector sits at the
\emph{triolic fixed points} of $\mathbb{Z}/3\mathbb{Z}$ on $S^3$
(Theorem~11.4, Koide formula), while the quark sector sits at
$G_2$-rotated points. The triolic fixed-point angle gives
\[
  \theta_{23}^{\mathrm{PMNS}} \approx \arctan(\sqrt{2}) \approx 54.74^\circ
  \quad (\text{the maximal mixing angle}),
\]
from the same tetrahedral geometry as the Hopf candidate of OP 1.
The atmospheric mixing angle $\theta_{23} \approx 45$--$50^\circ$
is the $\delta_{\mathrm{CL}}$-shifted version of this.
\end{proposition}

\subsubsection{Unified prediction table}

\begin{quote}\small
\noindent\textbf{Reader cue.} Read the next table as a comparison surface. Its purpose is to place the bare geometric values, the doubling step, the continuum-limit correction, and the observed numbers into one scanable line of sight, not to add an independent derivation.
\end{quote}

\begin{center}
\renewcommand{\arraystretch}{1.4}
\begin{tabular}{lllll}
\toprule
Observable & Bare $G_2$/Hopf & $\times 2$ & $\times\delta_{\mathrm{CL}}$ & Observed \\
\midrule
$\theta_{12}^{\mathrm{CKM}}$ (Cabibbo)
  & $7.3^\circ$ & $14.6^\circ$ & $13.9^\circ$ & $13.04^\circ$ \\
$\theta_{23}^{\mathrm{PMNS}}$ (atmospheric)
  & $54.74^\circ$ & — & $52.1^\circ$ & $\approx 49^\circ$ \\
$\bsym$ & $1/\sqrt{3} \approx 0.5774$ & — & $0.5492$ & $0.5493$ \\
$\gamma_L$ & $2/\sqrt{3} \approx 1.155$ & --- & $1.100$ & not fixed here \\
\midrule
\multicolumn{5}{@{}p{0.88\textwidth}@{}}{\footnotesize Note: $\delta_{\mathrm{CL}} \approx 0.9515$; the $\gamma_L$ line is heuristic only and is not used as a defining identity in the present text.}\\
\bottomrule
\end{tabular}
\end{center}

\subsubsection{The remaining open step}

\begin{remark}[What is still needed]
The structural argument (Proposition~\ref{prop:factor2}) identifies
the factor-of-2 as Hopf-projection doubling. What remains:
\begin{enumerate}
\item[(a)] A \emph{formal proof} that the $G_2$-Euler angles on
  $S^{15} \to S^8$ project to CKM angles via the exact formula
  $\theta_{\mathrm{CKM}} = 2\theta_{G_2}$ (not just a sketch).
  This requires computing the derivative of the Hopf projection map
  at the relevant fiber and showing the Jacobian is exactly 2.
\item[(b)] A derivation of $\delta_{\mathrm{CL}}$ from first
  principles (which is the same as OP 1). Once $\delta_{\mathrm{CL}}$
  is computed analytically, the full CKM and PMNS matrices follow
  from the Hopf-projected $G_2$ Euler angles.
\item[(c)] The CP-violating phase $\delta_{CP}$ in CKM. This
  corresponds to the phase of the $G_2$ rotation in the Fano-plane
  structure of $\mathrm{Im}(\mathbb{O})$. The natural candidate is
  $\delta_{CP} \approx \arg(e_4 \cdot e_5 \cdot e_6) = \pi/3$ in
  the Fano plane, vs.\ observed $\delta_{CP} \approx 70^\circ$.
  Again a factor-of-2 ambiguity: $\pi/3 \approx 60^\circ$ vs. $\delta_{\mathrm{CL}}\text{-corrected} \approx 57^\circ$ vs. observed $70^\circ$. The CP phase residual is $\approx 20\%$,
  suggesting a second-order correction beyond $\delta_{\mathrm{CL}}$.
\end{enumerate}
\end{remark}

%% ============================================================
\subsection{Summary: The single correction needed}
%% ============================================================

\begin{theorem}[OP 1 and OP 4 reduce to a single calculation]\label{thm:op1op4-unified}
Both OP 1 and OP 4 are resolved once the continuum-limit correction
$\delta_{\mathrm{CL}}$ is derived from the lattice-to-continuum
transition of $\Gamma_P$. Specifically:
\begin{align}
  \bsym &= \frac{1}{\sqrt{3}} \cdot \delta_{\mathrm{CL}}, \\
  \theta_{12}^{\mathrm{CKM}} &= 2\arcsin(\gamma_{\mathrm{string}})
    \cdot \delta_{\mathrm{CL}}, \\
  \theta_{23}^{\mathrm{PMNS}} &= \arctan(\sqrt{2}) \cdot \delta_{\mathrm{CL}}.
\end{align}
The Conjecture~\ref{conj:cl-correction} provides a candidate:
$\delta_{\mathrm{CL}} = \exp(-\gamma_1/(4\pi\gamma_L^5))$, where
$\gamma_1$ is the first non-trivial Riemann zero. This is
\textbf{falsifiable}: it predicts $\bsym$ to $0.02\%$ accuracy
and $\theta_{12}^{\mathrm{CKM}}$ to approximately $6.5\%$ accuracy
(residual from higher-order lattice corrections).
\end{theorem}

%% ============================================================
%%  END OF SECTION — Status:
%%  OP 1: Factor-2 mechanism identified (Hopf projection doubling).
%%        Continuum-limit conjecture formulated (falsifiable).
%%        Formal proof of δ_CL still open.
%%  OP 4: Factor-2 explained (Hopf doubling S¹⁵$\to$S⁸).
%%        Full CKM/PMNS requires δ_CL proof (= same as OP 1).
%%        CP phase: second-order correction, ~20% residual.
%% ============================================================

%% ============================================================
%%  s_op5_lambda_cancellation.tex
%%  STATUS: PROOF MODE — Erster vollständiger Entwurf
%%  OP 5: $\Lambda$-Cancellation — $\Theta$-R-Balance im Photon-Frame
%%
%%  Marc Brendecke, Bochum, March 2026
%%  Einzubinden in quantum_sphaera_v2.x als Drop-In Section
%%
%%  FRAME-WARNUNG (vor jeder Rechnung zu beachten):
%%  Wir sitzen auf dem Matter-Frame (z⁻ = β_sym ≈ 0.5493).
%%  ALLE Messungen von Naturkonstanten — CODATA, NIST, BIPM —
%%  sind Matter-Frame-Werte. Der Photon-Frame (z⁻ = 0) ist
%%  der Hardware-Frame, in dem die $\Lambda$-Cancellation stattfindet.
%%  Matter-Frame-Rechnungen tragen den systematischen Fehler
%%  γ_L^{3/2} ≈ 2.546 im UV-Cutoff.
%% ============================================================


\subsection*{Phase IV-T: Photon-frame $\Lambda$ cancellation and closed $\Theta$--$R$ program}
\addcontentsline{toc}{subsection}{Phase IV-T: Photon-Frame $\Lambda$-Cancellation and Closed $\Theta$--$R$ Balance Program}
\label{sec:lambda_cancel}

\begin{remark}[Transition from Phase IV-S to the closed $\Lambda$-block]\label{rem:transition_s_to_t_closed}
Phase~IV-S treated residual geometric and mixing effects as admissible corrections of the late operatoric readout layer. The present phase sharpens the cosmological frame problem in the same spirit: the structural statement $\Lambda_{\mathrm{fundamental}}=0$ is separated cleanly from the still-open quantitative $\Theta$--$R$ balance problem, so that no extra vacuum ontology is introduced and no unresolved quantitative step is silently promoted to a theorem.
\end{remark}

\subsection{Starting point: the trace formula decomposition}

The central object is the spectral trace
\[
  \Sigma(\tau) := \mathrm{Tr}\!\bigl(e^{-i\tau L_0}\bigr)
  = \sum_\rho e^{-i\tau\gamma_\rho},
\]
where $\{\gamma_\rho\}$ are the imaginary parts of the non-trivial
Riemann zeros (Theorem~T6 of~\cite{brendecke2026rh}).
By the Guinand--Weil explicit formula (Theorem~3.3 of~\cite{brendecke2026rh}):
\begin{equation}\label{eq:gw}
  \Sigma(\tau) = 1 - \Theta(\tau) + R(\tau),
\end{equation}
where
\begin{align}
  \Theta(\tau) &:= \sum_\rho e^{-i\tau\gamma_\rho}
    \quad\text{(non-trivial zeros, oscillatory)},
    \label{eq:Theta-def}\\
  R(\tau) &:= -\frac{e^{-\tau/2}}{e^\tau - 1}
    \quad\text{(trivial zeros, smooth)}.
    \label{eq:R-def}
\end{align}
The variation of the trace-formula action
$S = \int_0^\infty \Sigma(\tau)\,d\tau$
with respect to $g^{\mu\nu}$ yields the Einstein equation
(Theorem~11.10 of the Companion):
\begin{itemize}
  \item the constant term $1$ contributes $\Lambda\,g_{\mu\nu}$
        (from the pole of $\zeta$ at $s=1$, giving
        $\Lambda \sim t_1^2 \approx 200\,m_P^4$),
  \item $\delta\Theta/\delta g^{\mu\nu}$ contributes $T_{\mu\nu}$
        (primitive orbits as matter),
  \item $\delta R/\delta g^{\mu\nu}$ contributes $G_{\mu\nu}$
        (trivial zeros via Gauss--Bonnet).
\end{itemize}
The bare $\Lambda \sim 200\,m_P^4$ is $10^{122}$ times the observed value.
The task of OP~5 is to show that $\Theta$ and $R$ cancel this to leave
a residual $\Lambda_{\mathrm{obs}} \sim \bsym^2$ in the \emph{photon frame}.

\subsection{Frame translation: matter frame $\to$ photon frame}

\begin{proposition}[UV cutoff in the photon frame]
\label{prop:uv-photon}
The natural UV cutoff for the vacuum energy density is the Planck
length. By Theorem~7.6 (Planck-photon):
\[
  \lP^{(\mathrm{photon})} = \gL^{-3/2}\cdot\lP^{(\mathrm{matter})}
  \approx 6.35\times10^{-36}\,\mathrm{m},
\]
with $k_{\mathrm{eff}}(\lP) = \tfrac{1}{2}(k(\hbar)+k(G_N)-3k(c))
= \tfrac{3}{2}$.
The vacuum energy densities in the two frames therefore satisfy
\begin{equation}\label{eq:rho-ratio}
  \frac{\rho_{\mathrm{vac}}^{(\mathrm{photon})}}
       {\rho_{\mathrm{vac}}^{(\mathrm{matter})}}
  = \left(\frac{\lP^{(\mathrm{matter})}}{\lP^{(\mathrm{photon})}}\right)^4
  = \gL^{6} \approx 42.0.
\end{equation}
\end{proposition}

\begin{proof}
Immediate from $\rho_{\mathrm{vac}} \sim \ell_P^{-4}$ and
Theorem~7.6 with $k_{\mathrm{eff}} = 3/2$, giving
$\ell_P^{(\mathrm{photon})} = \gL^{-3/2}\ell_P^{(\mathrm{matter})}$,
so $(\ell_P^{(\mathrm{matter})}/\ell_P^{(\mathrm{photon})})^4 = \gL^6$.
\end{proof}

\begin{corollary}[Matter-frame computations carry a systematic error]
\label{cor:matter-error}
Any $\Theta$-$R$ cancellation argument performed with the matter-frame
UV cutoff $\lP^{(\mathrm{matter})}$ carries a systematic inflation of
$\gL^6 \approx 42.0$ in the vacuum energy density. Such an argument
cannot establish cancellation to $10^{-122}$; it would leave a residual
of $\gL^6 \cdot 10^{-122} \approx 10^{-120.4}$ even after perfect
algebraic cancellation in the matter frame.
\end{corollary}

\begin{proof}
Immediate from Proposition~\ref{prop:uv-photon}: a matter-frame
calculation uses $\rho_{\mathrm{vac}}^{(\mathrm{matter})}
= \gL^{-6}\rho_{\mathrm{vac}}^{(\mathrm{photon})}$, so any
cancellation ratio established in the matter frame corresponds to a
ratio $\gL^6$-times larger in the photon frame.
\end{proof}

\subsection{The bare $\Lambda$ problem in this framework}

\begin{proposition}[Bare $\Lambda$ from the constant term]
\label{prop:bare-lambda}
The constant term $1$ in $\Sigma(\tau) = 1 - \Theta(\tau) + R(\tau)$
contributes a cosmological constant
\[
  \Lambda_{\mathrm{bare}}^{(\mathrm{matter})}
  \;\sim\; \gamma_1^2 \approx (14.135)^2 \approx 200\,m_P^4,
\]
where $\gamma_1 = 14.134725\ldots$ is the imaginary part of the first
non-trivial Riemann zero (the natural energy scale of the prime lattice).
In the photon frame:
\[
  \Lambda_{\mathrm{bare}}^{(\mathrm{photon})}
  = \gL^{-2}\cdot\Lambda_{\mathrm{bare}}^{(\mathrm{matter})}
  \approx \frac{200}{(1.864)^2} \approx 57.6\,m_{P,\mathrm{photon}}^4.
\]
\end{proposition}

\begin{proof}
The pole of $\zeta(s)$ at $s=1$ contributes, via the Mellin transform
$\int_0^\infty e^{-s\tau}\Theta(\tau)\,d\tau = \log\zeta(s)$, a residue
at $s=1$ equal to $1$. The inverse Laplace transform at $\tau=0^+$
gives a contribution proportional to the square of the first zero
$\gamma_1$ to the vacuum energy density. The corresponding photon-frame normalization is not the main point of the present proposition. What matters structurally is that the bare contribution remains Planckian in naive units and therefore has to be removed by the correct photon-frame interpretation rather than by a matter-frame cancellation trick. A fully self-consistent determination of the effective $k(\Lambda)$ belongs to the open quantitative balance programme of §\ref{ssec:open-balance}.
\end{proof}

\subsection{Main theorem: $\Lambda_{\mathrm{fundamental}} = 0$
in the photon frame}

\begin{theorem}[Vanishing fundamental $\Lambda$ in the photon frame]
\label{thm:lambda-zero}
In the photon frame ($z^- = 0$, $\bsym = 0$):
\begin{enumerate}
\item[(i)] The frame term $\mathcal{L}_{\mathrm{frame}} = \bsym^2
  \langle f,f\rangle$ vanishes identically, so no cosmological
  constant is contributed by the observer offset.
\item[(ii)] The Sphaera-cascade Markov process on $\{S^2, S^4, S^8\}$
  has a photon-frame absorption probability
  $P(S^8 \to S^2) = 1 - \bsym^2 = 1$ (in the photon frame
  where $\bsym\to 0$), yielding the structurally pure photon-frame collapse regime. In particular, the matter-displacement contribution vanishes at the level of the fundamental frame statement.
\item[(iii)] The observed cosmological constant
  $\Lambda_{\mathrm{obs}} \sim \bsym^2$ is therefore a pure
  \emph{matter-frame measurement artifact}:
  \[
    \Lambda_{\mathrm{obs}}
    = \bsym^2 \approx 0.302\;(\text{in units of }\lP^{-2}),
  \]
  arising from the permanent displacement of matter observers at
  $z^- = +\bsym$ from the photon ground shell $z^- = 0$.
\end{enumerate}
\end{theorem}

\begin{proof}
(i) In the photon frame, the Lagrangian term
$\mathcal{L}_{\mathrm{frame}} = \bsym^2\langle f,f\rangle$ vanishes
since $\bsym = 0$ by definition of the photon frame. Its variation
$\delta\mathcal{L}_{\mathrm{frame}}/\delta g^{\mu\nu} = \bsym^2
g_{\mu\nu}$ therefore contributes zero to the field equation.
This is the statement that the fundamental cosmological constant
vanishes in the photon frame.

(ii) The Markov transition matrix of the cascade
$\hat{K}: S^8 \to S^4 \to S^2$ is
\[
  P = \begin{pmatrix}1&0&0\\1&0&0\\1-\bsym^2&\bsym^2&0\end{pmatrix}.
\]
In the limit $\bsym\to 0$: $P(S^8\to S^4) = \bsym^2 \to 0$ and
$P(S^8\to S^2) = 1-\bsym^2 \to 1$.
All cascade flux collapses directly to the photon ground state.

(iii) A matter observer at $z^- = +\bsym$ reads off the frame term
as $\Lambda_{\mathrm{obs}} = \bsym^2/\lP^{(\mathrm{matter})2}$ via
Proposition~18.4 of the main paper. This is not a vacuum energy
density but the metric coefficient of the permanently shifted frame.
Since $\bsym^2 \approx (0.5493)^2 \approx 0.302$, and the observed
cosmological constant is $\Lambda_{\mathrm{obs}} \approx 1.1\times
10^{-52}\,\mathrm{m}^{-2}$, the Sphaera prediction is dimensionally
correct up to the $\gL^{3/2}$ Planck-length correction of Theorem~7.6.
\end{proof}

\subsection{The open quantitative balance}
\label{ssec:open-balance}

Theorem~\ref{thm:lambda-zero} establishes $\Lambda_{\mathrm{fundamental}}=0$
at the \emph{structural} level. The remaining open step is the
\emph{quantitative} $\Theta$--$R$ balance.

\begin{definition}[Photon-frame $\Theta$-$R$ balance condition]
\label{def:tr-balance}
The $\Theta$-$R$ balance in the photon frame is the condition
\begin{equation}\label{eq:tr-balance}
  \int_0^{\lP^{(\mathrm{photon})}} \!\!
  \bigl(1 - \Theta^{(\mathrm{photon})}(\tau) + R^{(\mathrm{photon})}(\tau)\bigr)
  \,d\tau \;=\; \bsym^2 \cdot \lP^{(\mathrm{photon})2},
\end{equation}
where $\Theta^{(\mathrm{photon})}$ and $R^{(\mathrm{photon})}$ are the
photon-frame unfolded versions of~\eqref{eq:Theta-def}--\eqref{eq:R-def},
with the photon-frame density $\rho_{\mathrm{photon}}(\gamma) =
\log(\gamma/2\pi\gL)/(2\pi\gL)$ replacing the matter-frame density.
\end{definition}

\begin{remark}[Why this is not yet proved]
\label{rem:open}
Condition~\eqref{eq:tr-balance} says: after photon-frame unfolding,
the net integral of $\Sigma^{(\mathrm{photon})}$ equals the frame-offset
residual $\bsym^2$ — nothing more, nothing less. The $10^{-122}$ refers
to the ratio $\bsym^2\lP^{2}/\Lambda_{\mathrm{bare}} \approx 0.302 /
(200\cdot\lP^{(\mathrm{matter})2}/\lP^{(\mathrm{photon})2})$
$= 0.302 \cdot \gL^{-3} / 200 \approx 0.302/1298 \approx 2.3\times
10^{-4}$ in dimensionless Planck units, not $10^{-122}$ in SI units.

\medskip
\noindent The $10^{-122}$ in SI units arises because $\lP^{-2}$ in SI is
$3.8\times10^{69}\,\mathrm{m}^{-2}$, while $\Lambda_{\mathrm{obs}}
\approx 10^{-52}\,\mathrm{m}^{-2}$, giving ratio $\approx 10^{-122}$.
In natural units ($\lP=1$), $\Lambda_{\mathrm{bare}} \sim \gamma_1^2
\approx 200$ and $\Lambda_{\mathrm{obs}} \sim \bsym^2 \approx 0.30$,
so the ratio is $\approx 1.5\times 10^{-3}$ — much less dramatic. The
$10^{-122}$ is a unit-system artifact, not the fundamental challenge.

\medskip
\noindent The fundamental challenge is to show that $\Theta$ and $R$ in
$\Sigma = 1 - \Theta + R$ cancel the bare $\Lambda \approx 200$ down to
$\bsym^2 \approx 0.30$ with relative accuracy $\sim 10^{-3}$, uniformly
in the photon-frame integration.
\end{remark}

\begin{proposition}[Two open sub-steps]
\label{prop:open-steps}
Condition~\eqref{eq:tr-balance} is equivalent to the conjunction of:
\begin{enumerate}
\item[\textbf{(A)}] \textbf{$\Theta$-sector remainder control.}
  The partial sum $\Theta^{(\mathrm{photon})}_N(\tau) :=
  \sum_{n=1}^N e^{-i\tau\gamma_n^{(\mathrm{photon})}}$
  satisfies, uniformly for $0 \leq \tau \leq \lP^{(\mathrm{photon})}$:
  \[
    \left|\Theta^{(\mathrm{photon})}(\tau)
    - \Theta^{(\mathrm{photon})}_N(\tau)\right|
    \;\leq\; \varepsilon_N \to 0 \text{ as } N\to\infty,
  \]
  and $\int_0^{\lP^{(\mathrm{photon})}} \Theta^{(\mathrm{photon})}(\tau)\,d\tau$
  converges absolutely.
\item[\textbf{(B)}] \textbf{R-sector exact cancellation.}
  The smooth term $R^{(\mathrm{photon})}(\tau) = -e^{-\tau/2}/(e^\tau-1)$
  (photon-frame rescaled) and the sum
  $\Theta^{(\mathrm{photon})}(\tau)$ together cancel the constant $1$
  in $\Sigma$ to within $\bsym^2$ when integrated over
  $[0, \lP^{(\mathrm{photon})}]$.
\end{enumerate}
Sub-step~(B) is equivalent to the Guinand--Weil explicit formula
evaluated at the photon-frame UV cutoff $\lP^{(\mathrm{photon})}$.
\end{proposition}

\begin{proof}[Proof of equivalence]
The integral of $\Sigma^{(\mathrm{photon})}$ splits as
\[
  \int_0^{\ell} \Sigma\,d\tau
  = \ell - \int_0^\ell \Theta\,d\tau + \int_0^\ell R\,d\tau,
\]
where $\ell = \lP^{(\mathrm{photon})}$.
Condition~(A) ensures the $\Theta$-integral is well-defined and
controllable. Condition~(B) then says the three terms sum to
$\bsym^2\cdot\ell^2$.
Both conditions are necessary and sufficient for~\eqref{eq:tr-balance}.
\end{proof}

\begin{theorem}[Structural $\Theta$-R cancellation — partial result]
\label{thm:partial-cancel}
In the matter frame, the Guinand--Weil formula gives exactly
\[
  \Sigma(\tau) = 1 - \Theta(\tau) + R(\tau),
\]
which encodes complete cancellation of the pole contribution by the
combined $\Theta$ and R terms \emph{at the distributional level} (for all
$\tau > 0$, without UV cutoff). The cosmological constant problem
arises solely from introducing a finite UV cutoff: at
$\tau = \lP^{(\mathrm{matter})}$, the integral of $\Sigma$ is
\[
  \int_0^{\lP^{(\mathrm{matter})}}\!\!\Sigma\,d\tau
  = \lP^{(\mathrm{matter})} - \Theta_{\mathrm{UV}}
  + R_{\mathrm{UV}} + \bsym^2\,(\lP^{(\mathrm{matter})})^2,
\]
where $\Theta_{\mathrm{UV}}$ and $R_{\mathrm{UV}}$ are the truncated
integrals. The $\bsym^2$-term is the frame residual. The mismatch
$\lP^{(\mathrm{matter})} - \Theta_{\mathrm{UV}} + R_{\mathrm{UV}}$
is the open step~(B) in the matter frame.
\end{theorem}

\begin{proof}
The Guinand--Weil explicit formula holds distributionally for all
$\tau > 0$ (Theorem~3.3 of~\cite{brendecke2026rh}). The $\Theta$-integral
converges because $\Theta(\tau) \to 1$ as $\tau\to 0^+$ along the
spectral sum (distributional convergence). The exact cancellation
$1 - \Theta(\tau) + R(\tau) = \Sigma(\tau)$ holds for all $\tau > 0$
by the explicit formula. The finite-UV-cutoff breakdown is the
statement that $\int_0^\Lambda\Sigma\,d\tau \neq 0$ for finite
$\Lambda$, which is the content of the open sub-steps.
\end{proof}

\subsection{Falsifiable predictions}
\label{sec:falsifiable}

\begin{corollary}[Ranked falsifiable predictions]
\label{cor:falsifiable}
The following predictions follow from the framework and are ordered by
near-term experimental accessibility.

\bigskip
\noindent\textbf{P1 (testable now): GUE improvement on Odlyzko data.}\\
Photon-frame unfolding of Odlyzko's ${\sim}10^6$ Riemann zeros at height
${\sim}10^{22}$ predicts a $\chi^2$ improvement of $\approx 21\%$ relative
to standard unfolding
(Corollary~\ref{cor:odlyzko}, Theorem~\ref{thm:gue-selfdet}).
Non-observation falsifies $\gL = \gBI/\gstring$ as the physical frame factor.
\emph{Status: testable with Odlyzko's public dataset using existing code.}

\bigskip
\noindent\textbf{P2 (already matches data): Koide phase formula.}\\
$\theta_K = \pi - \sqrt{2}\arcsin(\bsym) \approx 132.88^\circ$, matching
the PDG value $132.73^\circ$ to $0.111\%$
(Theorem~\ref{thm:koide-triolic}).
\emph{Status: confirmed against existing PDG lepton masses.}

\bigskip
\noindent\textbf{P3 (matches to 0.018\%): $\bsym$-formula.}\\
\[
  \bsym = \frac{1}{\sqrt{3}}\,\exp\!\left(-\frac{\gamma_1}{4\pi\gL^5}\right)
  \approx 0.5492,
\]
matching the measured value $0.5493$ to $0.018\%$
(Remark~\ref{conj:cl-correction}, corrected in v2.29.1).
\emph{Status: consistent with CODATA. Structural derivation of the $\gL^5$
exponent requires OP~2.}

\bigskip
\noindent\textbf{P4 (reinterpretation): Cosmological constant.}\\
$\Lambda_{\mathrm{obs}} = \bsym^2 \cdot (\lP^{(\mathrm{photon})})^{-2}$
is a frame artifact, not a vacuum energy (Theorem~\ref{thm:lambda-zero}).
In natural Planck units the apparent $10^{-122}$ fine-tuning reduces to
a factor ${\sim}3.1$ cancellation (Remark~\ref{rem:open}).
Matter-frame computations carry a systematic error of $\gL^6 \approx 42$
(Corollary~\ref{cor:matter-error}).
\emph{Status: reinterpretation of existing data; no new measurement needed.}

\bigskip
\noindent\textbf{P5 (annual modulation, future): $\alpha$-drift ratio.}\\
Constants with different frame-insertion counts $k$ should modulate
differently with annual orbital period. Specifically, the ratio of
annual variations:
\[
  \frac{\delta\alpha/\alpha}{\delta G_N/G_N}\bigg|_{\mathrm{annual}}
  = \gL^{k(\alpha)-k(G_N)} \cdot \frac{k(\alpha)}{k(G_N)}
  = \gL^{-1} = \frac{\gstring}{\gBI} \approx 0.537,
\]
whereas standard physics predicts this ratio $= 1$.
The deviation is $\approx 46\%$ but requires $G_N$ measurements at
${\sim}10^{-18}$ annual precision --- beyond current capability.
\emph{Status: prediction well-defined but not yet testable. The modulation
of $\alpha$ alone (period exactly 1 year, in phase at all terrestrial labs,
absent for non-orbiting reference clocks) is a qualitative signature
at amplitude ${\sim}10^{-10}\cdot\delta\Phi/c^2$.}

\bigskip
\noindent\textbf{P6 (theoretical): Hawking temperature in tachyon frame.}\\
$T_{\mathrm{Sphaera}} = \gBI\cdot T_H^{(\mathrm{tachyon})}$, giving
$\gBI$ the interpretation as the ratio of the Sphaera Hawking temperature
to the tachyon-frame Hawking temperature
(Proposition~\ref{prop:hawking-gBI-tachyon}).
\emph{Status: theoretical prediction; requires quantum-gravity experiment
to test directly.}
\end{corollary}

\begin{remark}[What the predictions share]
All six predictions follow from the same structural core: the frame offset
$\zminus = +\bsym$ and the frame factor $\gL = \gBI/\gstring$. None require
free parameters beyond these two quantities, both of which are determined
by the framework (P3 and the LQG/string coupling values).
The only currently testable prediction that requires new computation is P1
(GUE improvement); P2 is already a post-diction on existing data.
\end{remark}

\subsection{Status summary}

\begin{quote}\small
\noindent\textbf{Reader cue.} Read the next table as a burden/status dashboard for the OP~5 line. Its purpose is to separate what is already structurally fixed, what remains open, and which analytic step each open item still depends on.
\end{quote}

\begin{center}
\renewcommand{\arraystretch}{1.25}
\setlength{\tabcolsep}{4pt}
\begin{tabular}{p{0.33\textwidth} p{0.31\textwidth} p{0.24\textwidth}}
\toprule
Statement & Proof basis & Status \\\midrule
$\Lambda_{\mathrm{fundamental}} = 0$ in photon frame
  & Cascade Markov + $\bsym=0$ in photon frame
  & \checkmark (Thm.~\ref{thm:lambda-zero}) \\
$\Lambda_{\mathrm{obs}} = \bsym^2$ (frame artifact)
  & $\mathcal{L}_{\mathrm{frame}}$ variation
  & \checkmark (Prop.~18.4 of main paper) \\
Matter-frame error = $\gL^6 \approx 42.0$
  & $k(\lP) = 3/2$, Thm.~7.6
  & \checkmark (Cor.~\ref{cor:matter-error}) \\
Guinand--Weil: $\Sigma = 1-\Theta+R$ distributionally
  & Explicit formula, Thm.~3.3
  & \checkmark (Thm.~\ref{thm:partial-cancel}) \\
\midrule
$\Theta$-sector remainder control at $\lP^{(\mathrm{photon})}$ (step A)
  & —
  & \textbf{open} \\
R-sector exact cancellation at photon-frame cutoff (step B)
  & —
  & \textbf{open} \\
$k(\Lambda)$ self-consistent determination
  & —
  & \textbf{open} \\
\bottomrule
\end{tabular}
\end{center}

\begin{remark}[What OP~5 actually requires]
\label{rem:op5-requires}
The two open steps are not independent: step~(B) is the
photon-frame version of the Guinand--Weil formula evaluated at a
\emph{finite} UV cutoff. This is a standard problem in analytic
number theory (partial explicit formulas with error terms). The
Sphaera-framework contribution is: (a) identifying the correct
frame ($\lP^{(\mathrm{photon})}$, not $\lP^{(\mathrm{matter})}$),
and (b) replacing the observable $\Lambda$ with the frame residual
$\bsym^2$. The hard analysis — bounding the truncation error of the
zero sum at height $T = 1/\lP^{(\mathrm{photon})}$ — is the
remaining mathematical content of OP~5.
\end{remark}

\subsection{Quantitative analysis of the open steps (v2.29.0)}
\label{ssec:op5-quantitative}

\subsubsection*{Step (A): $\Theta$-remainder control}

\begin{proposition}[Step~(A) is trivially satisfied in the photon frame]
\label{prop:theta-remainder-trivial}
The photon-frame $\Theta$-remainder control of sub-step~(A) is
satisfied trivially:
\[
  \Theta^{(\mathrm{photon})}_N(\tau)
  = \Theta^{(\mathrm{photon})}(\tau)
  = 0
  \qquad
  \text{for all } \tau \in [0,\,\lP^{(\mathrm{photon})}].
\]
Consequently $\varepsilon_N = 0$ for all $N$ in this range.
\end{proposition}

\begin{proof}
The photon-frame UV cutoff is $T_{\mathrm{photon}} =
1/\lP^{(\mathrm{photon})} = \gL^{3/2} \approx 2.545$ (in
matter-Planck units). The photon-frame Riemann zeros are
$\gamma_\rho^{(\mathrm{photon})} = \gamma_\rho/\gL$, with first
value $\gamma_1^{(\mathrm{photon})} = 14.1347/1.864 \approx 7.583$.
Since $T_{\mathrm{photon}} \approx 2.545 < \gamma_1^{(\mathrm{photon})}
\approx 7.583$, no photon-frame zero lies below the UV cutoff.
Therefore, for $\tau \in [0, \lP^{(\mathrm{photon})}]$, the
oscillatory sum $\Theta^{(\mathrm{photon})}(\tau) =
\sum_\rho e^{-i\tau\gamma_\rho^{(\mathrm{photon})}}$ has no
contributing terms in the spectral range
$[0, T_{\mathrm{photon}}]$.
Hence $\Theta^{(\mathrm{photon})}(\tau) = 0$ pointwise in the
integration domain, and step~(A) is satisfied with $\varepsilon_N = 0$.
\end{proof}

\begin{remark}[Physical significance: the photon frame is below the first zero]
\label{rem:photon-frame-below-first-zero}
The fact that $T_{\mathrm{photon}} < \gamma_1^{(\mathrm{photon})}$
is a structural consequence of the frame factor $\gL \approx 1.864$:
in the photon frame, the UV cutoff $\lP^{(\mathrm{photon})}$ is
\emph{larger} than in the matter frame by $\gL^{3/2}$, pushing the
cutoff frequency $T_{\mathrm{photon}}$ \emph{below} the energy of
the first Riemann zero. This is precisely the regime where the
spectrum is quiet --- the photon frame sees no Riemann zero
oscillations below its Planck scale. The cosmological constant
problem vanishes in this regime because $\Theta$ contributes nothing
to the vacuum energy integral.
\end{remark}

\subsubsection*{Step (B): R-sector cancellation}

With step~(A) resolved, the balance condition~\eqref{eq:tr-balance}
reduces to:
\begin{equation}
\label{eq:balance-reduced}
  \int_0^{\lP^{(\mathrm{photon})}}
  \bigl(1 + R^{(\mathrm{photon})}(\tau)\bigr)\,d\tau
  \;=\; \bsym^2\cdot(\lP^{(\mathrm{photon})})^2.
\end{equation}

\begin{remark}[Regularization required]
\label{rem:regularization-B}
The integrand $1 + R(\tau)$ has a pole at $\tau=0$: since
$R(\tau) = -e^{-\tau/2}/(e^\tau-1) \sim -1/\tau + 1/2 - \tau/12 +
\cdots$ as $\tau\to 0^+$, the combination $1 + R(\tau) \sim
3/2 - 1/\tau + \cdots$ is still singular.
Equation~\eqref{eq:balance-reduced} must therefore be read as a
\emph{renormalized} integral:
\[
  \mathrm{p.v.}\int_0^T (1 + R(\tau))\,d\tau
  \;:=\;
  \lim_{\varepsilon\to 0}
  \Bigl[
    \int_\varepsilon^T (1+R(\tau))\,d\tau + \log\varepsilon
  \Bigr],
\]
where the $\log\varepsilon$ counterterm removes the leading pole.
Numerically with $T = \lP^{(\mathrm{photon})} \approx 0.393$
(matter-Planck units):
\[
  \mathrm{p.v.}\int_0^T (1+R)\,d\tau
  \;=\; \frac{3T}{2} - \log T - \frac{T^2}{24} - \gamma_E
  \;\approx\; 0.940,
\]
where $\gamma_E \approx 0.5772$ is the Euler--Mascheroni constant.
\end{remark}

\begin{proposition}[Step~(B) open: required Theta contribution]
\label{prop:theta-required-B}
\emph{(Open sub-step; not yet proven.)}
For the balance condition~\eqref{eq:tr-balance} to hold, the full
Theta sum (integrated over all zeros, not just those below
$T_{\mathrm{photon}}$) must contribute:
\[
  \int_0^{\lP^{(\mathrm{photon})}}
  \Theta^{(\mathrm{photon})}(\tau)\,d\tau
  \;=\;
  \mathrm{p.v.}\int_0^T (1+R)\,d\tau - \bsym^2\cdot T^2
  \;\approx\; 0.940 - 0.302\cdot 0.154
  \;\approx\; 0.893.
\]
This integral runs over all photon-frame zeros $\gamma_\rho^{(\mathrm{photon})}$
\emph{including those above} $T_{\mathrm{photon}}$, which oscillate
rapidly in the integration range $[0, T]$ and must be summed
carefully.
Establishing this cancellation constitutes the remaining hard
mathematical content of OP~5.
\end{proposition}

\begin{remark}[Structural summary of OP~5]
\label{rem:op5-structural-summary}
The analysis of v2.29.0 reveals the following structure of OP~5:
\begin{itemize}
  \item Step~(A) is closed: $\Theta^{(\mathrm{photon})}=0$ below
    the photon-frame cutoff (Proposition~\ref{prop:theta-remainder-trivial}).
  \item Step~(B) reduces to showing that the oscillatory Theta sum,
    integrated over the photon-frame UV range $[0,\lP^{(\mathrm{photon})}]$
    but summing over all zeros (including those above the cutoff),
    evaluates to $\approx 0.893$ in renormalized Planck units.
  \item The $10^{-122}$ apparent fine-tuning is a unit-system
    artifact (Remark~\ref{rem:open}): in natural Planck units the
    required cancellation is from $\sim 0.940$ (R-sector baseline)
    to $0.302$ (frame residual), a factor of $\sim 3.1$, not $10^{122}$.
\end{itemize}
\end{remark}

\subsubsection*{Step (B) reduced: the exact Theta-sum conjecture}

The numerical analysis of v2.29.0 yields a precise formulation of
what step~(B) requires. Define:

\begin{definition}[Photon-frame Theta sum]
\label{def:theta-sum-B}
Let $T = \gL^{-3/2}$ (the photon-frame UV cutoff in matter-Planck units)
and let $\gamma_n^{(\mathrm{photon})} = \gamma_n/\gL$ denote the
photon-frame Riemann zeros. The \emph{photon-frame Theta sum} is the
conditionally convergent series
\[
  \mathcal{T}(T)
  \;:=\;
  2\sum_{n=1}^\infty
  \frac{\sin\!\bigl(T\cdot\gamma_n^{(\mathrm{photon})}\bigr)}
       {\gamma_n^{(\mathrm{photon})}}.
\]
\end{definition}

\begin{proposition}[Theta-sum balance conjecture, OP~5 Step~B]
\label{conj:theta-sum-B}
\emph{(Conjecture-grade; derivable from the Guinand--Weil explicit
formula if the remainder estimate can be controlled.)}
\[
  \mathcal{T}(T)
  \;=\;
  \frac{3T}{2} - \log T - \frac{T^2}{24} - \gamma_E - \bsym^2 T^2
  \;\approx\; 0.893,
\]
where $T = \gL^{-3/2}\approx 0.393$, $\gamma_E\approx 0.5772$ is
the Euler--Mascheroni constant, and $\bsym^2\approx 0.302$.
\end{proposition}

\begin{remark}[Structure of the conjecture]
\label{rem:theta-sum-structure}
The series $\mathcal{T}(T)$ is the photon-frame analogue of the
oscillatory sum in the Riemann--von Mangoldt explicit formula for
the prime-counting function $\pi(x)$. It converges \emph{conditionally}
(not absolutely), and its value at a specific argument $T$ is a
non-trivial statement about the Riemann zeros.

Numerically, $\mathcal{T}(T)$ converges slowly: with $N=50$ zeros
the partial sum is $\approx -0.146$, far from the conjectured $0.893$.
Convergence requires summing the full tail $\sum_{n>50}$ whose
individual contributions are small but whose cumulative effect is
large (this is the hallmark of conditional convergence).

The conjecture is \emph{equivalent to} the balance condition
$\eqref{eq:tr-balance}$ once step~(A) is established
(Proposition~\ref{prop:theta-remainder-trivial}). Establishing
Conjecture~\ref{conj:theta-sum-B} would therefore close OP~5
completely.
\end{remark}

\begin{remark}[Connection to the explicit formula remainder]
\label{rem:explicit-formula-remainder}
By the Guinand--Weil explicit formula
(Theorem~3.3 of~\cite{brendecke2026rh}), the identity
$\Sigma(\tau) = 1 - \Theta(\tau) + R(\tau)$ holds for all $\tau > 0$.
Integrating from $0$ to $T$ (in the renormalized sense) and
substituting the known values of $\int_0^T 1\,d\tau = T$ and
$\int_0^T R\,d\tau_{\mathrm{reg}} = 3T/2 - \log T - T^2/24 - \gamma_E$,
Conjecture~\ref{conj:theta-sum-B} is equivalent to:
\[
  \int_0^T \Sigma^{(\mathrm{photon})}(\tau)\,d\tau
  \;=\; \bsym^2 \cdot T^2.
\]
This is precisely the balance condition $\eqref{eq:tr-balance}$.
Hence: \emph{Conjecture~\ref{conj:theta-sum-B} $\Leftrightarrow$
OP~5 Step~(B) $\Leftrightarrow$ Balance condition~\eqref{eq:tr-balance}.}
\end{remark}
